%% air_freight_routing.tex
%% Air Freight Multi-Commodity Routing — MVP Formal Model
%% Phase 1 — must be approved before any code starts
%% See model/air_graph_construction.md for the graph layer this builds on.

\documentclass[11pt,a4paper]{article}

\usepackage[T1]{fontenc}
\usepackage{amsmath,amssymb,amsthm}
\usepackage{bbm}
\usepackage{booktabs}
\usepackage{geometry}
\usepackage{hyperref}
\usepackage{array}
\usepackage{xcolor}
\usepackage{enumitem}
\usepackage{longtable}
\usepackage{graphicx}
\usepackage{float}

\geometry{margin=2.5cm}
\setlength{\parskip}{0.4em}
\hypersetup{colorlinks=true, linkcolor=blue}

\newcommand{\ceil}[1]{\left\lceil #1 \right\rceil}
\newcommand{\floor}[1]{\left\lfloor #1 \right\rfloor}
\newcommand{\defeq}{\triangleq}
\newcommand{\ctype}{\text{ctype}}

\theoremstyle{definition}
\newtheorem{defn}{Definition}
\newtheorem{remark}{Remark}

\title{\textbf{Air Freight Multi-Commodity Routing}\\[0.4em]
\large MVP Formal Model --- Phase 1, v3 (O-D-arc graph)}
\author{AI Freight Agent}
\date{Draft --- 2026-05-23}

\begin{document}
\maketitle

\begin{abstract}
Formal specification of the air freight multi-commodity routing MILP. The model
routes HAWBs (House Air Waybills) on an \emph{O-D-arc graph} (see
\texttt{model/air\_graph\_construction.md} and
\texttt{docs/air\_graph\_construction.drawio}). Arcs are O-D segments of three
air types --- \emph{MAWB-arc} (forwarder-consolidated carriage of one O-D
segment under one offer), \emph{co-load arc} (forwarder buys per-kg space from
another consolidator), and \emph{deconsolidation-dwell arc} (at a
forwarder-operated hub CFS) --- plus ground arcs for pickup, cartage, CFS
dwell, customs clearance, and final delivery. The central commercial object is
the \emph{MAWB} $= (\text{MAWB-arc},\, g)$: one carriage contract over an O-D
segment, identified by a MAWB-arc and a consolidation group \(g\). HAWBs
consolidate by sharing a MAWB-arc within the same \(g\). Physical flight legs
sit \emph{beneath} a MAWB-arc as metadata; the MILP does not index over flights
directly. Air capacity is procured through a catalog of offers (TACT / SCR /
CCR / BSA / NAC / co-load / spot), surfaced to the MILP as per-arc rate
families \texttt{flat\_rate}, \texttt{min\_flat\_breaks}, \texttt{per\_uld\_pivot},
and \texttt{coload\_per\_kg} (\S\ref{sec:supply-option-catalog}). The IATA
next-break-down rule is the \texttt{min\_flat\_breaks} family, encoded with a
per-MAWB break-selection binary and a 3-inequality disaggregation
(\S\ref{sec:lin-bucket}). Block-Space-Agreement take-or-pay is priced
per-flight via a pivot floor (\texttt{settlement\_basis} \(=\)
\texttt{per\_flight}) or via an exogenous sunk allowance \(A_c\)
(\texttt{settlement\_basis} \(=\) \texttt{equalized}). Operational scope:
cargo-type compatibility (GEN / DGR / PER / VAL / AVI / HUM); cargo-ready
window; effective cutoff stack (DCO + AMS + ICS2 + ACI) folded per-arc;
embargo and lithium-battery PI gating; tenant-defined service products with a
\emph{quadratic} tardiness penalty scaled by service-product weight
\(w_{\text{sp}(k)}\) and per-HAWB shipment-value coefficient \(\mu_k\)
(PWL-linearized via convex tangent cuts, \S\ref{sec:lin-tardiness}); a
\emph{guaranteed-feasibility fallback arc} per HAWB (direct origin \(\to\)
destination, arriving at \(T_k^{\text{abs}}\) --- now mandatory-finite for every
HAWB, semantically the ``latest valid arrival time'' --- at dollar cost
\(C^{\text{fallback}}\), selected only when no real route exists; the MILP
carries no hard time bound, and fallback-arc usage replaces solver-infeasibility
as the rescue signal, \S\ref{sec:fallback-arc}); locked commitments resolved at
preprocessing (not as MILP constraints); and per-HAWB carrier eligibility
resolved from four rule layers
(tenant blacklist, shipper-lane allow/deny list, service-product carrier list,
commodity overlay) by intersecting allows and unioning denies,
so an explicit deny at any layer blocks the carrier even when a higher-layer
rule would allow it (\S\ref{sec:carrier-policy}).
Per-shipment hard budget caps are excluded by design --- a hard cap on a
committed must-serve shipment can produce spurious infeasibility; budget is a
quoting-layer concern. Probabilistic transit times, multi-tier per-ULD
over-pivot rate structures, lock-buyout as a MILP decision, split-shipment
across flights, and full 3D bin-packing inside ULDs are deferred
(\S\ref{sec:deferred}).
\end{abstract}

\tableofcontents
\newpage

%%--------------------------------------------------------------------
\section{Problem Statement}
\label{sec:problem}
%%--------------------------------------------------------------------

Given a set of HAWBs $K$ to route this solve --- each with origin door, destination
door, actual weight $w_k$, volume $v_k$, cargo class,
temperature regime, service product binding $\text{sp}(k)$, cargo-ready window
$[t_k^{\text{rdy,early}}, t_k^{\text{rdy,late}}]$, soft deadline $T_k^{\text{SLA}}
= t_k^{\text{rdy,early}} + T^{\text{SLA}}_{\text{sp}(k)}$, latest valid arrival
time $T_k^{\text{abs}}$ (mandatory-finite for every HAWB; semantically the
cargo-death point for PER, customer-cancellation point for GEN), and declared
shipment value $\text{value}_k$ --- find a minimum-cost routing of each HAWB
through the O-D-arc graph such that:
\begin{itemize}[noitemsep]
  \item every HAWB enters at its origin door and exits at its destination door
        (flow conservation, C.1);
  \item HAWBs consolidating on a MAWB-arc share a MAWB $(a, g)$ iff they share
        the consolidation group $g$, where $g$ is a single-valued function of
        cargo class and temperature regime (consolidable cargo) or a singleton
        key for non-consolidable cargo (VAL / HUM / AVI) --- see
        \S\ref{sec:mawb-hawb} and \texttt{air\_graph\_construction.md \S4};
  \item every MAWB respects the per-contract allotment of ULD positions
        $N_{a,u}$ (C.5), the per-ULD physical capacity $W_u, V_u$ (C.5b), and
        any specified per-offer cap $\text{cap}_a$ (C.5c);
  \item cargo cutoffs $\text{CO}_a^*$ at every intermediate POL and hub-MCT
        timing are enforced \emph{at graph build} by forward-time-window
        propagation (\S\ref{sec:fwd-time-propagation}): an arc is admitted to
        the per-HAWB subgraph only if the propagated arrival window at its
        head is consistent with any hard time window (flight cutoff) there.
        The MILP carries \emph{no} time-propagation, cutoff, or hub-MCT
        constraint family;
  \item every HAWB is guaranteed a feasible routing via a per-HAWB fallback
        arc (direct $O_k \to D_k^{\text{node}}$, arriving at $T_k^{\text{abs}}$,
        dollar cost $C^{\text{fallback}}$; \S\ref{sec:fallback-arc}) --- the
        MILP carries no hard time bound, and fallback-arc usage replaces
        solver-infeasibility as the rescue signal;
  \item tardiness $\tau_k$ relative to the soft deadline $\Delta_k =
        \min(T_k^{\text{dead}}, T_k^{\text{SLA}})$ enters the objective
        \emph{quadratically} with combined weight $W_k = w_{\text{sp}(k)} \cdot
        \mu_k$, where $\mu_k = \text{value}_k / V^{\text{ref}}$ is the
        shipment-value coefficient (C.10, PWL-linearized in
        \S\ref{sec:lin-tardiness});
  \item already-committed (locked) HAWBs are resolved at the \emph{preprocessing}
        layer --- fully locked HAWBs are removed from $K$; partially locked
        HAWBs enter $K$ with origin re-pointed to current observed position +
        subgraph truncated to forward arcs + initial time fixed
        (\S\ref{sec:locked-commitments}, C.12); the MILP itself is
        lock-agnostic;
  \item cargo-type compatibility (cargo\_type\_ok), embargo (embargo\_ok),
        lithium battery PI gating (lithium\_ok), and the resolved
        service-product / carrier-policy cascade are enforced in the
        per-shipment subgraph pre-filter
        (\S\ref{sec:prefilter}); the MILP optimizes over arcs that have already
        survived these predicates.
\end{itemize}

\paragraph{What is \emph{not} a hard constraint in this MVP.}
\begin{itemize}[noitemsep]
  \item \textbf{Per-shipment budget cap.} A hard per-shipment cost cap can render
        an accepted must-serve shipment infeasible --- operationally wrong.
        Removed entirely (\S\ref{sec:deferred}). Cost minimization happens
        through the objective alone; budget is a quoting-layer concern.
  \item \textbf{Flight-level physical capacity $W_f, V_f$.} The forwarder cannot
        observe other parties' bookings on the same flight; carrier overbooking
        and offload is the airline's responsibility, not a forwarder planning
        concern. Capacity binds only at the per-contract / per-offer / per-ULD
        layer (\S\ref{sec:supply-option-catalog}).
  \item \textbf{Hub minimum connect time.} Absorbed into per-arc scalars at graph
        build (\S\ref{sec:graph-construction}); forward-time-window propagation
        (\S\ref{sec:fwd-time-propagation}) enforces the timing. No MILP
        constraint family.
  \item \textbf{Hard transit-time SLA $T_p^{\text{SLA}}$.} Promoted to a soft
        quadratic penalty $W_k \cdot \tau_k^2$ in the objective.
  \item \textbf{Hard backstop $T_k^{\text{abs}}$.} Removed as a constraint.
        $T_k^{\text{abs}}$ is retained as the arrival time of the per-HAWB
        fallback arc (\S\ref{sec:fallback-arc}), making it the parameter that
        bounds the worst-case quadratic penalty in the objective rather than a
        hard feasibility cutoff. The MILP has \emph{no} hard time bound and is
        always feasible by construction.
\end{itemize}

\paragraph{MVP scope --- key choices.}
\begin{itemize}[noitemsep]
  \item Each arc carries one offer; $\text{rate\_family}_a \in \{\texttt{flat\_rate},
        \texttt{min\_flat\_breaks}, \texttt{per\_uld\_pivot},
        \texttt{coload\_per\_kg}\}$ (\S\ref{sec:supply-option-catalog}).
  \item MAWB $=$ (MAWB-arc, $g$). Membership of HAWB $k$ in MAWB $(a, g)$ is
        \emph{implicit}: $k$ rides MAWB $(a, g)$ iff $x_{k,a} = 1$ and $g(k) =
        g$. No explicit $h_{k,m}$ assignment variable; no MAWB-count decision;
        no permutation symmetry.
  \item Consolidation group is a partition by construction:
        $g(k) = (\text{cargo\_class}(k), \text{temperature}(k))$ if
        $\text{cargo\_class}(k) \in \{\text{GEN, DGR, PER}\}$ (consolidable);
        else $g(k) = (\text{cargo\_class}(k), \text{HAWB-id}(k))$
        (non-consolidable: VAL, HUM, AVI each get a singleton group). See
        \texttt{air\_graph\_construction.md \S4} for the partition analysis.
  \item DGR is one coarse group at the MAWB layer; fine-grained pairwise DG
        segregation is a ULD-layer concern, deferred (\S\ref{sec:deferred}).
  \item Surcharges (Path A per-arc + Path B per-flight) carry over from the
        prior model (\S\ref{sec:surcharge}); the math is unchanged by the graph
        rewrite.
  \item ULD interchange policy (Star / SkyTeam / oneworld pool, bilateral
        Cargo SPA) influences which \emph{MAWB-arcs} the graph generator emits
        (a same-AWB through-segment under interline is one MAWB-arc, not two);
        it does not appear in the MILP body. See
        \S\ref{sec:through-uld-graph}.
  \item Sell rate and margin are upstream of this MILP (quote engine at booking
        time); the MILP outputs cost only.
\end{itemize}

%%--------------------------------------------------------------------
\section{Time-zone Convention}
\label{sec:tz-convention}
%%--------------------------------------------------------------------

All timestamps, time-window parameters, the destination-arrival decision
variable $t_k^{D_k^{\text{node}}}$, the graph-build forward-time-window
intervals, and constraint-side time bounds in this model are expressed in
\textbf{Coordinated Universal Time (UTC)}. This convention is normative for
every internal computation; conversions to local time happen only at I/O
boundaries (operator display, customer notifications, external API payloads
keyed on origin-local or destination-local).

\paragraph{Why UTC internally.}
\begin{itemize}[noitemsep]
  \item IATA SSIM Chapter~6 (Airport Coordination Procedures) standardizes slot
        coordination messaging on UTC --- the regulatory channel between
        airlines and slot coordinators uses UTC for all dates, days, and
        times.\footnote{IATA Standard Schedules Information Manual,
        Chapter 6 --- Airport Coordination Procedures.
        \url{https://www.iata.org/en/store/publications/manuals-standards-and-regulations/standard-schedule-information-manual-chapter-6airport-coordination-procedures__ssim-ch6/}}
  \item Flight planning and Air Traffic Control universally use UTC (``Zulu time'').
  \item Date-line crossings (e.g., TPE$\to$SFO) and daylight-saving transitions create
        ambiguity in local time and disappear as categories in UTC arithmetic.
  \item Cross-leg arithmetic ($\text{ETA}_f + \text{MCT} \leq \text{ETD}_g$) is direct in
        UTC; in local time it would require hop-by-hop conversion at every constraint.
\end{itemize}

\paragraph{Ingestion of SSIM-style sources.}
SSIM messages carry a Time Mode field with value \texttt{U} (UTC) or \texttt{L} (local);
the Local mode records include a paired UTC offset field
(\texttt{utc\_local\_time\_variation\_departure} and the arrival analog).\footnote{Octallium,
``Flight Schedules From SSIM'' (technical analysis, 2025).
\url{https://www.octallium.com/blog/2025/02/27/flight-schedules-from-ssim/}} At
ingestion, any record in Time Mode \texttt{L} is converted to UTC using the embedded
offset; only the resulting UTC timestamps enter the model. Where source data omits an
offset, ingestion fails fast with a structured error --- the model does not silently
infer a timezone.

\paragraph{Cutoffs and customs filing deadlines.}
The cutoff stack ($\text{CGC}_{f,\ctype}$, $\text{DCO}_f$, $\text{AMS}_f$,
$\text{ICS2}_f$, $\text{ACI}_f$ per leg) and the per-HAWB document-prep time
adjustment $\text{prep\_time}_k$ are all UTC timestamps. The graph generator
folds the stack into the per-MAWB-arc scalar $\text{CO}_a^*$ at build time;
the forward-time-window propagation (\S\ref{sec:fwd-time-propagation})
admits arcs against this scalar. The carrier-side cutoff is published in
local time at the origin airport; ingestion converts at intake.

\paragraph{Daylight Saving Time.}
IATA publishes seasonal schedules (Summer + Winter); UTC offsets in flight
records change with the season. Working in UTC absorbs DST entirely --- no DST
math anywhere in the model. Schedules near a DST transition simply carry the
post-transition UTC offset in the SSIM record, and ingestion produces correct
UTC timestamps either way.

\paragraph{Date-line crossings.}
A flight whose ETD and ETA fall on the same UTC date but different local dates
is unambiguous in UTC. SSIM handles this via the operational suffix `Z' for
flights using the same flight number on the same UTC date, and via the
Overnight Indicator. The MILP sees only UTC timestamps; the date-line is
invisible to constraint evaluation.

\paragraph{User-facing display.}
Operator dashboards and customer-facing screens convert UTC back to local time
at the relevant station (origin for cargo cutoff and pickup; destination for
ETA, deadline, and delivery). The conversion is a presentation-layer concern,
not a model concern.

%%--------------------------------------------------------------------
\section{Graph Construction (summary)}
\label{sec:graph-construction}
%%--------------------------------------------------------------------

The full graph-construction logic lives in
\texttt{model/air\_graph\_construction.md} (validated 2026-05-22). This section
records only what is load-bearing for the MILP that follows.

\paragraph{Symbols used in this section.}
\begin{table}[H]
\centering
\caption{Symbols used in the graph-construction section.}
\label{tab:graph-construction-syms}
\begin{tabular}{p{4.2cm}p{7.5cm}p{3cm}}
\toprule
Symbol & Brief & Defined \\
\midrule
$O,\, \text{CFS-O},\, \text{POL},\, H,$ $\text{CFS-H},\, \text{POD},\, \text{CFS-D},\, D$ & Eight node types (origin door $\to$ destination door) & Node types subsection \\
$A^{\text{MAWB}},\, A^{\text{coload}},\, A^{\text{ground}}$ & MAWB-eligible, co-load, and ground arc partitions & \S\ref{sec:variables} \\
$\mu_a$ & Realized air transit time on arc $a$ (hours; includes internal MCT on multi-leg arcs) & Per-arc scalars \\
$\text{CO}_a^*$ & Effective cargo cutoff at $\text{tail}(a)$ on POL-leaving air arcs: $\max(\text{DCO}, \text{AMS}, \text{ICS2}, \text{ACI}) - \text{prep\_time}_k$ & Per-arc scalars \\
$\text{ETA}_a$ & Scheduled arrival at $\text{head}(a)$ (UTC); anchors destination arrival (C.10a) & Per-arc scalars \\
$\text{STD}_a$ & Scheduled departure of air arc $a$ from $\text{tail}(a)$ (UTC); $\text{CO}_a^* \leq \text{STD}_a$ and $\text{ETA}_a = \text{STD}_a + \mu_a$ & Per-arc scalars \\
$\delta_a$ & Per-arc dwell on ground / connection arcs (hours) & Per-arc scalars, \S\ref{sec:ground-arc-params} \\
$\text{transit}(k, a)$ & Per-HAWB transit on arc $a$ & \S\ref{sec:variables} \\
$t^{\text{lo}}_n,\, t^{\text{hi}}_n$ & Propagated arrival-time window at node $n$ for HAWB $k$ & \S\ref{sec:fwd-time-propagation} \\
$t_k^{\text{rdy,early}},\, t_k^{\text{rdy,late}}$ & Cargo-ready window at origin & \S\ref{sec:hawb-params} \\
$T_k^{\text{abs}}$ & Latest valid arrival time (fallback-arc arrival; reachability bound) & \S\ref{sec:hawb-params} \\
$\Pi$ & ULD-interchange triples $(c_1, c_2, u)$ admissible under Cargo SPA / alliance pool & \S\ref{sec:through-uld} \\
$a^{\text{fb}}_k$ & Per-HAWB fallback arc $(O_k \to D_k^{\text{node}})$, always emitted, exempt from pre-filter & \S\ref{sec:fallback-arc} \\
$\text{arr\_dest}(k, a)$ & Precomputed destination-arrival scalar on terminal arcs (C.10a input) & \S\ref{sec:fallback-arc}, C.10a \\
$C^{\text{fallback}}$ & Tenant-global per-HAWB fallback dollar cost & \S\ref{sec:hawb-params} (Tenant-global) \\
$N_k,\, A_k$ & Per-HAWB node / arc set surviving pre-filter & \S\ref{sec:variables} \\
$k,\, a,\, n,\, f$ & HAWB / arc / node / flight indices & \S\ref{sec:variables} \\
\bottomrule
\end{tabular}
\end{table}

\subsection{Node types}

Eight node types: origin door \(O\), origin CFS \(\text{CFS-O}\), origin
gateway airport \(\text{POL}\), hub airport \(\text{H}\), hub CFS
\(\text{CFS-H}\), destination gateway airport \(\text{POD}\), destination CFS
\(\text{CFS-D}\), destination door \(D\). \(\text{CFS-O/D/H}\) may be
off-airport or on-airport; the difference is captured by the cartage-arc
time/cost. \(\text{CFS-H}\) exists \emph{only} at hubs where the forwarder (or
its agent) operates a warehouse --- so deconsolidation / re-consolidation is
possible only at those hubs; all other hubs allow carrier-side connections
only.

\subsection{Arc types}

\begin{table}[H]
\centering
\caption{Arc types in the O-D-arc graph.}
\label{tab:arc-types}
\begin{tabular}{lp{2.0cm}p{6.5cm}}
\toprule
Arc type & Carries MAWB? & Notes \\
\midrule
Pickup ($O \to \text{CFS-O}$) & no & road leg. \\
Origin CFS dwell & no & consolidation build-up. \\
Origin cartage ($\text{CFS-O} \to \text{POL}$) & no & near-zero time/cost if CFS is on-airport. \\
MAWB-arc ($\text{POL} \to \text{POD}/\text{H}$, $\text{H} \to \text{POD}$) & \textbf{yes} & forwarder-consolidated O-D segment under one offer; one MAWB per group $g$. \\
Co-load arc (air O-D segment) & no & per-kg billing on each HAWB's own chargeable weight. \\
Deconsolidation-dwell arc (at $\text{CFS-H}$) & no & break-and-rebuild between two MAWBs at a forwarder-operated hub CFS. \\
Destination cartage ($\text{POD} \to \text{CFS-D}$) & no & near-zero on-airport. \\
Destination CFS dwell & no & deconsolidation. \\
Customs clearance dwell & no & per-HAWB $\delta^{\text{cust}}_k$ between $\text{CFS-D}$ and final delivery. \\
Final delivery ($\text{CFS-D} \to D$) & no & road leg. \\
Fallback arc ($O \to D$) & no & per-HAWB direct origin-to-destination arc emitted always; bypasses pre-filter, capacity, cutoff; arrival at $T_k^{\text{abs}}$; cost $C^{\text{fallback}}$; selected only when no real route exists (\S\ref{sec:fallback-arc}). \\
\bottomrule
\end{tabular}
\end{table}

\subsection{Per-arc scalars precomputed by the graph generator}

For each air arc $a$:
\begin{itemize}[noitemsep]
  \item $\mu_a$ (hours) --- realized air transit time. For a multi-leg arc
        (multi-stop, same-carrier-connection, or interline through-MAWB)
        $\mu_a$ already includes internal MCT between the legs.
  \item $\text{CO}_a^*$ (h, UTC) --- effective cargo cutoff at $\text{tail}(a)$
        for air arcs leaving a POL: $\max(\text{DCO}, \text{AMS}, \text{ICS2},
        \text{ACI})$ minus per-HAWB prep time. The forward-time-window
        propagation (\S\ref{sec:fwd-time-propagation}) admits an incoming arc
        only when its propagated arrival lower bound at $\text{tail}(a)$ is
        $\leq \text{CO}_a^*$.
  \item $\text{ETA}_a$ (h, UTC) --- scheduled arrival at $\text{head}(a)$, the
        ETA of the last physical leg of $a$. Anchors destination arrival
        (C.10a).
  \item $\text{STD}_a$ (h, UTC) --- scheduled departure of $a$ from its first
        physical leg. Satisfies $\text{CO}_a^* \leq \text{STD}_a$ (cutoff
        precedes departure) and $\text{ETA}_a = \text{STD}_a + \mu_a$. Used in
        the equivalence proof (\S\ref{app:fwd-prop-equiv}); not a separate MILP
        input beyond $\text{CO}_a^*$, $\text{ETA}_a$.
\end{itemize}
Cross-MAWB transitions at a hub are realized by traversing a
deconsolidation-dwell arc, whose $\delta_a$ carries the deconsol / re-consol
time. Cross-carrier re-tendering at a hub without $\text{CFS-H}$ uses a
synthetic carrier-side connection-dwell arc emitted by the graph generator,
$\delta_a$ carrying the MCT. In all cases the dwell time enters the
forward-time-window propagation; no MILP MCT or time-propagation constraint
family exists.

\subsection{Two-phase construction}

Construction proceeds in two phases.
\begin{itemize}[noitemsep]
  \item \textbf{Phase 1 --- physical graph.} Nodes; all transport / dwell arcs;
        per-shipment forward-time-window propagation
        (\S\ref{sec:fwd-time-propagation}); per-shipment subgraph pre-filter
        (\S\ref{sec:prefilter}). Independent of consolidation logic --- no
        $g$, no MAWB objects. Validates transit times, cutoffs, dwells, and
        deadline reachability.
  \item \textbf{Phase 2 --- MAWB overlay.} Compute $g(k)$ for every HAWB.
        Instantiate MAWB candidates: one $(a, g)$ per MAWB-eligible air arc
        $a$ and each distinct $g$ present among HAWBs whose subgraph contains
        $a$. Co-load arcs are skipped --- they carry no MAWB.
\end{itemize}

\subsection{Arc enumeration policy on multi-leg flights (overlapping emission)}
\label{sec:arc-enumeration}

A multi-leg physical flight (e.g.\ $a \to b \to c$, freighter
TPE--ANC--JFK or pax-belly multi-stop) produces \emph{multiple} MAWB-arc
candidates in the graph, not one. The generator emits an arc for each
contiguous sub-segment for which the catalog provides an offer:
\begin{itemize}[noitemsep]
  \item single-leg arcs: $a \to b$, $b \to c$;
  \item the through arc $a \to c$ (when the carrier or interline pool
        publishes a through rate / through-AWB tariff over the full
        $a \to c$ O-D segment).
\end{itemize}
A 3-stop flight $a \to b \to c \to d$ analogously emits up to six
candidates: $\{a\!\to\!b,\, b\!\to\!c,\, c\!\to\!d,\, a\!\to\!c,\,
b\!\to\!d,\, a\!\to\!d\}$, with each present iff the catalog supplies an
offer for that sub-segment.

\paragraph{Why overlapping emission is load-bearing.}
The through arc $a \to c$ is \emph{not} the composition of the two
single-leg arcs $a \to b$ and $b \to c$. They differ on:
\begin{itemize}[noitemsep]
  \item \textbf{ULD / cargo handling at $b$.} Through-cargo on a single AWB
        rides through the hub without deconsolidation; segment-by-segment
        cargo deconsolidates at $b$'s CFS-H and rebuilds onto a new MAWB
        for $b \to c$. Different handling cost, different dwell, potentially
        different through-ULD policy under interchange
        (\S\ref{sec:through-uld-graph}).
  \item \textbf{Rate and pivot floors.} The through O-D segment can be sold
        under one BSA / TACT / GSA contract with one rate $\times$ one
        pivot for the full $a \to c$ chargeable weight, whereas
        single-leg arcs each carry their own contract, rate, and pivot
        floor independently. Two HAWBs riding the through arc share one
        through-MAWB; two HAWBs riding the single-leg arcs $a \to b$ and
        $b \to c$ activate two separate MAWBs each binding its own pivot.
  \item \textbf{Documentation.} The through arc represents one MAWB
        (one AWB number, one FWB/FHL filing pair). Segment-by-segment
        represents two MAWBs with separate filings.
\end{itemize}
The MILP picks among these emission options via $x_{k, a}$ on the candidate
arcs; HAWBs with different POD nodes naturally split across single-leg vs
through patterns inside the same physical flight schedule.

\paragraph{Worked example: TPE $\to$ HKG $\to$ LAX with two HAWBs.}
Single CX freighter rotation TPE--HKG--LAX. Forwarder operates a CFS-H at
HKG. Suppose the catalog supplies offers for three sub-segments:
\begin{itemize}[noitemsep]
  \item MAWB-arc 1: TPE $\to$ HKG. BSA contract $c_1$ on CX, rate $r_{c_1}$,
        pivot $\pi_1$.
  \item MAWB-arc 2: HKG $\to$ LAX. BSA contract $c_2$ on CX, rate
        $r_{c_2}$, pivot $\pi_2$.
  \item MAWB-arc 3: TPE $\to$ LAX (through). BSA contract $c_3$ on CX with
        through-rate $r_{c_3}$ and single pivot $\pi_3$. Through-ULD policy
        permits the ULD to remain intact at HKG.
\end{itemize}
Two HAWBs: HAWB-1 with $D_1^{\text{node}} = $ HKG door (final destination
HKG); HAWB-2 with $D_2^{\text{node}} = $ LAX door. Same consolidation
group $g$ for both.

The MILP can pick among (at least) two operationally distinct routings:
\begin{table}[H]
\centering
\caption{Routing combinations for two HAWBs on a multi-leg flight (TPE$\to$HKG$\to$LAX worked example).}
\label{tab:multi-leg-combos}
\begin{tabular}{lp{4.5cm}p{4.5cm}}
\toprule
Combo & HAWB-1 route & HAWB-2 route \\
\midrule
A (segment-by-segment) & $x_{1, \text{MAWB-arc 1}} = 1$, deconsols at HKG, ground to door. & $x_{2, \text{MAWB-arc 1}} = 1$; deconsol at HKG CFS-H; $x_{2, \text{MAWB-arc 2}} = 1$ from HKG; deconsol at LAX. \\
B (through for HAWB-2) & $x_{1, \text{MAWB-arc 1}} = 1$, deconsols at HKG, ground to door. & $x_{2, \text{MAWB-arc 3}} = 1$ through-AWB; rides intact through HKG; deconsol at LAX only. \\
\bottomrule
\end{tabular}
\end{table}

\paragraph{Cost comparison structure.}
\begin{itemize}[noitemsep]
  \item Combo A activates MAWB $(\text{arc 1}, g)$ shared by both HAWBs
        and MAWB $(\text{arc 2}, g)$ used by HAWB-2 alone. Two MAWBs
        instantiated; two pivot constraints
        $\text{pivot}_{\text{arc 1}, g}$ and $\text{pivot}_{\text{arc 2},
        g}$ each binding independently. HKG deconsol arc fires for HAWB-2:
        $\delta_a$ dwell, $c_a^{\text{flat}}$ rehandling charge.
  \item Combo B activates MAWB $(\text{arc 1}, g)$ for HAWB-1 alone and
        MAWB $(\text{arc 3}, g)$ for HAWB-2 alone. Two MAWBs instantiated;
        two pivot constraints, but the through-MAWB's pivot
        $\text{pivot}_{\text{arc 3}, g}$ binds on the \emph{full} TPE-LAX
        through-chargeable-weight against $\pi_3$. HAWB-2 skips the HKG
        deconsol arc (no $\delta_a$, no rehandling charge); only the
        flight's internal MCT is absorbed into $\mu_{\text{arc 3}}$.
\end{itemize}
The MILP picks based on total cost. Combo B wins when through-rate $r_{c_3}
\cdot \text{CW}_{\text{arc 3}, g_2}$ + saved HKG handling beats Combo A's
$r_{c_1} \cdot \text{CW}_{\text{arc 1}, g_2} + r_{c_2} \cdot
\text{CW}_{\text{arc 2}, g_2} + $ HKG deconsol cost. Combo A wins when
HAWB-1 and HAWB-2 share enough chargeable weight on the TPE-HKG leg that
density-mixing on MAWB-arc 1 (\S\ref{sec:cw-density-mixing}) brings
$\text{CW}_{\text{arc 1}, g}$ below the sum of per-HAWB chargeable weights
that Combo B would charge.

\paragraph{Pivot scope --- one pivot per MAWB-arc, not per leg.}
Each MAWB-arc carries one pivot floor (\S\ref{sec:bsa-params}). The
through arc $a \to c$ is contractually one offer with one rate against
one pivot covering the full O-D segment; \emph{not} two pivots tied to
the underlying legs. Per-leg flexibility is achieved by the graph
generator emitting separable single-leg arcs as parallel candidates, not
by subdividing a through arc's pivot. This is the structural reason
$\text{pivot}_{a, g}$ is indexed by $(a, g)$ and not by $(a, g, f)$ over
legs.

\subsection{Forward-time-window propagation (replaces MILP time constraints)}
\label{sec:fwd-time-propagation}

Arrival-time feasibility through the per-HAWB subgraph is enforced
\emph{at graph build}, not in the MILP. For each HAWB $k$, the generator
maintains an arrival-time window $[t^{\text{lo}}_n,\, t^{\text{hi}}_n]$ at
every node $n$ reachable from $O_k$, and admits arcs only when the
propagated window is consistent with any hard time window at the arc's head.

\paragraph{Window propagation rules.}
The propagation iterates over arcs, not nodes: each candidate outbound arc
$a'$ is admitted or rejected based on the propagated window at its tail and
its own per-arc time bound. This per-outbound admission is load-bearing ---
a POL node with multiple outbound flights (each with a different
$\text{CO}_{a'}^*$) is a common case, and a single per-node state cannot
represent the per-outbound admission criterion without losing arcs.
\begin{itemize}[noitemsep]
  \item \textbf{Origin.} $[t^{\text{lo}}_{O_k},\, t^{\text{hi}}_{O_k}] =
        [t_k^{\text{rdy,early}},\, t_k^{\text{rdy,late}}]$.
  \item \textbf{Non-air arc $a = (n_1, n_2)$ (ground / dwell).} Always
        admitted (no cutoff). Propagate $[t^{\text{lo}}_{n_2},\,
        t^{\text{hi}}_{n_2}] \leftarrow [t^{\text{lo}}_{n_1},\,
        t^{\text{hi}}_{n_1}] + \text{transit}(k, a)$.
  \item \textbf{Air arc $a = (n_1, n_2)$ leaving a POL node $n_1$ (cargo-cutoff admission).}
        Admit $a$ only if $t^{\text{lo}}_{n_1} \leq \text{CO}_a^*$
        (the propagated earliest arrival at the cutoff node clears the cargo
        cutoff). The check uses $\text{CO}_a^*$ of \emph{this} air arc; a
        sibling outbound air arc $a''$ leaving the same $n_1$ is admitted /
        rejected independently against its own $\text{CO}_{a''}^*$.
  \item \textbf{Dispatch-feasibility check on candidate outbound air arc.}
        For each candidate outbound air arc $a$ leaving a POL, the
        per-HAWB dispatch lead time $\lambda^{\text{disp}}_k$ enforces a
        backward-binding deadline at the origin door: admit $a$ only if
        $t_k^{\text{rdy,early}} \leq \text{CO}_a^* - \lambda^{\text{disp}}_k$.
        Equivalently, the truck must be able to depart the origin no later
        than $\text{CO}_a^* - \lambda^{\text{disp}}_k$ for the dispatch
        choreography (truck transit, CFS receiving, build, tender) to land
        cargo at the ramp by $\text{CO}_a^*$. This complements the
        forward-propagation cutoff check above: forward propagation
        validates that ground-arc dwells add up to a feasible arrival at
        the POL; the dispatch check validates that the truck has enough
        \emph{head-room} from cargo-ready to leave on time. A real-route
        air arc $a$ is admitted iff \emph{both} checks pass.
  \item \textbf{Air-arc head propagation (scheduled-arrival fix).} If $a$ is
        admitted, the arrival window at $\text{head}(a)$ along this branch is
        \emph{pinned} to the scheduled ETA: $[t^{\text{lo},(a)}_{n_2},\,
        t^{\text{hi},(a)}_{n_2}] = [\text{ETA}_a,\, \text{ETA}_a]$ ---
        \emph{not} $t^{\text{lo}}_{n_1} + \text{transit}$. The flight departs on
        schedule regardless of how early the cargo cleared the cutoff, so an
        air arc \emph{resets} the running time rather than accumulating it.
        This pinning is load-bearing for the equivalence in
        \S\ref{app:fwd-prop-equiv}: without it, early-arriving cargo would
        propagate a fictitiously early downstream time and over-admit arcs.
        Branching: each admitted outbound
        arc $a$ starts an independent downstream propagation from its head.
        Equivalently the propagation labels are per-(node, incoming-arc), not
        per-node alone. Where multiple incoming branches reach the same
        downstream node, the windows are unioned only for reachability checks
        (Destination reachability paragraph below), not for cutoff admission.
\end{itemize}
\paragraph{Equivalence to a MILP big-M time-propagation family.}
Per-outbound admission as defined above is mathematically equivalent, for
deterministic transits, to a big-M MILP with node-time variables
$t_k^n \in [t_k^{\text{rdy,early}}, T_k^{\text{abs}}]$ in which \emph{every}
propagation row is gated by the arc-selection variable $x_{k,a}$: each
\emph{ground/dwell} arc contributes $t_k^{n_2} \geq t_k^{n_1} +
\text{transit}(k, a) - M(1 - x_{k,a})$, and each \emph{air} arc $a$ at a POL
contributes a per-outbound cutoff gate $t_k^{n_1} \leq \text{CO}_a^* +
M(1 - x_{k,a})$ together with a gated scheduled-arrival fix $\text{ETA}_a -
M(1 - x_{k,a}) \leq t_k^{n_2} \leq \text{ETA}_a + M(1 - x_{k,a})$. Gating
\emph{all} rows (not just the cutoff) is required: an ungated air-head
equality $t_k^{n_2} = \text{ETA}_a$ would force a single $\text{ETA}$ on any
node that is the head of two candidate flights, and ungated ground rows would
couple unselected arcs' tail times --- both spurious. The scheduled-arrival
fix itself is also required: replacing it with a generic block-time row
$t_k^{n_2} \geq t_k^{n_1} + \text{transit}(k,a)$ lets early cargo propagate a
fictitiously early downstream time and over-admits arcs. Under this gated
formulation, the projection onto the route variables $x_k$ of the region cut
out by the time rows equals the graph-admitted route set, and the graph-build
version simply computes that projection to the feasible-arc set $A_k$ by
forward sweep rather than enumerating MILP rows. The proof, with the
scheduling-consistency assumption it needs ($\text{CO}_a^* \leq \text{STD}_a$,
$\text{ETA}_a = \text{STD}_a + \mu_a$), is in \S\ref{app:fwd-prop-equiv}.

\paragraph{Cross-MAWB hub transitions.}
Cross-MAWB transitions at a hub traverse a deconsolidation-dwell arc whose
$\delta_a$ carries the deconsol / re-consol time. Non-CFS-H hubs use a
synthetic carrier-side connection-dwell arc with $\delta_a = \text{MCT}$. In
both cases the dwell is included in the forward-propagation step above.

\paragraph{Destination reachability.}
An arc $a$ is retained in $A_k$ only if a forward path from $\text{head}(a)$
to $D_k^{\text{node}}$ exists with cumulative arrival $\leq T_k^{\text{abs}}$.
Arcs whose earliest-feasible end-to-end arrival at $D_k$ exceeds
$T_k^{\text{abs}}$ are dominated by the fallback arc
(\S\ref{sec:fallback-arc}) and may be pruned. \textbf{Pruning-safety
condition:} the dominance holds iff $C^{\text{fallback}}$ for this HAWB
satisfies the sizing rule in \S\ref{sec:hawb-params} (Tenant-global
parameters), which guarantees the fallback's combined cost ---
$C^{\text{fallback}} + W_k (T_k^{\text{abs}} - \Delta_k)^2$ --- exceeds
any realistic real-route cost at arrival $> T_k^{\text{abs}}$. The
per-tenant sizing formula makes this a per-HAWB property:
high-$W_k$ HAWBs (high-value, high service-product weight) get a larger
$C^{\text{fallback}}$, low-$W_k$ HAWBs get smaller, both maintaining
dominance. If a tenant under-sizes $C^{\text{fallback}}$ (sets it
globally below the per-HAWB max), late-arrival real routes may sneak
through pruning; the walking-skeleton metric
\texttt{fractional\_fallback\_usage} surfaces this regression.

\paragraph{Locked-prefix partial-locks.}
For partial-locked HAWBs (\S\ref{sec:locked-commitments}), the origin is
re-pointed to the current observed node with $[t^{\text{lo}}_{O_k},\,
t^{\text{hi}}_{O_k}] = \{t_k^{\text{obs}}\}$ (a singleton); propagation
proceeds from there.

\paragraph{Stochastic transit (P1 hook).}
Under deterministic arc times the propagated window is a deterministic
interval. The Finding S Ch~1 hook for stochastic transit \emph{does not}
naively lift to per-arc quantile arithmetic on the propagation: the sum of
per-arc P-quantiles is not the P-quantile of the sum (convolution
problem), so propagating $Q_p(\text{transit}(k, a))$ at each step
under-estimates end-to-end uncertainty. The supported P1 path is to
\emph{precompute end-to-end P-quantile per terminal arc} into
$\text{arr\_dest}(k, a)$ (\S\ref{sec:fallback-arc}, C.10a), leaving the
forward propagation on a safe deterministic statistic (e.g.\ mean transit)
used only for cutoff and reachability admission. The MILP still sees no
variance; see the cross-reference in
item~\ref{item:tt-quantile-binding} of \S\ref{sec:deferred}.

\paragraph{Why this replaces a constraint family.}
Pushing time-feasibility enforcement to graph build removes the
time-propagation variables ($t_k^n$ for intermediate $n$), the
time-propagation constraint family ($\sim 1{,}500$ rows at base scale), and
the cutoff constraint family ($\sim 300$ rows). Destination arrival enters
the MILP only as a definition equation in C.10a. The mathematical
equivalence (for deterministic transits) to the corresponding MILP big-M
family is detailed in the ``Equivalence to a MILP big-M time-propagation
family'' paragraph above.

\subsection{Fallback arc and feasibility guarantee}
\label{sec:fallback-arc}

For every HAWB $k \in K$ the graph generator emits one additional arc, the
\emph{fallback arc}:
\begin{equation*}
  a^{\text{fb}}_k = (O_k,\, D_k^{\text{node}}) \in A_k,
  \quad \text{tagged with arc type } \texttt{fallback}.
\end{equation*}
It is a direct edge from the origin door to the destination door, present in
every HAWB's subgraph regardless of pre-filter outcome. Its sole purpose is to
guarantee that the MILP is always feasible --- there is no input data
combination for which the optimizer can return \texttt{INFEASIBLE}.

\paragraph{Per-arc scalars on the fallback arc.}
\begin{itemize}[noitemsep]
  \item \textbf{Transit time:} $\text{transit}(k,\, a^{\text{fb}}_k) =
        T_k^{\text{abs}} - t_k^{\text{rdy,early}}$ (precomputed scalar).
        Forward propagation along $a^{\text{fb}}_k$ lands
        $t^{\text{lo}}_{D_k^{\text{node}}} = T_k^{\text{abs}}$ at the
        destination, consistent with the arrival-time scalar below.
  \item \textbf{Arrival time:} $\text{arr\_dest}(k,\, a^{\text{fb}}_k) =
        T_k^{\text{abs}}$ (precomputed scalar; enters the C.10a destination-arrival
        definition). The fallback arc fires C.10 tardiness for the gap
        $T_k^{\text{abs}} - \Delta_k$, contributing a defined, finite quadratic
        penalty $W_k \cdot (T_k^{\text{abs}} - \Delta_k)^2$ on top of
        $C^{\text{fallback}}$.
  \item \textbf{Dollar cost:} $C^{\text{fallback}}$, a tenant-config
        constant sized per the formula in \S\ref{sec:hawb-params}
        (Tenant-global parameters). Reported separately in the diagnostic
        output (\S\ref{sec:output-diagnostics}) so its magnitude does not
        drown the real routing cost in operator-facing screens.
\end{itemize}

\paragraph{Exemptions.}
The fallback arc \emph{bypasses}:
\begin{itemize}[noitemsep]
  \item the eight pre-filter predicates (\S\ref{sec:prefilter});
  \item per-contract allotment (C.5), per-ULD physical cap (C.5b), per-offer
        cap (C.5c) --- the fallback carries no MAWB and consumes no capacity;
  \item MAWB instantiation linkage (C.2) --- the fallback is not a MAWB-arc;
  \item the chargeable-weight aggregation (C.4) and bucket cost terms in the
        objective (no $z_{a^{\text{fb}}_k, g}$, no $\text{CW}$ on this arc).
\end{itemize}
It does participate in C.1 flow conservation like any other arc, and in the
C.10a destination-arrival definition via $\text{arr\_dest}(k, a^{\text{fb}}_k)
= T_k^{\text{abs}}$.

\paragraph{Per-shipment subgraph pre-filter interaction.}
The pre-filter (\S\ref{sec:prefilter}) applies to \emph{real} arcs only. If
the real arcs surviving pre-filter form an empty set ($A_k \setminus
\{a^{\text{fb}}_k\} = \emptyset$), the graph generator logs a pre-solve
warning (``HAWB $k$ has zero real arcs surviving pre-filter; will route via
fallback'') --- the MILP confirms by selecting $x_{k, a^{\text{fb}}_k} = 1$.
No pre-solve infeasibility branch exists.

\paragraph{Operational interpretation per cargo class.}
\begin{itemize}[noitemsep]
  \item \textbf{PER (and AVI, HUM, VAL).} Fallback selection signals that the
        HAWB is unplannable given the current supply set and human
        intervention is required (e.g.\ procure additional capacity, adjust
        the cargo-ready window with the shipper, escalate to a co-loader,
        invoke a spot booking outside the catalog). Cargo would not in fact
        be carried via the fallback arc; the arc is a planning artifact
        flagging the rescue case.
  \item \textbf{GEN.} $T_k^{\text{abs}}$ is the customer-cancellation point
        (the time past which the shipment loses commercial value or is
        refunded). Same operational semantics: fallback selection is a
        rescue-event signal, not an actual route.
\end{itemize}
In both cases the orchestrator inspects the post-solve fallback set
(\S\ref{sec:output-diagnostics}) and triggers structured rescue handling.

\subsection{Through-ULD policy --- absorbed at graph build}
\label{sec:through-uld-graph}

ULD interchange (Star / SkyTeam / oneworld pool, bilateral Cargo SPA) does not
appear as a MILP variable. It \emph{influences} how the graph generator emits
arcs: a same-AWB through-segment that crosses carriers under interline is a
\emph{single} MAWB-arc (one rate, one MAWB), with $\mu_a$ including internal
MCT and the per-leg cutoff stack collapsed into $\text{CO}_a^*$. Where no
interchange agreement exists, the generator emits two separate MAWB-arcs
joined by a synthetic carrier-side connection-dwell arc carrying the MCT. The
\emph{cost} difference (rehandling fees, re-stuffing) is captured in the
handling charge $c_a^{\text{flat}}$ of the joining ground arc (deconsol-dwell
arc or carrier-side connection arc); see \S\ref{sec:through-uld} for the
full case enumeration and the per-arc data the graph generator emits.

%%--------------------------------------------------------------------
\section{Commercial Structure: MAWB and HAWB}
\label{sec:mawb-hawb}
%%--------------------------------------------------------------------

Air freight is contracted, billed, and routed at two document layers ---
Master Air Waybill (MAWB) and House Air Waybill (HAWB). The economic mechanism
mid-market forwarders depend on (consolidation of many small HAWBs into one
MAWB to access lower weight-break rates) is invisible if the model treats every
shipment as an independent tender. This section defines the commercial
structure the formulation must represent.

\subsection{Definitions}

\begin{description}[leftmargin=2cm, labelwidth=1.8cm, style=nextline]
  \item[MAWB] \emph{Master Air Waybill.} The contract of carriage between the
        airline (or interline group under a Cargo SPA) and the freight
        forwarder for a contiguous O-D segment under one offer. 11-digit
        number, 3-digit IATA airline prefix (e.g., \texttt{160-12345678} for
        Cathay). The airline does not see what is inside the MAWB at the
        pricing level.
  \item[HAWB] \emph{House Air Waybill.} The contract of carriage between the
        freight forwarder and the underlying shipper. One HAWB per (shipper,
        consignee) pair. Issued under the forwarder's own numbering scheme.
        Airline does not see HAWBs in pricing, but customs at destination does
        (AMS / ICS2 / equivalent filings are per-HAWB).
\end{description}

\paragraph{Filing mechanics.}
The MAWB is filed electronically with the airline via IATA Cargo-IMP
(\texttt{FWB} = waybill data, \texttt{FHL} = list of HAWBs underneath) or
Cargo-XML / ONE Record. The cargo manifest lists all HAWBs under a given MAWB
and travels with the physical cargo. At destination, customs receives both the
MAWB (carriage record) and each HAWB (as the underlying consignment requiring
clearance).

\subsection{MAWB \(=\) (MAWB-arc, group)}
\label{sec:mawb-arc-g}

A MAWB is one carriage contract. In the routing graph it is identified by
\textbf{$(\text{MAWB-arc},\, g)$} --- one MAWB per (MAWB-arc, consolidation
group). The HAWBs assigned to a MAWB-arc partition into MAWBs by their group.
The index set $M = \{(a, g) : a \in A^{\text{MAWB}},\, g \in G_a\}$ is concrete
(enumerable arcs $\times$ enumerable groups) --- \emph{no symmetry, no
MAWB-count decision}.

A new MAWB begins wherever any of the following changes mid-journey:
\begin{enumerate}[noitemsep,label=(\alph*)]
  \item the offer / contract changes (BSA to TACT, or one BSA to another);
  \item the carrier changes \emph{without} interline (no Cargo SPA / MITA
        binding the two carriers under one AWB);
  \item the forwarder deconsolidates at a hub with $\text{CFS-H}$ (B3, C6, C7
        in \texttt{air\_graph\_construction.md \S6});
  \item co-consolidated HAWBs diverge to different downstream MAWB-arcs.
\end{enumerate}
Interline (Cargo SPA / MITA; alliance Cargo pools where applicable) lets one
AWB span carriers --- so the \emph{boundary is an offer boundary, not a carrier
boundary}.

\subsection{The consolidation group \(g(k)\)}
\label{sec:g-of-k}

Two HAWBs may share a MAWB only if they are on the \emph{same MAWB-arc} and
$g(k_1) = g(k_2)$. $g$ is a deterministic, single-valued function of a fixed
tuple of HAWB attributes:
\[
  g(k) =
    \begin{cases}
      \bigl(\text{cargo\_class}(k),\,\text{temperature}(k)\bigr)
        & \text{if } \text{cargo\_class}(k) \in \{\text{GEN, DGR, PER}\}, \\
      \bigl(\text{cargo\_class}(k),\,\text{HAWB-id}(k)\bigr)
        & \text{if } \text{cargo\_class}(k) \in \{\text{VAL, HUM, AVI}\}.
    \end{cases}
\]
$\text{VAL}$, $\text{HUM}$, $\text{AVI}$ are non-consolidable --- $g$ carries
the HAWB id, so each forms a singleton group and gets its own dedicated MAWB.

\paragraph{Partition by construction.}
The set of $g$ values induces a partition of the HAWBs: the groups are
\emph{pairwise disjoint and exhaustive}, automatically, because $g$ is a
single-valued total function. The weaker rule ``no group is a subset of
another'' is insufficient (it permits partial overlap); pairwise disjoint is
the correct rule. See \texttt{air\_graph\_construction.md \S4.3} for the
discussion. Adopt the discipline: $g$ is always the attribute tuple, never an
ad-hoc predicate.

\paragraph{Scope of the partition rule.}
It governs group \emph{definitions}. It does not apply to MAWB-membership sets
\emph{across} a journey: a downstream MAWB's HAWB set being a strict subset of
an upstream MAWB's (cargo peeling off at a hub) is normal and correct --- this
is cases C2, C6, C7 in the graph-construction doc.

\paragraph{DGR coarseness.}
All DG cargo lives in one DGR group at the MAWB layer. True pairwise DG
segregation is a ULD-layer concern (IATA segregation tables; non-transitive),
not a partition over groups. Fine-grained DG segregation is deferred
(\S\ref{sec:deferred}).


\subsection{Bucket chargeable weight: density mixing}
\label{sec:cw-density-mixing}

The airline applies the volumetric rule at the \emph{MAWB} level, not per
HAWB. This subsection defines the per-MAWB chargeable-weight aggregate
$\text{CW}_{a,g}$ that drives all rate-card billing on MAWB-eligible arcs,
and shows why aggregating physical quantities \emph{before} taking the
actual-vs.-volumetric $\max$ is the structural reason consolidation pays.

\paragraph{Nomenclature.}
Symbols used in this subsection. Decision variables and parameters are
formally collected in \S\ref{sec:variables} and \S\ref{sec:parameters}; they
appear in the table here so the subsection is self-contained.
\begin{table}[H]
\centering
\caption{Symbols used in the chargeable-weight density-mixing subsection.}
\label{tab:density-mixing-syms}
\begin{tabular}{p{2.5cm}p{2.2cm}p{8.4cm}}
\toprule
Symbol & Role & Meaning \\
\midrule
$k$ & index & HAWB (commodity / shipment) \\
$a$ & index & arc in the O-D-arc graph \\
$g$ & index & consolidation group key (\S\ref{sec:g-of-k}) \\
$(a, g)$ & pair & a MAWB candidate: one MAWB-arc carrying one consolidation group \\
$M$ & set & all MAWB candidates: $M = \{(a, g) : a \in A^{\text{MAWB}},\, g \in G_a\}$ (\S\ref{sec:sets}) \\
$A_k$ & set & pre-filtered subgraph of HAWB $k$ --- arcs surviving its per-shipment eligibility filter (\S\ref{sec:prefilter}) \\
$K_a$ & set & HAWBs whose subgraph contains $a$: $K_a = \{k : a \in A_k\}$ \\
$g(k)$ & function & consolidation group of HAWB $k$ (single-valued total function; \S\ref{sec:g-of-k}) \\
$w_k$ & param (kg) & gross actual weight of HAWB $k$ \\
$v_k$ & param (m\textsuperscript{3}) & physical volume of HAWB $k$ \\
$\varepsilon$ & param (dimensionless) & dunnage / strapping factor; default $\approx 0.05$ (i.e.\ $5\%$ uplift) \\
$167$ & constant (kg/m\textsuperscript{3}) & IATA volumetric divisor (TACT \S3.9): kg-per-m\textsuperscript{3} conversion at which volume becomes billable weight \\
$x_{k,a}$ & binary var & $1$ iff HAWB $k$ is routed through arc $a$; defined for $k \in K$,\, $a \in A_k$ \\
$\text{Wt}_{a,g}$ & cont.\ var (kg, $\geq 0$) & MAWB-level aggregate \emph{actual} weight with dunnage uplift \\
$\text{Wv}_{a,g}$ & cont.\ var (kg, $\geq 0$) & MAWB-level aggregate \emph{volumetric} weight (kg-equivalent) \\
$\text{CW}_{a,g}$ & cont.\ var (kg, $\geq 0$) & MAWB-level chargeable weight; billing reference for all rate families that bill on the MAWB \\
\bottomrule
\end{tabular}
\end{table}

\paragraph{Aggregates and chargeable-weight bound.}
For each MAWB $(a, g) \in M$:
\begin{align}
  \text{Wt}_{a,g} &\defeq (1 + \varepsilon) \sum_{k \in K_a : g(k) = g} w_k \cdot x_{k,a}
    \label{eq:mawb-wt} \\
  \text{Wv}_{a,g} &\defeq \sum_{k \in K_a : g(k) = g} v_k \cdot 167 \cdot x_{k,a}
    \label{eq:mawb-wv} \\
  \text{CW}_{a,g} &\geq \text{Wt}_{a,g}, \qquad
  \text{CW}_{a,g} \geq \text{Wv}_{a,g}
    \label{eq:mawb-cw}
\end{align}

\paragraph{Eq.~\ref{eq:mawb-wt} --- actual-weight aggregate $\text{Wt}_{a,g}$.}
Sums the gross weight $w_k$ over every HAWB that (i) has arc $a$ in its
pre-filtered subgraph ($k \in K_a$), (ii) belongs to consolidation group $g$
($g(k) = g$), and (iii) is actually routed through arc $a$ in this solution
($x_{k,a} = 1$). The $(1 + \varepsilon)$ multiplier adds the dunnage uplift
the airline charges on top of pure cargo weight --- kraft paper, edge
protectors, strapping, shrink wrap, ULD-mate dunnage. The summation indexing
$k \in K_a : g(k) = g$ enforces the partition: each HAWB on arc $a$
contributes to exactly one MAWB-aggregate, namely its own group's. There is
no $h_{k,(a,g)}$ assignment binary --- group membership is a function lookup,
not a decision.

\paragraph{Eq.~\ref{eq:mawb-wv} --- volumetric-weight aggregate $\text{Wv}_{a,g}$.}
Sums each HAWB's volume converted to a kg-equivalent at the IATA divisor of
$167\,\text{kg/m}^3$ (TACT \S3.9). This is the chargeable-weight floor an
airline applies to light, bulky cargo: a $1\,\text{m}^3$ shipment is billed
as if it weighed $167\,\text{kg}$ regardless of actual weight. The same
summation guard ($k \in K_a : g(k) = g$, weighted by $x_{k,a}$) applies.
Dunnage is \emph{not} multiplied into $\text{Wv}_{a,g}$: strapping adds
weight but does not displace appreciable cargo volume in industry billing
convention.

\paragraph{Eq.~\ref{eq:mawb-cw} --- chargeable-weight lower bounds.}
Two inequalities together force $\text{CW}_{a,g} \geq \max(\text{Wt}_{a,g},\,
\text{Wv}_{a,g})$. The airline bills on the larger of physical and
volumetric --- it charges on whichever yields more revenue. The bound is
written as two inequalities rather than an explicit $\max(\cdot, \cdot)$
operator for linearization: a max would otherwise require an auxiliary
binary (case split) or a PWL formulation. The inequality form is exact when
paired with the monotonicity invariant in the objective; see next paragraph.

\paragraph{Why the inequality relaxation is exact.}
Eq.~\ref{eq:mawb-cw} bounds $\text{CW}_{a,g}$ from below only --- nothing in
this constraint family forces $\text{CW}$ down to the max. The pin comes
from the objective: every cost term in the model that involves
$\text{CW}_{a,g}$ has a non-negative coefficient (the \emph{monotonicity
invariant}; see \S\ref{sec:objective}). Because the objective is a
minimization, the optimizer drives $\text{CW}_{a,g}$ to its lower envelope
--- which is exactly $\max(\text{Wt}_{a,g}, \text{Wv}_{a,g})$. No binary is
required for the $\max$. If a future rate-family extension introduces a
\emph{negative}-coefficient term on $\text{CW}_{a,g}$ (rebate, marketing
credit, volume-discount kicker), the monotonicity invariant breaks and the
inequalities must be tightened to equality via a PWL formulation; a
post-solve assertion suite checks this invariant.

\paragraph{Density mixing benefit.}
Aggregating physical quantities \emph{before} the $\max$ lets dense cargo
absorb the volumetric slack of light cargo, recovering chargeable-weight
savings that a per-HAWB billing model cannot see. Worked example: two HAWBs
on one MAWB.
\begin{table}[H]
\centering
\caption{Density-mixing worked example: two HAWBs on one MAWB.}
\label{tab:density-mixing-example}
\begin{tabular}{lcccc}
\toprule
HAWB & Actual wt & Volume & Volumetric ($\times 167$) & per-HAWB CW \\
\midrule
$k_1$ (machinery) & 200 kg & 0.5 m\textsuperscript{3} & 84 kg & 200 kg \\
$k_2$ (apparel)   & 60 kg  & 1.2 m\textsuperscript{3} & 200 kg & 200 kg \\
\midrule
\textbf{Naïve sum of per-HAWB CW} & & & & \textbf{400 kg} \\
\textbf{MAWB-level} ($\varepsilon = 0$): $\text{Wt} = 260$, $\text{Wv} = 284$ & & & & \textbf{284 kg} \\
\bottomrule
\end{tabular}
\end{table}
Per-HAWB billing charges each shipment at its own $\max(w_k, 167 v_k)$:
$\max(200, 84) + \max(60, 200) = 400\,\text{kg}$. MAWB-level billing
aggregates first: $\text{Wt} = 200 + 60 = 260$,\, $\text{Wv} = 84 + 200 =
284$,\, $\text{CW} = \max(260, 284) = 284\,\text{kg}$. The machinery's
spare volume ``soaks up'' the apparel's volumetric surplus inside the same
MAWB. The $116\,\text{kg}$ saving ($\approx 29\%$) is captured only by
aggregating before the $\max$, and is the core reason consolidation is
modeled at the MAWB level, not per HAWB.

%%--------------------------------------------------------------------
\section{Parameters}
\label{sec:parameters}
%%--------------------------------------------------------------------

\paragraph{Currency convention.}
All monetary parameters in this section --- rate-card amounts, contract
floors, surcharges, ground-arc costs --- are denominated in a single canonical
settlement currency, USD. Source rates routinely arrive in other currencies:
per IATA TACT rules, air cargo rates are quoted in the local currency of the
origin country, and local ground charges, customs fees, and surcharges follow
their country of incurrence. Each rate record in the supply and surcharge
catalogs carries an explicit \texttt{currency} field; values are converted to
USD at solve time through a single FX rate table. The conversion mechanism is
specified in \texttt{data\_model.md~\S7} and is not restated here. Every
constraint and objective term operates on post-conversion USD values.

\subsection{Per-HAWB parameters}
\label{sec:hawb-params}

\begin{table}[H]
\centering
\caption{Per-HAWB parameters.}
\label{tab:hawb-params}
\begin{tabular}{p{3.5cm}p{2.2cm}p{6.5cm}p{1.5cm}}
\toprule
Parameter & Symbol & Description & Unit \\
\midrule
Volume & $v_k$ & Cargo volume (sum over all pieces) & m\textsuperscript{3} \\
Actual weight & $w_k$ & Gross weight & kg \\
Chargeable weight (per-HAWB billing reference, not used in bucket capacity) & $cw_k$ & $\max(w_k,\, v_k \cdot 167)$ & kg \\
Piece dimensions & $\text{pieces}(k)$ & Max single-piece $L \times W \times H$ for ULD-fit pre-filter & cm \\
Cargo-ready window (early) & $t_k^{\text{rdy,early}}$ & Earliest moment cargo is available at the origin door & h (UTC) \\
Cargo-ready window (late) & $t_k^{\text{rdy,late}}$ & Latest moment within the shipper's window & h (UTC) \\
Per-HAWB origin-side dispatch lead time & $\lambda^{\text{disp}}_k$ & Backward-binding offset: the truck must depart the origin door no later than $\text{CO}_a^* - \lambda^{\text{disp}}_k$ to make the cutoff $\text{CO}_a^*$ of any candidate outbound air arc $a$. Accumulates: pickup-to-CFS road transit + origin-CFS receiving + ULD build + tender-to-ramp lead time + queue contingency. Tenant-calibrated lookup keyed on (origin region, cargo type, MAWB consolidation profile); typical 4--8h on mid-sized origin lanes, longer for DGR / VAL builds. Used in the forward-time-window propagation's pickup-feasibility check (\S\ref{sec:fwd-time-propagation}). & h \\
Latest valid arrival time (``hard backstop'') & $T_k^{\text{abs}}$ & \textbf{Mandatory-finite for every HAWB.} Semantically the cargo-death point for PER / AVI / HUM / VAL (spoilage, expiration, security window) and the customer-cancellation point for GEN (commercial value lost; shipment refunded). \emph{Not} a hard MILP constraint --- it is the arrival time of the per-HAWB fallback arc (\S\ref{sec:fallback-arc}) and bounds the worst-case quadratic tardiness penalty (C.10). For GEN, ingestion applies a tenant-configured default (e.g.\ $t_k^{\text{rdy,early}} + 30$ d) when the shipper does not specify. \textbf{Pre-MILP guard:} ingestion asserts $T_k^{\text{abs}}$ is finite and $\geq \Delta_k$; failure surfaces as a structured intake error, not silent corruption of the C.14 $\tau_k$ bound or the fallback $\text{arr\_dest}$ scalar. & h (UTC) \\
Service-product soft deadline & $T_k^{\text{SLA}}$ & $= t_k^{\text{rdy,early}} + T^{\text{SLA}}_{\text{sp}(k)}$ (\S\ref{sec:service-products}) & h (UTC) \\
Shipper-contractual deadline & $T_k^{\text{dead}}$ & Optional shipper-contracted door-to-door deadline; absent when the service product alone governs. Used only as one of the two terms in $\Delta_k$ (next row). & h (UTC) \\
Effective soft deadline & $\Delta_k$ & $\min(T_k^{\text{dead}},\, T_k^{\text{SLA}})$ if $T_k^{\text{dead}}$ exists; otherwise $\Delta_k = T_k^{\text{SLA}}$. Used in C.10. & h (UTC) \\
Declared shipment value & $\text{value}_k$ & Shipper-declared cargo value used to scale the tardiness penalty. Required at ingestion. & USD \\
Shipment-value coefficient & $\mu_k$ & $\mu_k = \text{value}_k / V^{\text{ref}}$ (dimensionless). Scales tardiness penalty in proportion to declared value. & --- \\
Combined tardiness weight & $W_k$ & $W_k = w_{\text{sp}(k)} \cdot \mu_k$. Per-HAWB coefficient in the quadratic tardiness term $W_k \cdot \tau_k^2$ (C.10, \S\ref{sec:lin-tardiness}). & USD / h\textsuperscript{2} \\
Origin door node & $O_k \in N_k$ & --- & --- \\
Destination door node & $D_k^{\text{node}} \in N_k$ & --- & --- \\
Port of discharge node & $\text{POD}_k \in N_k$ & HAWB-specific POD --- the air-arc head node where $k$'s air leg lands. Used in C.10a to define the terminal-arc set $\mathcal{A}^{\text{last}}_k$. Determined by the chosen route; multiple POD candidates may exist in $A_k$ pre-solve. & --- \\
Post-POD ground tail transit sum & $\Delta^{\text{post}}_k$ & Sum of ground-arc dwells from $\text{POD}_k$ to $D_k^{\text{node}}$ on $k$'s destination tail: destination cartage, CFS dwell, customs clearance, final delivery. Used in C.10a to compute $\text{arr\_dest}(k, a) = \text{ETA}_a + \Delta^{\text{post}}_k$ for air terminal arcs. & h \\
Per-HAWB customs dwell & $\delta^{\text{cust}}_k$ & Dwell time on the customs-clearance arc; depends on destination country, cargo type, commodity class, broker performance, pre-filed entry status & h \\
Per-HAWB AWB release dwell & $\delta^{\text{rel}}_k$ & Dwell at destination CFS between cargo arrival and AWB release. Depends on release type: \texttt{express\_release} $\approx 0$h; \texttt{telex} 1--4h; \texttt{original\_awb\_surrender} 24--72h. Folded into $\delta_a$ of the destination CFS dwell arc at graph build. & h \\
Cargo class & $\ctype_k$ & Canonical enum $\{\text{GEN, DGR, PER, VAL, AVI, HUM}\}$ following IATA Cargo-IMP & categorical \\
Temperature regime & $\text{temperature}(k)$ & $\{\text{ambient, chilled, frozen, pharma}\}$ (subdivides PER; other classes default to ambient) & categorical \\
HS code & $\text{hs}_k$ & 6-digit Harmonized System code & --- \\
Incoterm & $\text{inco}_k$ & EXW / FCA / CIF / DDP / \ldots & --- \\
Service-product binding & $\text{sp}(k) \in P$ & The tenant service product the HAWB is contracted under & --- \\
Lithium spec (optional) & $\text{lithium\_spec}_k$ & UN number, PI code, section, Wh, SOC compliance, DDR flag (\S\ref{sec:lithium}) & --- \\
\bottomrule
\end{tabular}
\end{table}

\paragraph{Cargo-ready window.}
The shipper offers cargo for pickup within $[t_k^{\text{rdy,early}},
t_k^{\text{rdy,late}}]$, not at a single moment. The arrival time at the
origin door is bounded by this window (\S\ref{sec:domain}). Anchors throughout
the rest of the model (subgraph optimistic arrival, big-M tightening, SLA
computation) reference $t_k^{\text{rdy,early}}$ as the canonical ``cargo first
available'' time --- service-level commitments are measured from the shipper's
earliest readiness, not from when the forwarder elects to pick up.

\paragraph{Chargeable weight (per-HAWB).}
$cw_k = \max(w_k, v_k \cdot 167)$ is the IATA per-HAWB billing reference.
\textbf{It does \emph{not} appear in any bucket-capacity constraint} (item 13-A
fix --- see C.5b). Light/bulky cargo's volumetric weight is captured by the
volume capacity C.5b-v, not by double-counting $cw_k$ against the weight
capacity. $cw_k$ is used only in the co-load per-kg cost term (\S\ref{sec:objective})
and as a sanity reference --- billing aggregations on MAWB-arcs use the
density-mixed $\text{CW}_{a,g}$ from \S\ref{sec:cw-density-mixing}.

\paragraph{Tenant-global parameters supporting C.10 and the fallback arc.}
Two scalars are tenant-configured and apply to every HAWB in a solve:
\begin{itemize}[noitemsep]
  \item $V^{\text{ref}}$ (USD) --- shipment-value reference scale used to
        normalize $\mu_k$. Setting $\mu_k = \text{value}_k / V^{\text{ref}}$
        means a HAWB at the reference value carries $\mu_k = 1$; a 10$\times$
        more valuable HAWB carries $\mu_k = 10$. Recommended default
        $V^{\text{ref}} = \$10{,}000$ (\texttt{CALIBRATION NEEDED}).
  \item $C^{\text{fallback}}$ (USD) --- dollar cost of routing one HAWB via
        the fallback arc (\S\ref{sec:fallback-arc}). Must be large enough
        that the MILP prefers any feasible real route over the fallback.
        \emph{Set to the smallest constant that achieves dominance}, not a
        round large number --- an unnecessarily large $C^{\text{fallback}}$
        weakens the LP relaxation and enables fractional fallback smearing
        at the LP root. Recommended sizing rule:
        \[
          C^{\text{fallback}} \;\geq\; \max_{k \in K}\, W_k \cdot
          (T_k^{\text{abs}} - \Delta_k)^2 \;+\; 10 \cdot \max_{k \in K,\, a \in A_k}
          \text{real\_cost}(k, a)
        \]
        where $\text{real\_cost}(k, a)$ is the worst-case per-HAWB
        realistic routing cost (sum of arc cost contributions for an
        all-real-arcs path). For a tenant with worst real-route cost
        $\sim \$10$K per HAWB and worst-case $W_k (T_k^{\text{abs}} -
        \Delta_k)^2 \sim \$50$K, $C^{\text{fallback}} \sim \$150$K is
        sufficient; the prior $\$1$M placeholder was unnecessarily
        large. The sizing must be \emph{per-tenant} (depends on the
        tenant's $V^{\text{ref}}$, service-product weights $w_p$, and
        backstop horizons), so $C^{\text{fallback}}$ is a tenant-config
        value computed at deployment from the catalog, not a global
        constant. \texttt{CALIBRATION NEEDED} for the actual tenant
        values. The walking-skeleton metric
        \texttt{fractional\_fallback\_usage} (\S\ref{sec:walking-skeleton-instrumentation})
        surfaces LP smearing if $C^{\text{fallback}}$ is set too low.
\end{itemize}

\subsection{Ground-arc parameters (pickup, cartage, CFS dwell, customs, delivery)}
\label{sec:ground-arc-params}

For each ground arc $a \in A^{\text{ground}}$:

\begin{table}[H]
\centering
\caption{Per-arc parameters for ground arcs (pickup, cartage, CFS dwell, customs, delivery).}
\label{tab:ground-arc-params}
\begin{tabular}{p{3.5cm}p{2.2cm}p{6.5cm}p{1.5cm}}
\toprule
Parameter & Symbol & Description & Unit \\
\midrule
Dwell / transit time & $\delta_a$ & Per-arc dwell or transit (pickup truck, cartage, CFS build-up, deconsol-dwell, AWB release dwell, final delivery). Includes any in-transit hub customs adjustment for deconsol arcs. & h \\
Flat per-HAWB handling charge & $c_a^{\text{flat}}$ & e.g.\ per-AWB CFS handling fee, deconsol fee. May be 0. & USD/HAWB \\
Per-kg handling charge & $c_a^{\text{kg}}$ & e.g.\ per-kg cartage. May be 0. & USD/kg \\
\bottomrule
\end{tabular}
\end{table}

Total per-HAWB ground-arc cost on arc $a$: $(c_a^{\text{flat}} + c_a^{\text{kg}}
\cdot w_k) \cdot x_{k,a}$. The customs-clearance dwell arc (\(a \in
A^{\text{cust}}\)) additionally carries $\delta^{\text{cust}}_k$ per-HAWB; see
\(\text{transit}(k, a)\) in \S\ref{sec:sets}. AWB release dwell
$\delta^{\text{rel}}_k$ is folded into the destination CFS dwell arc's
$\delta_a$ at graph build (per-HAWB lookup from $\text{release\_type}$).

\paragraph{Origin / destination trucking is a per-HAWB approximation (scope boundary).}
This air model does \emph{not} route or consolidate first/last-mile trucking.
Origin pre-carriage (pickup, origin cartage) and destination on-carriage
(destination cartage, final delivery) enter only as a per-HAWB \emph{cost
estimate}: $c_a^{\text{flat}}$ (per-AWB) plus $c_a^{\text{kg}} \cdot w_k$, where
the ground arc encodes the specific origin/destination leg (hence the
lane/distance) and the coefficients are estimated from the HAWB's size/weight
class and that lane. Because each HAWB carries its own estimate, the per-HAWB
attribution $\times\, x_{k,a}$ is correct by construction --- there is no
shared-truck pooling in this model, so no flat charge is multi-counted across
HAWBs sharing a move. Detailed truck routing and multi-HAWB consolidation
(where shared-vehicle flat costs and consolidation savings are actually
modeled) are deferred to a downstream \emph{trucking planner} that runs
\emph{after} the air legs are firmed up --- origin-side first, destination-side
later --- closer to execution time when the realized HAWB set per lane is
known. The estimate here is deliberately good-enough for choosing air legs; it
is not the trucking plan.

\paragraph{Sourcing $\delta^{\text{cust}}_k$ and $\delta^{\text{rel}}_k$.}
For MVP, both resolve from tenant-configurable lookup tables keyed on
shipment attributes:
\begin{itemize}[noitemsep]
  \item $\delta^{\text{cust}}_k = \text{lookup}^{\text{cust}}(\text{destination\_country},
        \ctype_k,\, \text{commodity\_class})$.
  \item $\delta^{\text{rel}}_k = \text{lookup}^{\text{rel}}(\text{release\_type})$.
\end{itemize}
Initial values are tenant-calibrated from historical clearance times observed
in the forwarder's own broker performance data. The customs filing
\emph{deadlines} (AMS / ICS2 / ACI) are absorbed into the air-arc cutoff
scalar $\text{CO}_a^*$ (\S\ref{sec:air-arc-params}); $\delta^{\text{cust}}_k$
captures the \emph{post-arrival} clearance hold, distinct from the pre-flight
filing.

\paragraph{In-transit hub customs.}
The time impact of in-transit declarations (US T\&E, EU summary declaration,
etc.) is an additive component of the existing deconsolidation-dwell arc's
$\delta_a$, computed by the graph generator from hub type + cargo type +
in-transit flag. No new arc type, no new MILP constraints.

\subsection{Per-arc air parameters}
\label{sec:air-arc-params}

For each air arc $a \in A^{\text{MAWB}} \cup A^{\text{coload}}$:

\begin{table}[H]
\centering
\caption{Per-arc air parameters.}
\label{tab:air-arc-params}
\begin{tabular}{p{3cm}p{2.2cm}p{7cm}p{1.5cm}}
\toprule
Parameter & Symbol & Description & Unit \\
\midrule
Realized air transit time & $\mu_a$ & Graph-build-precomputed scalar. For multi-leg arcs (same-carrier connection or interline through-MAWB), $\mu_a$ already includes \emph{internal MCT} between legs --- no separate MILP constraint family for hub MCT. & h \\
Effective cargo cutoff at $\text{tail}(a)$ & $\text{CO}_a^*$ & Air arcs leaving a POL: $\max(\text{DCO}, \text{AMS}, \text{ICS2}, \text{ACI})$ across all legs of $a$, minus per-HAWB prep time. Graph-build-precomputed scalar. & h (UTC) \\
Rate family & $\text{rate\_family}_a$ & $\{\texttt{flat\_rate}, \texttt{min\_flat\_breaks}, \texttt{per\_uld\_pivot}, \texttt{coload\_per\_kg}\}$. \texttt{coload\_per\_kg} applies only to $A^{\text{coload}}$. & enum \\
Quotation currency & $\text{currency}_a$ & ISO 4217; converted to USD at MILP build. & --- \\
\bottomrule
\end{tabular}
\end{table}

Family-specific:

\begin{table}[H]
\centering
\caption{Rate-family-specific parameters per air arc.}
\label{tab:rate-family-params}
\begin{tabular}{lp{4.6cm}p{6.5cm}}
\toprule
Family & Symbols & Cost \\
\midrule
\texttt{flat\_rate} & $m_a$ (USD/kg), $\text{min\_chg}_a$ (USD per instantiated MAWB), $\text{cap}_a$ (kg, optional) & $\max(\text{min\_chg}_a \cdot z_{a,g},\, m_a \cdot \text{CW}_{a,g})$. Covers NAC, block rates, and dynamic spot quotes sold at one flat rate with no weight-break discount. \\[0.3em]
\texttt{min\_flat\_breaks} & $\{(\text{break}_{a,b},\, \text{rate}_{a,b})\}_{b \in B_a}$, $\text{cap}_a$ (kg, optional) & IATA next-break-down: $\min_b \text{rate}_{a,b} \cdot \max(\text{CW}_{a,g},\, \text{break}_{a,b})$. Covers TACT GCR, SCR, CCR. $B_a$ ordered ascending; $\text{break}_{a,b}$ is the minimum billed weight if break $b$ is selected; $\text{rate}_{a,b}$ is the per-kg rate. \\[0.3em]
\texttt{per\_uld\_pivot} (BSA) & $r_a$ (USD/kg) = $r_c$ for the tagging contract; $\pi_a$ (kg) pivot weight; $U_a \subseteq U$; $N_{a,u}$ (count) contracted ULD-position allotment; $W_u, V_u$ (kg, m\textsuperscript{3}) ULD physical capacity; $\text{settlement}_a \in \{\texttt{per\_flight}, \texttt{equalized}\}$ & $r_a \cdot \max(\text{CW}_{a,g},\, \pi_a \cdot \sum_u \eta_{a,g,u})$. BSA take-or-pay; pivot floor binds per flight for \texttt{per\_flight} settlement, or via the sunk allowance mechanism (C.13a) for \texttt{equalized}. \\[0.3em]
\texttt{coload\_per\_kg} & $m_a^{\text{cl}}$ (USD/kg), $\text{cap}_a^{\text{cl}}$ (kg, optional) & Per-kg, billed on each HAWB's own $cw_k$. \\
\bottomrule
\end{tabular}
\end{table}

$\text{cap}_a / \text{cap}_a^{\text{cl}}$ are bounds on \emph{actual} weight
summed across HAWBs on the arc (across consolidation groups), enforced by
C.5c. Offers without a specified cap (typical TACT/NAC, typical informal
co-load) are treated as uncapped at planning time; capacity check happens at
the carrier's request/confirm step at booking.

\paragraph{Per-contract parameters (per-ULD-pivot only).}
\begin{table}[H]
\centering
\caption{Per-contract parameters (per-ULD-pivot rate family only).}
\label{tab:per-contract-params}
\begin{tabular}{p{2.8cm}p{2.5cm}p{7cm}p{1.5cm}}
\toprule
Parameter & Symbol & Description & Unit \\
\midrule
Per-kg contract rate & $r_c$ & Equal to $r_a$ for every $a \in A_c^{\text{MAWB}}$ (consistency assumption: all arcs tagged to one contract carry the same rate) & USD/kg \\
Remaining sunk allowance & $A_c$ & For $c \in C^{\text{eq}}$: $A_c = P_{c,t} - (\text{chargeable weight already consumed on $c$ this period})$, supplied exogenously per solve by the consumed-weight accumulator (\S\ref{sec:bsa-equalized-settlement}). Undefined for \texttt{per\_flight} settlement. & kg \\
\bottomrule
\end{tabular}
\end{table}

\paragraph{Per-service-product parameters.}
\begin{table}[H]
\centering
\caption{Per-service-product parameters.}
\label{tab:service-product-params}
\begin{tabular}{p{2.8cm}p{2.5cm}p{7cm}p{1.5cm}}
\toprule
Parameter & Symbol & Description & Unit \\
\midrule
SLA transit-time bound & $T^{\text{SLA}}_p$ & For service product $p \in P$ & h \\
Tardiness weight & $w_p$ & Quadratic soft-tardiness penalty coefficient for service product $p$; combined with the per-HAWB value coefficient $\mu_k$ into $W_k = w_p \cdot \mu_k$ which scales $\tau_k^2$ in the objective (C.10, \S\ref{sec:lin-tardiness}). \texttt{CALIBRATION NEEDED} --- defaults pending. & USD / h\textsuperscript{2} \\
\bottomrule
\end{tabular}
\end{table}

\paragraph{Other.}
\begin{table}[H]
\centering
\caption{Other per-MAWB parameters (dunnage factor; MAWB fixed-charge).}
\label{tab:other-params}
\begin{tabular}{p{2.8cm}p{2.5cm}p{7cm}p{1.5cm}}
\toprule
Parameter & Symbol & Description & Unit \\
\midrule
Dunnage factor & $\varepsilon$ & Dimensionless; default $\approx 0.05$ ($\approx 5\%$); used in $\text{Wt}_{a,g}$ aggregation (Eq.~\ref{eq:mawb-wt}) & --- \\
MAWB fixed-charge & $c^{\text{MAWB}}_{\text{fix}}$ & Per-active-MAWB documentation overhead: customs manifest filing (AMS / ENS), MAWB issuance handling, and forwarder admin. \textbf{Scoped to ``new AWB number issued''} --- one charge per distinct MAWB-arc that triggers a fresh AWB filing, not per arc-activation when a single AWB spans multiple arcs (e.g.\ a future arc-splitting model edit must guard against double-charging the through-AWB case). Placeholder $\$50$/MAWB; \texttt{MARKET RESEARCH NEEDED}. Wired to $z_{a,g}$ in the objective (\S\ref{sec:objective}). & USD/MAWB \\
\bottomrule
\end{tabular}
\end{table}

\paragraph{Per-flight metadata (graph-build only).}
Physical flights $f$ exist in \texttt{air\_graph\_construction.md} but are
\emph{not} indexed by the MILP. The following per-flight attributes are
consumed by the graph generator to compute per-arc scalars, run the
subgraph pre-filter, and instantiate predicate values:
\begin{table}[H]
\centering
\caption{Per-flight metadata consumed by the graph generator (not indexed by the MILP).}
\label{tab:flight-metadata}
\begin{tabular}{p{2.8cm}p{2.5cm}p{7cm}p{1.5cm}}
\toprule
Attribute & Symbol & Description & Unit \\
\midrule
Scheduled departure time & $\text{ETD}_f$ & UTC; used in embargo activeness window (Eq.~\ref{eq:active-embargoes}) and rate-validity windows. & h (UTC) \\
Scheduled arrival time & $\text{ETA}_f$ & UTC; aggregated into the per-arc $\text{ETA}_a$ scalar (\S\ref{sec:graph-construction}) and used directly in stochastic-transit P1 hooks. & h (UTC) \\
Aircraft type & $\text{ac}(f)$ & Categorical: specific aircraft model (e.g.\ B777F, A330-200F, B747-400F, A330-300-pax-belly, etc.). Aggregates into the broader $\text{ac\_type}(f)$ enum. & --- \\
Aircraft-type bucket & $\text{ac\_type}(f)$ & $\{\texttt{FREIGHTER},\, \texttt{PAX\_BELLY}\}$; used in embargo records (Eq.~\ref{eq:active-embargoes}), lithium acceptance (Eq.~\ref{eq:lithium-accept}), and service-product $\text{ac\_type\_ok}$ predicate (Eq.~\ref{eq:sp-ac-ok}). & enum \\
DGR capability & $\text{dgr}(f) \in \{0, 1\}$ & 1 if $f$'s operator + aircraft accept dangerous goods on this leg; used in $\text{cargo\_type\_ok}$ (Eq.~\ref{eq:cargo-type-ok}). & boolean \\
PER capability & $\text{per}(f) \in \{0, 1\}$ & 1 if $f$ accepts perishable cargo (cold-chain handling, etc.); used in $\text{cargo\_type\_ok}$. & boolean \\
VAL capability & $\text{val\_capable}(f) \in \{0, 1\}$ & 1 if $f$ accepts valuable cargo (secure handling). & boolean \\
AVI capability & $\text{avi\_capable}(f) \in \{0, 1\}$ & 1 if $f$ accepts live animal cargo. & boolean \\
HUM capability & $\text{hum\_capable}(f) \in \{0, 1\}$ & 1 if $f$ accepts human remains cargo. & boolean \\
Operating carrier & $\text{carrier}^{\text{op}}(f)$ & IATA carrier code of the operating airline (physical aircraft operator); used in carrier-policy cascade (Eq.~\ref{eq:carrier-ok-resolved}). & IATA \\
Marketing carrier & $\text{carrier}^{\text{mk}}(f)$ & IATA carrier code of the marketing airline (AWB-prefix carrier; codeshare-aware); used in carrier-policy cascade. & IATA \\
Cutoff stack & $\text{DCO}_f,\, \text{AMS}_f,\, \text{ICS2}_f,\, \text{ACI}_f$ & Per-leg carrier cargo cutoff, plus customs filing deadlines; aggregated into $\text{CO}_a^*$ at graph build (\S\ref{sec:graph-construction}). & h (UTC) \\
Lithium acceptance matrix & $\text{lithium\_accept}_f$ & Function on (PI, Section, ac\_type); see Eq.~\ref{eq:lithium-accept} and \S\ref{sec:lithium}. & --- \\
\bottomrule
\end{tabular}
\end{table}
This metadata is consumed by the graph generator to:
\begin{itemize}[noitemsep]
  \item compute the per-arc scalars $\mu_a$ and $\text{CO}_a^*$;
  \item determine the arc \emph{type} (MAWB-eligible vs co-load) and admissible
        ULD types $U_a$;
  \item resolve the cutoff stack against the destination region;
  \item apply embargo, lithium, cargo-type, service-product, and
        carrier-policy predicates during the per-shipment subgraph pre-filter
        (\S\ref{sec:prefilter}).
\end{itemize}

\subsection{ULD specifications}
\label{sec:uld-spec}

ULD specifications are governed by IATA ULD Regulations
(ULDR).\footnote{IATA Unit Load Device Regulations (ULDR).
\url{https://www.iata.org/en/publications/manuals/uld-regulations/}} The
model needs a richer per-ULD attribute set than volume/weight alone --- physical
fit at the aircraft door, floor-load tolerance, deck position, tare-vs-cargo
weight separation, and certification all influence whether a (HAWB, MAWB-arc,
ULD type) triple is operationally feasible.

\paragraph{IATA ULD identification format.}
The standard identifier is \texttt{XXX-NNNNN-CC}: 3 leading letters (type
code), 4--5 numerals (serial), 2 trailing letters (owning airline IATA
code).\footnote{ULD CARE, ``ULD Identification Codes Demystified.''
\url{https://www.uldcare.com/articles/library/care/uld-identification-codes-demystified/}}
The 3-letter type code has structural meaning:
\begin{itemize}[noitemsep]
  \item \textbf{1st letter}: family. \texttt{P} = Pallet, \texttt{A} = Aircraft
        container (rigid walls), \texttt{R} = Reefer, \texttt{U} = Non-certified.
  \item \textbf{2nd letter}: base dimensions. \texttt{M} =
        96\textquotedbl{}$\times$125\textquotedbl{} (main-deck size), \texttt{K} =
        60.4\textquotedbl{}$\times$61.5\textquotedbl{}, \texttt{L} =
        60.4\textquotedbl{}$\times$125\textquotedbl{}.
  \item \textbf{3rd letter}: contour. \texttt{C} = standard cube, \texttt{E} =
        contoured, \texttt{P} = lower-deck contour.
\end{itemize}
Example: \textbf{PMC} = Pallet, 96\textquotedbl{}$\times$125\textquotedbl{} base,
cube contour --- a flat aluminum pallet $244 \times 318$ cm with cargo stacked
on top and secured by net (main-deck freighter ULD). Contrast \textbf{LD3
(AKE/AVE)}: rigid-walled container used in wide-body belly compartments, cargo
stuffed inside through a door.

\paragraph{Reference catalog (illustrative; MVP synthetic instances).}
\begin{table}[H]
\centering
\small
\caption{Reference ULD catalog (illustrative; for MVP synthetic instances).}
\label{tab:uld-reference}
\begin{tabular}{lllrrrrl}
\toprule
Family & Type code & Base (cm) & MGW (kg) & Tare (kg) & Cargo (kg) & Vol (m\textsuperscript{3}) & Deck \\
\midrule
LD3 & AKE/AVE & 153 $\times$ 153 & 1{,}588 & $\approx$80 & 1{,}508 & 4.5 & lower \\
LD7 & AKH/AAY/ALE & 244 $\times$ 318 & 3{,}175 & $\approx$120 & 3{,}055 & 11.1 & lower \\
LD11 & AMP & 244 $\times$ 318 & 6{,}804 & $\approx$240 & 6{,}564 & 7.0 & lower / main \\
PMC & PMC/PAG/P6P & 244 $\times$ 318 & 6{,}804 & $\approx$130 (pallet+net) & 6{,}674 & 7.5 (cube contour) & main \\
\bottomrule
\end{tabular}
\end{table}
Tare values are nominal; airlines publish exact tare on each ULD's
identification plate. The model uses the nominal tare from the catalog; exact
tare matters at billing, not at planning.

\paragraph{Operational parameter set per $u \in U$.}
\begin{table}[H]
\centering
\caption{Operational parameter set per ULD type $u \in U$.}
\label{tab:uld-params}
\begin{tabular}{p{3cm}p{2.2cm}p{7cm}p{1.5cm}}
\toprule
Parameter & Symbol & Description & Unit \\
\midrule
Family name & --- & LD3 / LD7 / LD11 / PMC / AKE / \ldots & --- \\
Type code & --- & IATA 3-letter type & --- \\
Base dimensions & --- & Width $\times$ length of footprint & cm \\
Internal volume & $V_u$ & Usable interior after contour & m\textsuperscript{3} \\
External dimensions & --- & $L \times W \times H$ including frame (door clearance) & cm \\
Max gross weight & $\text{MGW}_u$ & Total ULD weight loaded (cargo + tare) & kg \\
Tare weight & $\text{tare}_u$ & Empty ULD weight & kg \\
Cargo weight capacity & $W_u$ & $\text{MGW}_u - \text{tare}_u$; effective payload --- used in C.5b-w & kg \\
Floor load limit & $\phi_u$ & Max kg per m\textsuperscript{2} on the floor & kg/m\textsuperscript{2} \\
Contour profile & --- & Aircraft compatibility code & --- \\
Deck position class & --- & \texttt{main} / \texttt{lower} / \texttt{either} & --- \\
Pressurized hold & --- & Live animals, sensitive perishables & boolean \\
Fill rate cap & $\eta_u$ & Practical max fraction of $V_u$ usable due to nets, cargo shape (default $0.97$) & --- \\
Build-up time & $\beta_u^{\text{up}}$ & Forwarder labor to build a $u$-type ULD & min \\
Breakdown time & $\beta_u^{\text{down}}$ & Time to break down at destination CFS & min \\
ULDR certification & --- & IATA ULDR certification level & --- \\
\bottomrule
\end{tabular}
\end{table}

\paragraph{Floor-load constraint.}
Each ULD has a floor-load limit $\phi_u$ (typically $\sim 730$
kg/m\textsuperscript{2} for standard ULDs).\footnote{Hansatic Air Freight ULD
Guide. \url{https://hansatic.com/en/guides/air-freight-uld-guide}. The binding
value is the minimum of the ULD limit and the aircraft floor structure
limit.} A HAWB with a heavy piece on a small footprint (e.g., a 1{,}500 kg
machine on a 0.8 m\textsuperscript{2} base $\to \sim 1{,}875$
kg/m\textsuperscript{2}) is operationally infeasible on that ULD even if
aggregate weight fits. Captured in the subgraph pre-filter step 8
(\S\ref{sec:prefilter}) alongside the contour-fit check.

\paragraph{Aircraft compatibility.}
$U_a$ (admissible ULD types on \texttt{per\_uld\_pivot} arc $a$) is derived
from the underlying aircraft type and the published aircraft--ULD
compatibility matrix (e.g., 777F accepts PMC main-deck + AKE/AKH/AAY
lower-deck; A330 belly accepts AKE/AKH only). For multi-leg MAWB-arcs the
intersection of per-leg admissible sets is used. The graph generator emits
$U_a$ as part of the arc data.

\subsection{Procurement types and the supply-option catalog (per arc)}
\label{sec:supply-option-catalog}

The forwarder procures air capacity through several structurally distinct
channels. In v3 every offer is attached to an \emph{arc} (a MAWB-arc or a
co-load arc), and each arc carries one offer. The arc's
$\text{rate\_family}_a$ attribute selects how the cost is computed.

\paragraph{Supply types.}
Five operational procurement channels feed the arc catalog (the
\texttt{COLOADER} channel is realized through co-load arcs at the graph
layer; from the MILP's point of view it is just a \texttt{coload\_per\_kg}
arc):
\begin{table}[H]
\centering
\caption{Supply types feeding the per-arc offer catalog.}
\label{tab:supply-types}
\begin{tabular}{p{2.5cm}p{2.0cm}p{2.0cm}p{2.5cm}p{4.5cm}}
\toprule
Channel & Allocation held by & Tenant contract with & Arc type & Typical $\text{rate\_family}_a$ \\
\midrule
\texttt{DIRECT\_BSA\_HARD} & Tenant & Carrier & MAWB-arc & \texttt{per\_uld\_pivot} \\
\texttt{DIRECT\_BSA\_SOFT} & Tenant & Carrier & MAWB-arc & \texttt{per\_uld\_pivot} or \texttt{min\_flat\_breaks} \\
\texttt{GSA} & GSA & GSA & MAWB-arc & \texttt{min\_flat\_breaks} or \texttt{flat\_rate} \\
\texttt{COLOADER} & Co-loader & Co-loader & Co-load arc & \texttt{coload\_per\_kg} \\
\texttt{DYNAMIC\_SPOT} & None & Airline / platform & MAWB-arc & \texttt{flat\_rate} or \texttt{min\_flat\_breaks} \\
\bottomrule
\end{tabular}
\end{table}

\paragraph{Rate-family mapping --- MVP-active.}
\begin{itemize}[noitemsep]
  \item \texttt{flat\_rate} and \texttt{min\_flat\_breaks} rate the MAWB
        aggregate $\text{CW}_{a,g}$ --- fully MVP-active. Family 2 is the IATA
        next-break-down rule (Eq.~\ref{eq:rate-fn-minflat}); volume-discount
        cost shapes live here, not in any ``cumulative PWL'' family. Air cargo
        is not rated by accumulating per-segment marginals.
  \item \texttt{per\_uld\_pivot} rates per ULD position via $\eta_{a,g,u}$ and
        the per-flight pivot floor (P.10 / C.13b). MVP-active for the
        single-pivot shape with $\text{settlement}_a \in \{\texttt{per\_flight},
        \texttt{equalized}\}$. A genuine multi-tier per-ULD over-pivot rate
        structure (several over-pivot brackets) was \emph{not} found in the
        tariff sources reviewed and is recorded as a contingency only
        (\S\ref{sec:deferred}).
  \item \texttt{coload\_per\_kg} costs each HAWB at $m_a^{\text{cl}} \cdot
        cw_k$ on a co-load arc; no MAWB object is instantiated. Co-loader
        consolidation is opaque to the model.
\end{itemize}

\paragraph{Per-offer cap $\text{cap}_a / \text{cap}_a^{\text{cl}}$.}
Where the catalog offer specifies a cap on actual weight summed across HAWBs
on the arc (BSA secondary cap, GSA tier limit, co-loader quote ceiling), C.5c
enforces it. Offers without a specified cap are treated as uncapped at
planning time --- capacity check happens at the carrier's request/confirm step.

\paragraph{Dynamic spot --- rate validity.}
\texttt{DYNAMIC\_SPOT}-tagged offers carry $\text{expiry}_a$: the latest ETD
of any leg of $a$ at which the quote remains valid. Validated at the data
layer (graph generator); the MILP does not need to express it.

\paragraph{Cardinality at MVP scale.}
For a typical mid-market lane (TPE--JFK GEN): per direct or hub-routed O-D
segment, $\sim 5$ candidate MAWB-arcs ($\sim 1$ BSA + $\sim 1$ TACT + $\sim 1$
GSA + $\sim 1$ dynamic spot + $\sim 1$ NAC \texttt{flat\_rate}) plus typically
$\sim 1$ co-load arc. The break-selection binaries $\gamma_{a,g,b}$ on
\texttt{min\_flat\_breaks} arcs index off the $|K|$ axis: the IATA tiers of a
TACT card produce one binary per $(\text{arc}, g, b)$, not per HAWB. This is
the key scale-friendly property of the $(a, g)$ formulation.

\subsection{Per-contract BSA parameters and the equalized-allowance mechanism}
\label{sec:bsa-params}

A Block Space Agreement (BSA) reserves a fixed allocation of ULD positions on
a contracted set of MAWB-arcs at a contracted rate. For each contract $c \in
C$:

\begin{table}[H]
\centering
\caption{Per-contract BSA parameters (for $c \in C$).}
\label{tab:bsa-params}
\begin{tabular}{p{3.5cm}p{2.2cm}p{6.5cm}p{1.5cm}}
\toprule
Parameter & Symbol & Description & Unit \\
\midrule
Tagged arcs & $A_c^{\text{MAWB}} \subseteq A^{\text{MAWB}}$ & MAWB-arcs whose offer is contract $c$ & --- \\
Per-arc allotment & $N_{a,u}$ & ULD-position count of type $u$ on arc $a \in A_c^{\text{MAWB}}$ & ULDs \\
Pivot weight & $\pi_a$ & Minimum chargeable weight per ULD per flight under $c$ on arc $a$ & kg/ULD \\
Contract per-kg rate & $r_c$ & Equal to $r_a$ for every $a \in A_c^{\text{MAWB}}$ & USD/kg \\
Settlement basis & $\text{settlement}_a$ & $\{\texttt{per\_flight}, \texttt{equalized}\}$ --- per-arc attribute (a contract may have arcs of either basis, but typically uniform) & --- \\
Remaining sunk allowance (equalized only) & $A_c$ & Exogenous per solve from the consumed-weight accumulator (\S\ref{sec:bsa-equalized-settlement}); undefined for \texttt{per\_flight} & kg \\
Period commitment (equalized only) & $P_{c,t}$ & Period pivot weight $\times$ committed ULD count for period $t$; the take-or-pay base & kg \\
\bottomrule
\end{tabular}
\end{table}

\paragraph{Per-flight settlement.}
$\text{settlement}_a = \texttt{per\_flight}$: the pivot binds per offer
--- one pivot per MAWB-arc, by construction --- against the chargeable
weight booked on that MAWB. The cost is $r_a \cdot \max(\text{CW}_{a,g},\,
\pi_a \cdot \sum_u \eta_{a,g,u})$, linearized in C.13b. \textbf{For
multi-leg MAWB-arcs} (interline through-AWB; same-carrier through-cargo
sold at one rate), the offer is one contractual pivot floor over the full
O-D segment, not one pivot per leg. Per-leg pivot flexibility comes from
the graph generator emitting separable single-leg MAWB-arcs as parallel
candidates alongside the through arc (\S\ref{sec:arc-enumeration}), each
of which then carries its own pivot.

\paragraph{Equalized settlement --- the IATA equalization clause.}
\label{sec:bsa-equalized-settlement}
$\text{settlement}_a = \texttt{equalized}$: the pivot is averaged over an
allocation period via an equalization clause. The take-or-pay is on the
\emph{period total chargeable weight} tendered on contract $c$'s arcs,
against the period commitment $P_{c,t}$ (the period pivot weight $\times$ the
committed ULD count).

Because the MILP solves \emph{per batch} and cannot see the whole period, the
equalized clause is \emph{not} an in-MILP period constraint. Instead the
contract is rendered as a two-segment cost on the chargeable weight tendered
on $c$ \emph{this solve}:
\begin{itemize}[noitemsep]
  \item Weight $0 \to A_c$ is free (sunk: already paid via the period
        commitment).
  \item Weight above $A_c$ costs $r_c$/kg (the marginal overage rate).
\end{itemize}
$A_c$ is the \emph{remaining sunk allowance} this solve: $A_c = P_{c,t} -
(\text{chargeable weight already consumed on $c$ this period})$, maintained
by an external consumed-weight accumulator that runs alongside the MILP. The
accumulator reduces to a running total per contract; no price, no
subgradient, no tuning, no forecast in the cost loop. C.13a applies. The
constant sunk cost $r_c \cdot P_{c,t}$ is booked by the accumulator outside
the MILP --- it is a sunk constant that does not affect the optimizer's
\emph{argmin} (only the objective value's level).

\paragraph{No hard period-minimum.}
BSA period take-or-pay is not a hard count minimum on the per-batch MILP ---
\texttt{per\_flight} settlement prices it per flight (C.13b), and
\texttt{equalized} settlement absorbs it via $A_c$ (C.13a). A hard
period-minimum would risk spurious infeasibility when the current batch has
nothing eligible to load on contract $c$.

\subsection{Surcharge parameters}
\label{sec:surcharge}

Real air-freight quotes layer 10--20 surcharges per shipment per arc on top of
the base freight rate. Fuel (FSC) and security (SSC) are the most visible;
the full stack includes terminal handling, customs filing (AMS / ICS2 / ACI),
AWB issuance, screening, DGR handling, war risk, peak-season surcharge,
perishable / cool-chain handling, valuables handling, overweight surcharge,
and lithium handling. Together these typically add 30--50\% on top of the
base rate for a tier-2 forwarder quote; omitting them produces systematic
under-quoting.

\paragraph{Storage.}
Surcharges live in the \texttt{surcharge\_catalog} table per tenant
(\texttt{data\_model.md \S6}), versioned via append-only $+$ per-run-snapshot
just like policy rules (\S4) and spot rate snapshots (\S5). Each catalog row
carries: \texttt{surcharge\_type} (enumerated), \texttt{scope} (applicability
JSONB: carrier, origin, destination, ac\_type, cargo\_type, region, date
range), and \texttt{calculation} (basis: \texttt{per\_kg} / \texttt{per\_shipment}
/ \texttt{per\_awb} / \texttt{per\_uld} / \texttt{per\_cbm} / \texttt{per\_piece}
/ \texttt{percent\_of\_freight}, rate, currency, optional min/max).

\paragraph{Per-arc resolution.}
Pre-solve, for each HAWB $k$ and each arc $a \in A_k$, the resolver
materializes the applicable surcharges:
\begin{equation}
  S_{\text{arc}}(k, a) \;\defeq\; \{\, s \in \text{surcharge\_catalog} :
    \text{scope\_match}(s, k, a) \wedge s.\text{status} = \text{active} \wedge
    \text{now} \in [s.\text{effective\_from}, s.\text{effective\_to})\,\}
  \label{eq:surcharge-applicable}
\end{equation}
The per-arc surcharge cost is the sum of resolved amounts:
\begin{equation}
  \text{surcharge\_cost}(k, a) \;\defeq\; \sum_{s \in S_{\text{arc}}(k, a)}
    \text{amount}(s, k, a)
  \label{eq:surcharge-cost}
\end{equation}
where $\text{amount}(s, k, a)$ is computed per the surcharge's calculation
basis and clipped by the optional $[\text{min\_charge}, \text{max\_charge}]$
band.

\paragraph{Path A --- per-HAWB-per-arc bases.}
Resolve into $\text{surcharge\_cost}(k, a)$:
\begin{itemize}[noitemsep]
  \item \texttt{per\_kg}: $\text{rate} \cdot cw_k$ (or $w_k$ if
        \texttt{weight\_basis} = \texttt{actual}).
  \item \texttt{per\_shipment}: $\text{rate}$ (flat).
  \item \texttt{per\_awb}: $\text{rate}$ (flat per AWB; for a single-HAWB
        shipment equivalent to \texttt{per\_shipment}).
  \item \texttt{per\_cbm}: $\text{rate} \cdot v_k$.
  \item \texttt{per\_piece}: $\text{rate} \cdot \text{piece\_count}_k$.
  \item \texttt{percent\_of\_freight}: $\text{rate} \cdot
        \text{base\_freight}_k$ (with \texttt{percent\_of\_field} =
        \texttt{base\_freight}, pre-computed per arc).
\end{itemize}
$cw_k$, $v_k$, $w_k$, $\text{piece\_count}_k$, $\text{base\_freight}_k$ are
all known constants; $\text{surcharge\_cost}(k, a)$ for Path A is a constant
once resolved.

\paragraph{Path B --- per-ULD build/breakdown bases.}
The \texttt{per\_uld} basis applies once per loaded ULD per
build/breakdown \emph{event}, regardless of which HAWBs share those ULDs.
A ULD is built once at its origin and broken down once at its destination.
The naive per-flight accounting view,
\begin{equation}
  \text{flight\_uld\_surcharge\_cost}(f) \;\defeq\;
  \sum_{s \in S_{\text{flight}}(f)} \text{rate}_s
    \cdot \sum_{a \in A^{\text{MAWB}} : f \in \text{legs}(a)}
    \sum_{g \in G_a} \sum_{u \in U_a} \eta_{a,g,u}
  \label{eq:flight-uld-surcharge}
\end{equation}
(where $S_{\text{flight}}(f)$ is the set of \texttt{per\_uld}-basis surcharges
applicable to flight $f$), is correct for single-leg arcs but
\emph{overcounts multi-leg through-arcs}: a ULD riding intact through a hub
is built and broken exactly once, yet summing this identity over the arc's
legs charges a build/breakdown at every intermediate departure. The MILP
therefore uses the per-arc-endpoint coefficient (Eq.~\ref{eq:sigma-uld})
below, not this per-flight sum.

\paragraph{What legitimately belongs in $S^{\text{build}}/S^{\text{break}}$ (grounding).}
Only costs genuinely proportional to ULD \emph{count} qualify: \textbf{ULD
build-up labor} (origin palletization/containerization) and \textbf{breakdown
labor} (destination de-palletization). These are discrete per-unit activities
(ULD~CARE / IATA Master Operating Plan). Three common mis-attributions are
explicitly \emph{excluded}: (i) the \textbf{linehaul / BUC rate} is per-kg with
a pivot, not per-ULD --- it lives in the bucket cost (\S\ref{sec:lin-bucket}),
and folding it here would drop the weight dependence and double-charge; (ii)
\textbf{terminal handling (THC)} is per-kg-with-minimum in standard tariffs,
carried as a Path-A surcharge, not per-ULD; (iii) \textbf{ULD demurrage /
control} is per-ULD-\emph{per-day} and conditional on late return --- a
separate, stochastic term, not a deterministic build coefficient (deferred).
Whether build-up is the forwarder's own CFS labor (Bulk Unitisation
Programme / BUP, forwarder-built) or a carrier/GHA charge (airline-built)
flips who pays; the MVP assumes a single mode per contract and documents it.
The $\sigma^{\text{uld}}_{a,u}$ coefficient is \textbf{synthetic / placeholder}
until sourced from a GHA rate card or TACT extract (per-ULD build/breakdown
labor rates are not generally public); \texttt{MARKET RESEARCH NEEDED}.

\paragraph{Objective form --- per-arc-endpoint coefficient on $\eta$.}
The per-ULD cost attaches to the ULD-count variable $\eta_{a,g,u}$ --- a
binary MAWB-activation $z_{a,g}$ cannot carry it, because build/breakdown
labor scales with the \emph{number} of ULDs, not merely with whether the
MAWB is active (the per-MAWB fixed overhead that \emph{does} scale with
$z_{a,g}$ is the separate $c^{\text{MAWB}}_{\text{fix}}$ term). The
coefficient is the build-up rate at the arc's origin plus the breakdown
rate at its destination, counted \emph{once} per MAWB-arc:
\begin{equation}
  \sigma^{\text{uld}}_{a,u} \;\defeq\;
    \sum_{s \in S^{\text{build}}_{\text{tail}(a)}} \text{rate}_s
    \;+\;
    \sum_{s \in S^{\text{break}}_{\text{head}(a)}} \text{rate}_s,
  \label{eq:sigma-uld}
\end{equation}
where $S^{\text{build}}_{n}$ and $S^{\text{break}}_{n}$ are the per-ULD
build-up and breakdown surcharges applicable at station $n$ for ULD type
$u$. The objective term is the manifestly linear $\sum_{(a,g) \in M,\,
u \in U_a} \sigma^{\text{uld}}_{a,u} \cdot \eta_{a,g,u}$
(\S\ref{sec:objective}). \textbf{Indexing by the arc endpoints
$\text{tail}(a), \text{head}(a)$ rather than by $\text{legs}(a)$ is
load-bearing}: a through-arc whose ULD is built and broken exactly once is
charged once, not once per leg. Segment-by-segment routing, where the ULD
is re-built at the hub, is represented as two separate single-leg MAWB-arcs
--- each with its own $\eta$ and its own $\sigma^{\text{uld}}$ --- so
genuine per-leg build/breakdown cycles are captured by the arc structure,
not by a per-leg coefficient. This keeps $F$ out of the MILP (it remains
graph-build-only) and surfaces the dependence on $\eta$.

\paragraph{Why the path split.}
Treating \texttt{per\_uld} as per-HAWB-per-arc was a bug in earlier drafts:
multiplying $\text{rate} \cdot \sum_u \eta$ by $x_{k,a}$ either over-charges
(every HAWB on the flight pays the full fee, $n$ times for $n$ HAWBs) or
introduces a bilinear $x \cdot \eta$ term requiring McCormick. Attribution to
the per-MAWB ULD-count variable $\eta_{a,g,u}$ (Eq.~\ref{eq:sigma-uld}) is
correct accounting and eliminates the bilinearity.

\paragraph{Common arc-type assignment.}
\begin{itemize}[noitemsep]
  \item Origin THC, AWB fee, origin screening, DGR handling: origin-side
        ground arcs (pickup, origin CFS dwell, origin cartage).
  \item Destination THC, AMS / ICS2 / ACI, destination handling: destination-side
        ground arcs (destination cartage, destination CFS dwell, customs
        dwell).
  \item FSC, SSC, war risk, PSS / GRI, percent-of-freight: air arcs (MAWB and
        co-load).
\end{itemize}
The \texttt{applies\_on\_arc\_types} scope field controls the gating; the
resolver attributes the cost only to matching arcs.

\paragraph{Worked example.}
For a 1{,}200 kg GEN shipment TPE $\to$ JFK on a forwarder-consolidated
TPE--HKG--JFK same-carrier-connection MAWB-arc (one MAWB at 1{,}200 kg
chargeable; PRM\_AIR\_EXP, May 2026 outside PSS season):
\begin{table}[H]
\centering
\caption{Surcharge worked example: 1{,}200~kg GEN TPE$\to$JFK on \texttt{PRM\_AIR\_EXP}.}
\label{tab:surcharge-example}
\begin{tabular}{lr}
\toprule
Surcharge & Per-arc (MAWB-arc) (USD) \\
\midrule
FSC (0.42/kg) & 1{,}008 \\
SSC (0.18/kg) & 432 \\
THC origin (Path A on origin cartage arc) & 65 \\
THC destination (Path A on destination cartage arc) & 90 \\
AMS (US filing, destination side) & 25 \\
AWB fee (Path A on origin CFS dwell arc) & 35 \\
Screening (0.06/kg, min 35) & 72 \\
\midrule
Subtotal & 1{,}727 \\
\bottomrule
\end{tabular}
\end{table}
$\sim 45\%$ uplift on a base rate of $\$3{,}500$--$\$4{,}000$ --- the MILP
must see this to quote correctly. See \texttt{data\_model.md \S6.6} for the
JSON catalog rows.

\subsection{ULD interchange and re-ULDing --- absorbed at graph build}
\label{sec:through-uld}

The decision of whether a built ULD can ride through a hub (``through-ULD'')
vs.\ being broken down and rebuilt (``re-ULD'') is a property of the
underlying flight pair and ULD type, evaluated by the graph generator. In v3
this is \emph{not} a MILP variable; instead the generator emits arcs whose
structure reflects the policy outcome.

\paragraph{ULD interchange set $\Pi$.}
$\Pi \subseteq \{(c_1, c_2, u)\}$ collects carrier-pair $\times$ ULD-type
triples for which $c_1$ and $c_2$ have a bilateral or alliance-level IATA
ULD-CPM agreement covering type $u$:
\begin{itemize}[noitemsep]
  \item Star Alliance ULD pool: AC, LH, NH, OZ, SK, UA, \ldots
  \item SkyTeam Cargo ULD pool: AF, KL, DL, KE, MU, \ldots
  \item oneworld Cargo ULD pool: AA, BA, CX, JL, QF, \ldots
  \item Bilateral non-alliance agreements.
\end{itemize}

\paragraph{Emission constructs (not mutually-exclusive alternatives).}
The four constructs below describe \emph{how} the generator emits arcs in
each through-vs-no-interchange configuration. They are \emph{not} a choice
between four output patterns for one physical flight --- per the
overlapping-emission policy of \S\ref{sec:arc-enumeration}, the generator
may emit \emph{multiple} of these constructs in parallel for the same
multi-leg physical flight when the catalog supplies offers under more than
one construct. The MILP then picks among the emitted candidates via
$x_{k, a}$. For each multi-leg routing $\text{POL} \to H \to \text{POD}$ on
operating carriers $(c_1, c_2)$ and ULD type $u$:
\begin{enumerate}[noitemsep,label=(\alph*)]
  \item \textbf{Through-flight (technical stop), same flight number, same
        $c$.} ULD never leaves the aircraft. Generator emits a
        \emph{single} MAWB-arc $\text{POL} \to \text{POD}$; $\mu_a$
        includes the technical stop turnaround. Co-existing constructs
        (b)--(d) may also be emitted from the same physical flight if the
        carrier separately tariffs the legs.
  \item \textbf{Same-airline through-cargo at hub} ($c_1 = c_2$ and the
        hub supports through-cargo for $u$). When the carrier prices the
        through-segment, generator emits a single MAWB-arc $\text{POL} \to
        \text{POD}$; when (also) priced per leg, generator emits the two
        single-leg arcs $\text{POL} \to H$ and $H \to \text{POD}$ joined
        by a carrier-side connection dwell arc \emph{as parallel
        candidates}. Both can co-exist when both rate structures are in
        the catalog.
  \item \textbf{Interline through-AWB} (Cargo SPA / MITA / alliance pool;
        $(c_1, c_2, u) \in \Pi$ and the hub supports cross-airline
        transfer). Generator emits a single MAWB-arc spanning the
        carriers (one AWB, one rate, one MAWB), with $\mu_a$ including
        the cross-airline through-cargo time. Single-leg fallback arcs
        may co-exist if the legs are also separately tariffed.
  \item \textbf{No interchange.} Generator emits two single-leg MAWB-arcs
        joined by a deconsolidation-dwell arc (if at a $\text{CFS-H}$) or
        a carrier-side connection-dwell arc (otherwise). The dwell's
        $\delta_a$ carries the longer re-ULDing time and the handling
        charge $c_a^{\text{flat}}$ carries the rehandling fee. No
        through-MAWB candidate exists in this case.
\end{enumerate}
In all cases the MILP sees only arcs and scalars --- no $\psi$ parameter,
no through-ULD decision variable. The cost differential between
constructs (rehandling fees \$150--\$500/ULD; dwell time differential
4--8 hours) is preserved in the per-arc data emitted to the MILP, and the
MILP's $x_{k, a}$ choice over the emitted candidates realizes the
selection between through and segment-by-segment routings.

%%--------------------------------------------------------------------
\section{Decision Variables, Sets and Indices}
\label{sec:variables}
\label{sec:sets}
%%--------------------------------------------------------------------

This section is the single source of truth for the MILP's structural symbols:
decision variables, sets, indices, and the function-style names that appear
across constraints, objective, and linearization. Equation-heavy downstream
sections recap the subset they need with cross-references back here. Some
sets (e.g.\ $C_k^{\text{allow}}$) are formally constructed in later
sections; their rows below forward-reference those sections.

\paragraph{Notation convention.}
Function-style names with arguments --- $\text{transit}(k, a)$,
$\text{prep\_time}_k$, $\text{surcharge\_cost}(k, a)$,
$\text{flight\_uld\_surcharge\_cost}(f)$, $\text{cargo\_type\_ok}(k, a)$,
$\text{embargo\_ok}(k, a)$, $\text{lithium\_ok}(k, a)$,
$\text{mode\_ok}(k, a)$, $\text{carrier\_ok}(k, a)$,
$\text{ac\_type\_ok}(k, a)$ --- denote \emph{pre-computable constants}, not
decision variables. They are evaluated during ingestion / graph construction /
subgraph pre-filtering from shipment attributes and supply data, and enter
the MILP as coefficients or predicates.

\subsection*{Sets and indices}

\begin{table}[H]
\centering
\caption{Sets and indices used by the MILP.}
\label{tab:sets-indices}
\begin{tabular}{lp{11.5cm}}
\toprule
Symbol & Description \\
\midrule
$K$ & Set of HAWBs (commodities, shipments) to route this solve, indexed by $k$ \\
$N,\, A$ & All nodes and arcs of graph $G(N,A)$ \\
$\text{tail}(a),\, \text{head}(a)$ & Endpoints of arc $a$ --- $a$ carries flow from $\text{tail}(a)$ to $\text{head}(a)$ \\
$A^{\text{MAWB}} \subseteq A$ & MAWB-eligible air arcs --- $\text{rate\_family}_a \in \{\texttt{flat\_rate}, \texttt{min\_flat\_breaks}, \texttt{per\_uld\_pivot}\}$ \\
$A^{\text{MFB}} \subseteq A^{\text{MAWB}}$ & Arcs with $\text{rate\_family}_a = \texttt{min\_flat\_breaks}$ \\
$A^{\text{coload}} \subseteq A$ & Co-load air arcs --- per-kg billing, no MAWB \\
$A^{\text{ground}} \subseteq A$ & Non-air arcs: pickup, cartage, CFS dwell, deconsol-dwell, customs-dwell, delivery \\
$A^{\text{cust}} \subseteq A^{\text{ground}}$ & Customs-clearance dwell arcs (between $\text{CFS-D}$ and final delivery); one per HAWB per route \\
$A_k \subseteq A$ & Pre-filtered subgraph of HAWB $k$ (\S\ref{sec:prefilter}) \\
$N_k \subseteq N$ & Nodes incident to $A_k$, with $O_k \in N_k$ and $D_k^{\text{node}} \in N_k$ guaranteed by the pre-filter \\
$\text{transit}(k, a)$ & Per-HAWB transit time on arc $a$ (hours): for ground arcs, $\delta_a + \delta^{\text{cust}}_k \cdot \mathbbm{1}[a \in A^{\text{cust}}]$; for air arcs, $\mu_a$ (a graph-build-precomputed scalar) \\
$F$ & Set of physical flight legs (graph-build only --- the MILP does not index over $F$) \\
$\text{legs}(a)$ & Physical flight legs realizing air arc $a$ (graph-build only; the MILP sees only the precomputed scalars $\mu_a, \text{CO}_a^*$) \\
$U$ & Set of ULD types $\{\text{LD3, LD7, PMC, AKE,} \ldots\}$, indexed by $u$ \\
$U_a \subseteq U$ & ULD types the forwarder can claim on \texttt{per\_uld\_pivot} arc $a$. Determined by the forwarder's BSA contract allotment ($N_{a,u} > 0$), bounded above by the ULD positions the underlying aircraft physically supports. Single-forwarder solve: no cross-tenant pooling \\
$C$ & Set of \texttt{per\_uld\_pivot} (BSA) contracts, indexed by $c$ \\
$C^{\text{eq}} \subseteq C$ & Equalized-settlement BSA contracts \\
$A_c^{\text{MAWB}} \subseteq A^{\text{MAWB}}$ & Arcs tagged to BSA contract $c$ \\
$A^{\text{pu}} = \bigcup_{c \in C} A_c^{\text{MAWB}} \subseteq A^{\text{MAWB}}$ & Per-ULD-pivot arcs \\
$g(k)$ & Group key of HAWB $k$ as defined by the function $g(\cdot)$ (\S\ref{sec:g-of-k}); the range of $g$ is only ever iterated via $G_a$ (below), never globally \\
$K_a = \{k : a \in A_k\}$ & HAWBs whose pre-filtered subgraph contains arc $a$. Computed at MILP build by inverting the per-shipment $A_k$ relation across $k$. Solve-specific (depends on this solve's HAWB batch) \\
$G_a = \{g(k) : k \in K_a\}$ & Distinct group keys present on arc $a$, derived by projecting $K_a$ through $g(\cdot)$. If 6 distinct group keys appear across the HAWBs whose subgraph contains $a$, then $|G_a| = 6$. Not chosen by the optimizer; computed at build time. Solve-specific \\
$M = \{(a, g) : a \in A^{\text{MAWB}},\, g \in G_a\}$ & MAWB candidates --- the central index set. One $(a, g)$ pair per (MAWB-eligible arc) $\times$ (group present on that arc), enumerated at build. Solve-specific \\
$B_a$ & Weight-break segments of a \texttt{min\_flat\_breaks} arc $a$, ordered ascending; indexed by $b$ \\
$P$ & Tenant catalog of service products, indexed by $p$ (\S\ref{sec:service-products}) \\
$\text{sp}(k) \in P$ & Service product bound to HAWB $k$ \\
$C_k^{\text{allow}},\, C_k^{\text{deny}}$ & Resolved hard-allowed / hard-denied carrier sets for HAWB $k$ from the policy cascade (\S\ref{sec:carrier-policy}, Eq.~\ref{eq:carrier-allow}) \\
$\text{chargeable}(c)$ & Derived continuous (\(\geq 0\)) for $c \in C^{\text{eq}}$: sum of MAWB-level chargeable weights tendered on contract $c$ this solve, $\sum_{(a,g):\, a \in A_c^{\text{MAWB}},\, g \in G_a} \text{CW}_{a,g}$. Used in C.13a to compute overage. \\
\bottomrule
\end{tabular}
\end{table}

\paragraph{Catalog sets vs.\ solve-specific sets.}
Two layers of statefulness in the sets above:
\begin{itemize}[noitemsep]
  \item \textbf{Catalog sets} are static across solves --- they come from
        the forwarder's supply catalog and configuration: $N, A, A^{\text{MAWB}},
        A^{\text{MFB}}, A^{\text{coload}}, A^{\text{ground}}, A^{\text{cust}},
        U, C, C^{\text{eq}}, A_c^{\text{MAWB}}, A^{\text{pu}}, B_a, P$. The
        graph generator builds them once per ingestion cycle.
  \item \textbf{Solve-specific sets} depend on \emph{this} solve's HAWB batch
        $K$ and are rebuilt at MILP build: $K, A_k, N_k, K_a, G_a, M,
        C_k^{\text{allow}}, C_k^{\text{deny}}$. A different
        batch produces a different $M$ even on identical supply. $K^{\text{fb}}$
        is post-solve, derived from the optimal $x_{k, a^{\text{fb}}_k}$.
\end{itemize}
The derivation chain on the solve-specific side is $A_k \;\to\; K_a \;\to\; G_a
\;\to\; M$: per-shipment pre-filter $\to$ invert across $k$ $\to$ project via
$g(\cdot)$ $\to$ enumerate $(a, g)$ candidates. The optimizer chooses
$x_{k,a}$ and $z_{a,g}$; it does not choose $A_k, K_a, G_a, M$.

\subsection*{Decision variables}

\begin{table}[H]
\centering
\caption{Decision variables of the MILP.}
\label{tab:decision-variables}
\begin{tabular}{p{3.5cm}p{2.5cm}p{6cm}p{1.5cm}}
\toprule
Variable & Type & Description & Index domain \\
\midrule
$x_{k,a}$ & binary & HAWB $k$ uses arc $a$ & $k \in K$, $a \in A_k$ \\
$z_{a,g}$ & binary & MAWB $(a, g)$ is instantiated (active) & $(a, g) \in M$ \\
$\text{CW}_{a,g}$ & continuous $\geq 0$ & Chargeable weight booked on MAWB $(a, g)$ & $(a, g) \in M$ \\
$\text{Wt}_{a,g}$ & continuous $\geq 0$ & Aggregate actual weight on MAWB $(a, g)$ (with dunnage) & $(a, g) \in M$ \\
$\text{Wv}_{a,g}$ & continuous $\geq 0$ & Aggregate volumetric weight on MAWB $(a, g)$ & $(a, g) \in M$ \\
$\gamma_{a,g,b}$ & binary & Break selector for MAWB $(a, g)$ at break $b$ (\texttt{min\_flat\_breaks} arcs only) & $(a, g) \in M,\, a \in A^{\text{MFB}},\, b \in B_a$ \\
$\text{BW}_{a,g,b}$ & continuous $\geq 0$ & Disaggregated billed weight at break $b$ & as $\gamma$ \\
$\eta_{a,g,u}$ & integer $\geq 0$ & ULDs of type $u$ claimed for MAWB $(a, g)$ & $(a, g) \in M,\, a \in A^{\text{pu}},\, u \in U_a$ \\
$\text{pivot}_{a,g}$ & continuous $\geq 0$ & Pivot-active billed weight $\max(\text{CW}_{a,g},\, \pi_a \cdot \sum_u \eta_{a,g,u})$ & $(a, g) \in M,\, a \in A^{\text{pu}},\, \text{settlement}_a = \texttt{per\_flight}$ \\
$\text{over}_c$ & continuous $\geq 0$ & Allowance overage for equalized contract $c$: $\max(0,\, \sum_{(a,g)} \text{CW}_{a,g} - A_c)$ & $c \in C^{\text{eq}}$ \\
$t_k^{D_k^{\text{node}}}$ & continuous $\geq 0$ & Arrival time of HAWB $k$ at its destination node; defined by C.10a from precomputed $\text{arr\_dest}(k, a)$ scalars. Intermediate-node arrival times do not appear in the MILP --- see \S\ref{sec:fwd-time-propagation}. & $k \in K$ \\
$\tau_k$ & continuous $\geq 0$ & Tardiness of HAWB $k$ against $\Delta_k$ & $k \in K$ \\
$\text{pen}_k$ & continuous $\geq 0$ & PWL outer-approximation of the quadratic tardiness penalty $W_k \cdot \tau_k^2$; appears in the objective in place of $\tau_k^2$ (\S\ref{sec:lin-tardiness}) & $k \in K$ \\
\bottomrule
\end{tabular}
\end{table}

\paragraph{Domain bounds.}
See C.14 (\S\ref{sec:domain}). The per-MAWB chargeable-weight upper bound
$\text{CW}^{\text{ub}}_{a,g}$ precomputed at MILP build,
\[
  \text{CW}^{\text{ub}}_{a,g} \;\defeq\; (1 + \varepsilon) \sum_{k \in K_a : g(k) = g} \max\bigl(w_k,\, v_k \cdot 167\bigr),
\]
ties bucket cost variables to MAWB activation (empty bucket $\Rightarrow
\text{CW} = 0$) and provides tight big-M values for the break-selection
disaggregation (\S\ref{sec:bigm-tightening}).

\subsection*{Post-solve derived quantities}

The optimal solution is post-processed to surface the diagnostic quantities
used by the orchestrator (\S\ref{sec:output-diagnostics}). They are not
decision variables; they are functions of $x, z, \text{pen}, \text{over},
\ldots$ at integer feasibility.

\begin{table}[H]
\centering
\caption{Post-solve derived quantities surfaced for diagnostics and operator-facing screens.}
\label{tab:post-solve-derived}
\begin{tabular}{lp{11.5cm}}
\toprule
Symbol & Description \\
\midrule
$K^{\text{fb}}$ & HAWBs routed via the fallback arc: $\{k : x_{k, a^{\text{fb}}_k} = 1\}$ \\
$z^{\text{routing}}$ & Real routing cost: sum of all objective terms \emph{excluding} fallback-arc cost and tardiness penalty. The number operators see on planning screens as ``cost to serve.'' \\
$z^{\text{tardiness}}$ & Total quadratic tardiness penalty $\sum_k \text{pen}_k$ (PWL approximation of $\sum_k W_k \tau_k^2$) \\
$z^{\text{fallback}}$ & Fallback usage cost $C^{\text{fallback}} \cdot |K^{\text{fb}}|$ \\
$z^{\text{total}}$ & Solver objective $= z^{\text{routing}} + z^{\text{tardiness}} + z^{\text{fallback}}$ \\
\bottomrule
\end{tabular}
\end{table}

%%--------------------------------------------------------------------
\section{Embargo Gating}
\label{sec:embargo}
%%--------------------------------------------------------------------

Embargoes are carrier-, route-, commodity-, and time-bounded prohibitions on
cargo acceptance. They are imposed by carriers (operational policy, security
incidents, peak capacity), by regulators (geopolitical, sanctions), or by
IATA / industry notice. The model represents embargoes as a parameter set
used at subgraph construction, not as MILP constraints --- the optimizer never
sees embargoed arcs.

\paragraph{Symbols used in this section.}
\begin{table}[H]
\centering
\caption{Symbols used in the embargo-gating section.}
\label{tab:embargo-syms}
\begin{tabular}{lp{9cm}p{3.5cm}}
\toprule
Symbol & Brief & Defined \\
\midrule
$E$ & Catalog of embargo records (one tuple per record) & this section (record schema) \\
$E_f$ & Embargoes active on flight $f$ at $\text{ETD}_f$ & Eq.~\ref{eq:active-embargoes} \\
$\text{match}(e, k, f)$ & Predicate: embargo $e$ applies to HAWB $k$ on flight $f$ & this section (Match predicate) \\
$\text{embargo\_ok}(k, a)$ & No active embargo matches $k$ on any leg of $a$ & Eq.~\ref{eq:embargo-feasibility} \\
$\text{legs}(a)$ & Physical flight legs of air arc $a$ & \S\ref{sec:variables} \\
$k,\, a,\, f$ & HAWB / arc / flight indices & \S\ref{sec:variables} \\
\bottomrule
\end{tabular}
\end{table}

\paragraph{Embargo record schema.}
Let $E$ be the set of embargo records. Each $e \in E$ is a tuple:
\begin{table}[H]
\centering
\caption{Embargo record schema (one tuple per record).}
\label{tab:embargo-schema}
\begin{tabular}{lp{9cm}}
\toprule
Field & Description \\
\midrule
$\text{carrier}(e)$ & IATA carrier code, or \texttt{*} wildcard \\
$\text{origin}(e)$ & Origin airport code, region prefix (e.g., \texttt{CN-*}), or \texttt{*} \\
$\text{dest}(e)$ & Destination airport code, region prefix, or \texttt{*} \\
$\text{hub}(e)$ & Transit hub airport code, or $\emptyset$ (applies regardless of routing) \\
$\text{ac\_type}(e)$ & \texttt{FREIGHTER} / \texttt{PAX\_BELLY} / \texttt{*} \\
$\text{cargo\_type}(e)$ & \texttt{GEN} / \texttt{DGR} / \texttt{PER} / \texttt{VAL} / \texttt{AVI} / \texttt{HUM} / \texttt{*} \\
$\text{commodity}(e)$ & 4-digit IATA commodity code, or \texttt{*} \\
$\text{dgr\_class}(e)$ & DGR class+subclass (e.g., \texttt{9-II}), or \texttt{*} \\
$\text{lithium\_pi}(e)$ & PI965 / PI966 / PI967 / PI968 / PI969 / PI970 / $\emptyset$ (\S\ref{sec:lithium}) \\
$\text{start}(e), \text{end}(e)$ & Date-time validity range \\
$\text{source}(e)$ & \texttt{carrier\_notice} / \texttt{iata} / \texttt{regulatory} / \texttt{tenant\_internal} \\
\bottomrule
\end{tabular}
\end{table}

\paragraph{Active embargoes for an arc.}
For each MAWB-arc or co-load arc $a$ and each underlying leg $f \in
\text{legs}(a)$:
\begin{equation}
  E_f \;\defeq\; \{ e \in E : \text{start}(e) \leq \text{ETD}_f \leq \text{end}(e) \}
  \label{eq:active-embargoes}
\end{equation}

\paragraph{Match predicate.}
Embargo $e$ matches HAWB $k$ on arc $a$ via leg $f$ iff every non-wildcard
field of $e$ matches the corresponding attribute --- carrier and aircraft type
from $f$; origin and destination from $a$'s endpoints; transit hub from the
intermediate nodes of $a$; cargo type / commodity / DGR class / lithium PI
from $k$'s attributes.

\paragraph{Arc embargo-feasibility.}
\begin{equation}
  \text{embargo\_ok}(k, a) \;\defeq\;
    \forall f \in \text{legs}(a):\;
    \neg \, \exists\, e \in E_f : \text{match}(e, k, f)
  \label{eq:embargo-feasibility}
\end{equation}
Evaluated at subgraph construction (\S\ref{sec:prefilter}); embargoed arcs are
pruned before the MILP.

\paragraph{Examples (illustrative records).}
\begin{table}[H]
\centering
\caption{Illustrative embargo records.}
\label{tab:embargo-examples}
\begin{tabular}{lp{10.5cm}}
\toprule
Type & Record \\
\midrule
Carrier-commodity & $\text{carrier} = \texttt{CX},\, \text{origin} = \texttt{TPE},\,
                    \text{dest} = \texttt{US-*},\, \text{hub} = \texttt{HKG},\,
                    \text{lithium\_pi} = \texttt{PI965},\,
                    \text{start} = 2023\text{-}06\text{-}01,\,
                    \text{end} = 2099\text{-}12\text{-}31$ \\
Seasonal Hajj & $\text{carrier} = \texttt{EK},\, \text{dest} = \texttt{JED},\,
                 \text{cargo\_type} = \texttt{GEN},\,
                 \text{start} = 2026\text{-}06\text{-}10,\,
                 \text{end} = 2026\text{-}07\text{-}22$ \\
CNY perishables & $\text{origin} = \texttt{HKG},\, \text{cargo\_type} = \texttt{PER},\,
                   \text{start} = 2026\text{-}02\text{-}08,\,
                   \text{end} = 2026\text{-}02\text{-}22$ \\
Aircraft-type DGR & $\text{carrier} = \texttt{*},\, \text{ac\_type} = \texttt{PAX\_BELLY},\,
                    \text{lithium\_pi} = \texttt{PI965},\,
                    \text{start} = 2024\text{-}04\text{-}01,\,
                    \text{end} = 2099\text{-}12\text{-}31$ \\
\bottomrule
\end{tabular}
\end{table}

\paragraph{Sourcing and scope.}
For MVP, $E$ is a manually-curated table maintained by tenant ops, seeded from
carrier embargo portal pages, IATA alerts (TACT supplements), and operational
knowledge. P1: agent capability to parse carrier notices into structured
records with human review. Embargoes are global per tenant; tenant-specific
carrier exclusions are handled separately via the carrier-policy cascade
(\S\ref{sec:carrier-policy}, distinct from embargoes). All embargoes are hard
prohibitions; soft / capacity-driven restrictions are handled in the rate /
capacity layer.

%%--------------------------------------------------------------------
\section{Lithium Battery Taxonomy and Per-Flight Acceptance}
\label{sec:lithium}
%%--------------------------------------------------------------------

Lithium batteries account for an estimated 30--40\% of TPE/HKG eastbound air
cargo in 2026. They are not a single commodity class: IATA DGR distinguishes
six packing instructions and three section tiers per PI, and each carrier
publishes per-route acceptance rules that vary by PI, section, and aircraft
type. Modeling lithium as a generic DGR boolean conflates flights that accept
PI967-contained-in-equipment with those that refuse PI965 IA standalone
cells.

\paragraph{Symbols used in this section.}
\begin{table}[H]
\centering
\caption{Symbols used in the lithium taxonomy section.}
\label{tab:lithium-syms}
\begin{tabular}{lp{9cm}p{3.5cm}}
\toprule
Symbol & Brief & Defined \\
\midrule
$\text{lithium\_spec}_k$ & Per-HAWB lithium attributes (UN \#, PI, Section, Wh, SOC, DDR); $\emptyset$ for non-lithium & this section \\
$\text{lithium\_accept}_f$ & Per-flight acceptance matrix on (PI, Section, $\text{ac\_type}$) & Eq.~\ref{eq:lithium-accept} \\
$\text{lithium\_ok}(k, a)$ & HAWB $k$ acceptable on every leg of $a$ & Eq.~\ref{eq:lithium-ok} \\
$\text{ac\_type}(f)$ & Aircraft type of flight $f$: \texttt{FREIGHTER} / \texttt{PAX\_BELLY} & \S\ref{sec:air-arc-params} \\
$\text{legs}(a)$ & Physical flight legs of air arc $a$ & \S\ref{sec:variables} \\
$k,\, a,\, f$ & HAWB / arc / flight indices & \S\ref{sec:variables} \\
\bottomrule
\end{tabular}
\end{table}

\paragraph{IATA lithium taxonomy.}
\begin{table}[H]
\centering
\caption{IATA lithium battery packing-instruction taxonomy.}
\label{tab:lithium-taxonomy}
\begin{tabular}{p{2cm}p{3.5cm}p{4.2cm}p{1.5cm}}
\toprule
UN \# & Chemistry & Configuration & PI \\
\midrule
UN3480 & Li-ion    & Standalone & PI965 \\
UN3481 & Li-ion    & Packed with equipment & PI966 \\
UN3481 & Li-ion    & Contained in equipment & PI967 \\
UN3090 & Li-metal  & Standalone & PI968 \quad (forbidden on passenger aircraft globally since 2015) \\
UN3091 & Li-metal  & Packed with equipment & PI969 \\
UN3091 & Li-metal  & Contained in equipment & PI970 \\
\bottomrule
\end{tabular}
\end{table}

Within each PI, the IATA Section tier determines regulatory burden:
\begin{description}[leftmargin=2cm, labelwidth=1.8cm, style=nextline]
  \item[Section IA] Large cells/batteries ($>100$ Wh Li-ion, $>2$ g Li metal).
                    Fully regulated.
  \item[Section IB] $\leq 100$ Wh Li-ion / $\leq 2$ g Li metal in larger
                    packaging. Partial relief.
  \item[Section II] Small batteries below threshold quantities. Most
                    operational relief, still regulated.
\end{description}

\paragraph{Per-HAWB lithium attributes.}
$\text{lithium\_spec}_k$ ($\emptyset$ for non-lithium):
\begin{table}[H]
\centering
\caption{Per-HAWB lithium specification fields.}
\label{tab:lithium-spec}
\begin{tabular}{lp{9cm}}
\toprule
Field & Description \\
\midrule
$\text{un\_number}$ & \texttt{UN3480} / \texttt{UN3481} / \texttt{UN3090} / \texttt{UN3091} \\
$\text{pi\_code}$ & PI965 / PI966 / PI967 / PI968 / PI969 / PI970 \\
$\text{section}$ & IA / IB / II \\
$\text{watt\_hours}$ & Maximum Wh per battery \\
$\text{soc\_compliant}$ & Boolean: shipper-certified SOC $\leq 30\%$ of rated capacity (required for UN3480, UN3090 standalone) \\
$\text{ddr}$ & Damaged / defective / recalled. Universally refused in MVP. \\
\bottomrule
\end{tabular}
\end{table}

\paragraph{Per-flight acceptance matrix.}
For each flight $f$ with $\text{ac\_type}(f) \in \{\texttt{FREIGHTER},
\texttt{PAX\_BELLY}\}$, the carrier publishes:
\begin{equation}
  \text{lithium\_accept}_f : (\text{PI}, \text{Section}, \text{ac\_type}) \to \{\text{True}, \text{False}\}
  \label{eq:lithium-accept}
\end{equation}
\textbf{Whitelist semantics:} default is \texttt{False} (refuse). Cost
asymmetry (false-accept implies 48 h roll, fire-risk liability) justifies a
default-deny posture.

\paragraph{Lithium feasibility predicate.}
For HAWB $k$ with $\text{lithium\_spec}_k \neq \emptyset$ and arc $a$:
\begin{equation}
  \begin{aligned}
    \text{lithium\_ok}(k, a) \;\defeq\;
    & \neg\, \text{lithium\_spec}_k.\text{ddr} \\
    \wedge\; & \bigl(\text{lithium\_spec}_k.\text{un\_number} \in \{\texttt{UN3480}, \texttt{UN3090}\}
                 \Rightarrow \text{lithium\_spec}_k.\text{soc\_compliant}\bigr) \\
    \wedge\; & \forall f \in \text{legs}(a):\;
                \text{lithium\_accept}_f\bigl[\text{lithium\_spec}_k.\text{pi\_code},\,
                  \text{lithium\_spec}_k.\text{section},\, \text{ac\_type}(f)\bigr] = \text{True}
  \end{aligned}
  \label{eq:lithium-ok}
\end{equation}
For non-lithium HAWBs, $\text{lithium\_ok}(k, a) \equiv \text{True}$ trivially.

\paragraph{Interaction with embargoes.}
Lithium acceptance is the \emph{positive} statement; embargoes are
\emph{negative}, possibly time-bounded. Both must hold: an arc is
lithium-feasible iff acceptance holds AND no active embargo matches. AND-ed
in the pre-filter (\S\ref{sec:prefilter}).

\paragraph{Per-flight aggregate lithium caps --- carrier-side, not modeled.}
Carriers publish per-flight aggregate limits on lithium quantity (e.g., max
$N$ kg PI965~II per flight) and enforce these at acceptance across all
forwarders booking the same flight. A single forwarder cannot model the
aggregate because they do not see other forwarders' bookings. The MILP's job
is correct per-HAWB declaration and acceptance check (above); aggregate
enforcement is the carrier's. If a forwarder's HAWB is bumped at acceptance
because the carrier's aggregate cap was reached after booking, this surfaces
as a booking-rejection event handled by the recovery workflow
(\S\ref{sec:locked-commitments}).

\paragraph{Sourcing.}
For MVP, $\text{lithium\_accept}_f$ is manually curated per carrier, seeded
from carrier cargo portals and IATA DGR section 9 (state and operator
variations). Annual refresh plus event-driven updates after carrier policy
changes (typically following lithium fire incidents).

\paragraph{Scope decisions for MVP.}
\begin{itemize}[noitemsep]
  \item All six PIs and three sections supported.
  \item SOC compliance is shipper-certified boolean; verification deferred.
  \item UN 38.3 test summary validation deferred (assumed valid in MVP).
  \item Per-flight aggregate lithium caps: carrier-side, not in MVP.
  \item PI965 Section IA double-packaging requirements handled at the CFS
        layer, not at routing.
  \item DDR universally refused; DDR specialty routing deferred.
  \item Out of scope: battery recycling (PI909), vehicle batteries (PI952),
        emerging chemistries (sodium-ion, solid-state).
\end{itemize}

%%--------------------------------------------------------------------
\section{Cargo Screening --- Cost, Not Eligibility}
\label{sec:screening}
%%--------------------------------------------------------------------

Regulators in the US (TSA CCSP, 49 CFR 1548/1549) and EU (Commission
Implementing Regulation 2015/1998 ACC3/RA3/KC3) mandate that cargo on
passenger flights into / through their airspace be screened (X-ray, ETD,
physical, or canine) prior to loading.\footnote{TSA 49 CFR 1549 Certified
Cargo Screening Program; Commission Implementing Regulation (EU) 2015/1998.}
Post-2010, US-destined PAX cargo requires 100\% physical screening regardless
of Known Shipper status. The practical consequence: in any commercially
feasible routing, every air-borne shipment is screened. Unscreened cargo is
not tendered.

\paragraph{Modeled as cost, not as eligibility.}
Because screening is universal in feasible operations, an arc-eligibility
predicate $\text{screening\_ok}(k, a)$ would be a tautology --- satisfied by
every shipment in any feasible solution. It adds variables and pre-filter
state without ever cutting a branch. Instead, the per-kg screening cost
($\approx \$0.05$--$0.30$/kg depending on facility type) is captured by the
existing surcharge mechanism (\S\ref{sec:surcharge}, Path A
\texttt{per\_kg} basis on origin-side ground arcs). No $\text{screening}(k)$
attribute on HAWBs; no $\text{screening\_ok}$ predicate; no
$\mathcal{R}^{\text{screen}}$ region set. The Screening line in the
\S\ref{sec:surcharge} worked example ($0.06$/kg, min \$35) is the canonical
encoding.

\paragraph{Deferred: multi-jurisdiction re-screening at transit.}
A multi-region routing (e.g., HKG $\to$ USA $\to$ EU) may require re-screening
at the US transit point to satisfy EU ACC3 chain-of-custody --- a
\emph{path-dependent} cost that cannot be cleanly localized to a single arc in
an arc-based formulation. This is intentionally not modeled in v3. Materiality
on transpacific lanes: the re-screen surcharge ($\approx \$0.10$--$0.15$/kg)
is $\lesssim 2\%$ of an all-in air rate of $\$6$--$\$9$/kg and narrower than
base-rate variability, so it does not flip routes. Revisit when migrating to a
\textbf{path-based formulation} (column generation / branch-and-price): a path
variable naturally carries the jurisdictional sequence, and a per-path
re-screen cost can be loaded at column-generation time without state-tracking
in the master problem.

%%--------------------------------------------------------------------
\section{Locked Commitments and Execution State}
\label{sec:locked-commitments}
%%--------------------------------------------------------------------

At any planning moment the rolling-horizon optimizer is invoked over a window
that may contain HAWBs at various stages of execution. Pickup may already
have dispatched; cargo may sit at the origin CFS; AWBs may be filed; ULDs may
be loaded on aircraft. Decisions \emph{already taken} cannot be unwound by the
optimizer. Locks are resolved at the \emph{preprocessing} layer
(\(\to\) C.12) --- the MILP itself is \textbf{lock-agnostic}. This section
defines the preprocessing semantics; the MILP body has no locked-commitments
constraint family.

\paragraph{Symbols used in this section.}
\begin{table}[H]
\centering
\caption{Symbols used in the locked-commitments section.}
\label{tab:locked-syms}
\begin{tabular}{lp{9cm}p{3.5cm}}
\toprule
Symbol & Brief & Defined \\
\midrule
$K^{\text{loc}} \subseteq K$ & HAWBs with at least one locked arc on entry to the solve & this section \\
$A_k^{\text{loc}} \subseteq A_k$ & Locked-arc prefix for HAWB $k$: arcs already executed or contractually committed & this section \\
$t_k^{\text{obs}}$ & Observed arrival time of HAWB $k$ at the current observed node (re-pointed origin) & this section \\
$\mathcal{E}^{\text{sup}}$ & Supply-side event stream invalidating locks (cancellation, equipment swap, allocation pull, cutoff shift) & this section \\
$K,\, A_k,\, O_k$ & HAWB set / pre-filtered subgraph / origin door & \S\ref{sec:variables} \\
\bottomrule
\end{tabular}
\end{table}
Fully locked HAWBs preprocess out of $K$ entirely (the MILP does not see
them). Partially locked HAWBs enter $K$ with origin re-pointed to the
observed node, $A_k$ truncated to forward arcs, and the cargo-ready window
at the re-pointed origin collapsed to $\{t_k^{\text{obs}}\}$.

\paragraph{Lifecycle states and lock posture.}
The 7-state shipment lifecycle (\texttt{data\_model.md}) maps to lock posture
as follows:
\begin{table}[H]
\centering
\caption{Shipment lifecycle states and lock posture --- the 7-state lifecycle (\texttt{data\_model.md}) mapped to MILP visibility.}
\label{tab:lifecycle-states}
\begin{tabular}{lp{6.8cm}p{4.5cm}}
\toprule
Lifecycle state & Lock semantics & MILP visibility \\
\midrule
\texttt{UNROUTED} & No commitments; fully open. & Full decision space \\
\texttt{SOFT\_PLANNED} & Tentative routing for cost preview; revocable. & Full decision space; current plan can be used as warm-start \\
\texttt{FIRM\_DEADLINE} & Dry-run window closed; about to commit. & Open if window not expired; near-locked otherwise \\
\texttt{FIRM\_PLANNED} & Routing committed to carriers; bookings issued. & Locked prefix = full intended route until execution events update it \\
\texttt{IN\_TRANSIT} & Cargo physically moving; arcs lock as events arrive. & Locked prefix grows monotonically \\
\texttt{DESTINATION\_PLANNING} & All air arcs done; destination ground arcs still open. & Air arcs locked; destination cartage / CFS / customs / delivery open \\
\texttt{DELIVERED} & Terminal. & Excluded from MILP entirely \\
\bottomrule
\end{tabular}
\end{table}

\paragraph{Preprocessing transformation.}
For each HAWB $k$ entering the routing batch:
\begin{itemize}[noitemsep]
  \item \textbf{Fully locked HAWB} (entire route committed; cargo may or may
        not have left origin): \emph{preprocess out of $K$ entirely}. The
        HAWB does not appear in the MILP. Its committed cost is recorded for
        reporting; the MILP has no decision to make about it.
  \item \textbf{Partially locked HAWB} (some upstream arcs executed,
        downstream open): HAWB enters $K$ with \emph{three rewrites}:
        \begin{enumerate}[noitemsep,label=(\arabic*)]
          \item Origin re-pointed to the current observed position (the node
                cargo currently sits at --- e.g.\ $\text{CFS-H}$ at a hub mid-journey).
          \item Subgraph truncated to forward arcs only ($A_k$ excludes any
                arc wholly upstream of the current position).
          \item Initial arrival-time window at the re-pointed origin
                collapsed to $\{t_k^{\text{obs}}\}$ (the observed arrival time
                at the current node); the forward-time-window propagation
                (\S\ref{sec:fwd-time-propagation}) proceeds from there.
        \end{enumerate}
        From there, standard MILP. No special constraints. No $b_k$, no
        buyout decision variable.
\end{itemize}

\paragraph{Lock-break is an orchestrator decision between MILP runs.}
If the rolling-horizon orchestrator decides to break a lock (weighing
buyout cost against predicted improvement), the HAWB is added back to the
\emph{next} routing instance with no lock --- same mechanism as
preprocessing-out, but with the lock cleared first. The MILP itself does not
decide lock breaks.

\paragraph{Pre-MILP reachability check.}
Before the MILP is built, a separate pass evaluates the current locked state
of each $k$ against forward reachability (capacity, cutoffs against the new
deadline, $T_k^{\text{abs}}$). If the locked state has no feasible forward
real-arc path, the HAWB is surfaced as a \emph{structured rescue event} to
the orchestrator (``lock on flight CX880 conflicts with allocation pull;
remediate'') in addition to entering the MILP --- the MILP itself will route
the HAWB via the fallback arc (\S\ref{sec:fallback-arc}) absent intervention,
so the pre-MILP warning is an early signal, not a feasibility gate. The
orchestrator then decides among accepting tardiness (do nothing --- the MILP
runs with the HAWB and C.10 absorbs the slip), breaking the lock (re-add
unlocked), or escalating.

\paragraph{Supply-side lock invalidation.}
Real operations generate supply-side events that invalidate previously-committed
routings:
\begin{itemize}[noitemsep]
  \item \textbf{Flight cancellation} --- carrier cancels (mechanical, weather,
        crew). IATA Cargo-IMP \texttt{FFM} amendment / \texttt{FSU} status.
  \item \textbf{Equipment swap} --- same flight number, different
        $\text{ac}(f)$, potentially different $U_a$ / $N_{a,u}$ on the
        affected arc.
  \item \textbf{Allocation pull / capacity reduction} --- carrier reduces the
        forwarder's contracted allotment.
  \item \textbf{Cargo cutoff shift} --- $\text{CO}_a^*$ recomputed earlier.
\end{itemize}
The supply event set $\mathcal{E}^{\text{sup}}$ is consumed by the
rolling-horizon controller (not the MILP). For each affected HAWB $k$ with
locked arcs, an event $e$ invalidates the lock when:
\begin{enumerate}[noitemsep,label=(\alph*)]
  \item Flight cancellation on a leg of the locked MAWB-arc.
  \item Equipment swap rendering the locked ULD assignment infeasible.
  \item Allocation pull below the committed quantity on the locked arc.
  \item Cutoff shift past the cargo's now-feasible ready time.
\end{enumerate}
On invalidation, the affected arc is removed from the locked prefix,
downstream arcs that depended on it are likewise unlocked, and the HAWB is
enqueued for a rescue re-solve.

\paragraph{Booking-confirmation rejection recovery.}
Carrier rejection (\texttt{NS}, operational refusal) is a workflow
concern, not a MILP constraint. The rejected arc is excluded from the
HAWB's subgraph on the next solve; no model change required.

%%--------------------------------------------------------------------
\section{Service Products and Service-Level Commitments}
\label{sec:service-products}
%%--------------------------------------------------------------------

Forwarders sell named service products to shippers --- bundled commitments
covering transit-time, allowed modes, allowed carriers, equipment, and
handling tier. A shipper buys ``Premium Air Express, 3-day door-to-door'' or
``Multimodal Economy, 25-day''; every routing decision for that shipment must
honor the bundle.

\paragraph{Catalog.}
$P$ is the per-tenant catalog, indexed by $p$. Each $p \in P$:
\begin{table}[H]
\centering
\caption{Service-product catalog fields (per $p \in P$).}
\label{tab:service-product-fields}
\begin{tabular}{lp{9.5cm}}
\toprule
Field & Description \\
\midrule
$\text{id}(p)$ & Stable identifier (e.g., \texttt{PRM\_AIR\_EXP}, \texttt{STD\_AIR}, \texttt{MM\_ECON}) \\
$\text{name}(p)$ & Display name \\
$\text{mode\_allow}(p)$ & Subset of $\{\text{AIR}, \text{OCEAN}, \text{GROUND}\}$ \\
$\text{carrier\_allow}(p)$ & Allowed operating carriers, or \texttt{ANY} \\
$\text{carrier\_deny}(p)$ & Denied operating carriers; deny wins \\
$\text{ac\_type\_allow}(p)$ & Allowed aircraft types (\texttt{FREIGHTER}, \texttt{PAX\_BELLY}, \ldots), or \texttt{ANY} \\
$T^{\text{SLA}}_p$ & Transit-time SLA: maximum door-to-door duration, in hours \\
$w_p$ & Quadratic soft-tardiness weight for service product $p$ (USD/h\textsuperscript{2}); combined with the per-HAWB value coefficient into $W_k = w_p \cdot \mu_k$ in the objective term $W_k \cdot \tau_k^2$ (C.10, \S\ref{sec:lin-tardiness}). \texttt{CALIBRATION NEEDED}. \\
$\text{handling\_tier}(p)$ & \texttt{priority} / \texttt{standard} / \texttt{economy} (consumed by the orchestrator, not the within-batch MILP) \\
$\text{cargo\_type\_min}(p)$ & Minimum cargo-handling certification (e.g., \texttt{COLD\_CHAIN}, \texttt{PHARMA}, \texttt{GEN}) \\
\bottomrule
\end{tabular}
\end{table}

\paragraph{Per-HAWB binding.}
Each $k \in K$ carries $\text{sp}(k) \in P$. At solve time, bundle attributes
pre-filter $A_k$ (\S\ref{sec:prefilter}) and contribute the soft-tardiness
weight to the objective.

\paragraph{Example catalog (illustrative).}
\begin{table}[H]
\centering
\caption{Illustrative service-product catalog (tenant-configurable).}
\label{tab:service-product-catalog}
\begin{tabular}{lllllr}
\toprule
ID & Mode & Carrier & Aircraft & SLA & Handling \\
\midrule
\texttt{PRM\_AIR\_EXP} & AIR & \{CX, BR, CV, LH, EK\} & FREIGHTER & 72h & priority \\
\texttt{STD\_AIR} & AIR & ANY & ANY & 144h & standard \\
\texttt{ECON\_AIR} & AIR & ANY & ANY & 240h & economy \\
\texttt{MM\_EXPEDITED} & AIR, OCEAN & ANY & ANY & 240h & priority \\
\texttt{MM\_STD} & AIR, OCEAN & ANY & ANY & 600h & standard \\
\texttt{MM\_ECON} & AIR, OCEAN & ANY & ANY & 840h & economy \\
\texttt{OCN\_EXP} & OCEAN & \{MSC, CMA, MAERSK\} & --- & 600h & priority \\
\texttt{OCN\_STD} & OCEAN & ANY & --- & 960h & standard \\
\bottomrule
\end{tabular}
\end{table}

\paragraph{Integration with subgraph construction.}
For each leg $f$ of arc $a$:
\begin{align}
  \text{mode\_ok}(k, a) \;&\defeq\; \text{AIR} \in \text{mode\_allow}(\text{sp}(k))
    \quad (\text{for air arcs; trivially true on AIR-allowed product})
    \label{eq:sp-mode-ok} \\
  \text{ac\_type\_ok}(k, a) \;&\defeq\;
    \forall f \in \text{legs}(a):\;
    \text{ac\_type\_allow}(\text{sp}(k)) = \texttt{ANY} \,\vee\, \text{ac}(f) \in \text{ac\_type\_allow}(\text{sp}(k))
    \label{eq:sp-ac-ok}
\end{align}
The carrier predicate is replaced by the policy-resolved version
$\text{carrier\_ok}(k, a)$ from \S\ref{sec:carrier-policy}. An arc is admitted
to $A_k$ only if all three predicates hold (in addition to cutoff, embargo,
lithium, cargo-type, and physical-fit checks).

\paragraph{SLA as a soft constraint via quadratic tardiness penalty.}
The transit-time SLA $T^{\text{SLA}}_p$ enters the model as the \emph{soft
quadratic} tardiness in C.10, with $\Delta_k = \min(T_k^{\text{dead}},
T_k^{\text{SLA}})$. The MILP minimizes $W_k \cdot \tau_k^2$ (PWL outer
approximation, \S\ref{sec:lin-tardiness}) where $\tau_k \geq 0$ measures
lateness against $\Delta_k$ and $W_k = w_p \cdot \mu_k$ combines the
service-product weight with the per-HAWB value coefficient.
\textbf{No hard SLA constraint, and no hard backstop, in v3 of the model.}
Feasibility is guaranteed by the per-HAWB fallback arc
(\S\ref{sec:fallback-arc}); $T_k^{\text{abs}}$ bounds the worst-case
quadratic penalty and the fallback arrival time but is not a feasibility
cutoff. This is the item 3 decision in the review
(\texttt{air\_review\_notes.md}): a hard SLA on a committed shipment can
produce spurious infeasibility; pricing it as a soft quadratic penalty lets
the optimizer elect a paid SLA breach when premium upgrade cost exceeds the
penalty, and prefer \emph{reasonable} lateness over concentrated breach when
slippage is unavoidable.

\paragraph{Finding S Ch 1 framing.}
When $T_k^{\text{SLA}}$ derives from a service product, this is a
\emph{planning bound}, not a contractual guarantee. The clock binds against
deterministic $t_k^{D_k^{\text{node}}}$ in MVP; a later phase will swap for
the Transit Time Service's P85--P90 quantile (hook noted, not implemented
here; see \texttt{transit\_time\_model.md}).

%%--------------------------------------------------------------------
\section{Carrier Policy and Rules Resolution}
\label{sec:carrier-policy}
%%--------------------------------------------------------------------

Beyond the service-product bundle, forwarders maintain a layered set of
carrier-level rules: tenant-wide blacklists (payment disputes, sanctions,
reliability); shipper-lane allow/deny lists (contractual restrictions per
shipper $\times$ lane); commodity-level rules (DGR-certified carriers for
hazmat). All rules are \emph{hard} allow/deny; there is no soft preference
layer. Operational outcomes a tenant might describe as ``preference''
(volume kickers, strategic carrier relationships, BSA allocation balancing)
are correctly modeled as cost --- volume kickers are rate-card extensions
on the underlying offer; strategic relationships should be priced into BSA
negotiations; BSA balancing is already realized by the equalized-allowance
accumulator (\S\ref{sec:bsa-equalized-settlement}). The model rejects
a separate ``preference'' dimension with a dollar-tradeoff knob because no
principled calibration exists for it.

\paragraph{Symbols used in this section.}
\begin{table}[H]
\centering
\caption{Symbols used in the carrier-policy cascade.}
\label{tab:carrier-policy-syms}
\begin{tabular}{lp{9cm}p{3.5cm}}
\toprule
Symbol & Brief & Defined \\
\midrule
$\ell$ & Layer index in the 4-layer cascade (blacklist / shipper-lane / service-product / commodity overlay) & this section (Cascade) \\
$\text{allow}_\ell(k),\, \text{deny}_\ell(k)$ & Per-layer allow / deny sets of carriers for HAWB $k$ & this section \\
$\mathcal{A}_k$ & Intersection of non-wildcard allow sets across layers & Eq.~\ref{eq:carrier-allow} \\
$\mathcal{D}_k$ & Union of deny sets across layers & Eq.~\ref{eq:carrier-allow} \\
$C_k^{\text{allow}}$ & Resolved hard-allowed carriers: $\mathcal{A}_k \setminus \mathcal{D}_k$ (deny wins) & Eq.~\ref{eq:carrier-allow} \\
$C_k^{\text{deny}}$ & Resolved hard-denied carriers: $\mathcal{D}_k$ & this section \\
$\text{carrier}^{\text{op}}(f),\, \text{carrier}^{\text{mk}}(f)$ & Operating vs marketing carrier of flight $f$ (codeshare aware) & this section (op vs mk) \\
$\text{carrier\_ok}(k, a)$ & Resolved per-leg carrier admissibility predicate on arc $a$ & Eq.~\ref{eq:carrier-ok-resolved} \\
$k,\, a,\, f$ & HAWB / arc / flight indices & \S\ref{sec:variables} \\
\bottomrule
\end{tabular}
\end{table}

\paragraph{Cascade and hierarchy.}
Rules are organized in a precedence order, evaluated outside the MILP by a
separate rules-engine component (cf.\ \texttt{build\_plan.md}). The output is
a per-HAWB resolved allow / deny pair. Levels, broadest to narrowest:
\begin{enumerate}[noitemsep]
  \item Tenant carrier blacklist: forwarder-wide hard deny.
  \item Shipper-lane allow/deny: per (shipper, lane).
  \item Service product bundle: $\text{carrier\_allow}(\text{sp}(k))$,
        $\text{carrier\_deny}(\text{sp}(k))$.
  \item Commodity-rule overlay: DGR / cold-chain / pharma-grade carrier
        restrictions.
\end{enumerate}

\paragraph{Conflict semantics: deny wins.}
Let $\mathcal{D}_k = \bigcup_{\ell} \text{deny}_\ell(k)$ and
$\mathcal{A}_k = \bigcap_{\ell : \text{allow}_\ell(k) \neq \texttt{ANY}} \text{allow}_\ell(k)$.
The resolved hard-allowed carrier set:
\begin{equation}
  C_k^{\text{allow}} \;\defeq\; \mathcal{A}_k \setminus \mathcal{D}_k
  \label{eq:carrier-allow}
\end{equation}
An explicit deny at any layer is honored. Rationale: denies represent
intentional safety/compliance/contract decisions; silently overriding them
by a higher-tier allow would be a defect. $C_k^{\text{deny}} = \mathcal{D}_k$
is exposed for audit / belt-and-suspenders gating.

\paragraph{Subgraph integration.}
\begin{equation}
  \text{carrier\_ok}(k, a) \;\defeq\;
    \forall f \in \text{legs}(a):\;
    \bigl(C_k^{\text{allow}} = \texttt{ANY} \vee \text{carrier}^{\text{op}}(f) \in C_k^{\text{allow}}\bigr)
    \wedge \text{carrier}^{\text{op}}(f) \notin C_k^{\text{deny}}
  \label{eq:carrier-ok-resolved}
\end{equation}
The service-product carrier fields feed the cascade as layer 3; they no
longer gate the subgraph directly. The per-leg AND handles interline arcs
correctly.

\paragraph{Operating vs marketing carrier scope.}
Each rule's scope JSONB carries a \texttt{carrier\_basis} field:
\begin{itemize}[noitemsep]
  \item \texttt{op} (default for compliance/regulatory rules): evaluate
        against $\text{carrier}^{\text{op}}(f)$. Used when the rule concerns
        physical aircraft operation, ULD pool ownership, or operational
        safety.
  \item \texttt{mk} (default for shipper-contract rules): evaluate against
        $\text{carrier}^{\text{mk}}(f)$. Used when the shipper's MSA names a
        carrier on the AWB regardless of underlying operation (``no UA on
        our pharma lane'' typically means UA-prefix AWB, even when operated
        by LH/NH under codeshare).
  \item \texttt{either} (when set): rule applies if either op or mk matches.
\end{itemize}
Default migration mapping: tenant blacklist $\to$ \texttt{op}; shipper-lane
$\to$ \texttt{mk}; service-product carrier list $\to$ \texttt{op};
commodity overlay $\to$ \texttt{op}.

\paragraph{Single-pass solve.}
With only hard allow/deny in the cascade, the MILP is solved \emph{once}
with cost as the sole objective (the objective of \S\ref{sec:objective}).
No Pass 2, no $\varepsilon$-constraint, no lexicographic mechanism.
Carrier-policy enforcement happens entirely at the subgraph pre-filter
(\S\ref{sec:prefilter}, predicate 5) via $\text{carrier\_ok}(k, a)$ ---
ineligible arcs never enter the MILP.

\paragraph{Why no soft-preference layer.}
A separate ``preference'' dimension would require a dollar-tradeoff weight
or a cost-ceiling slack that has no principled calibration: tenants
cannot defensibly state ``I would pay \$X more to use carrier CX over
carrier BR'' for any concrete $X$. The legitimate operational behaviors
that ``preference'' might be reaching for --- volume-kicker rebates,
strategic-relationship value, BSA allocation balancing --- are all
correctly modeled as cost: volume kickers extend the rate card on the
underlying offer (\S\ref{sec:supply-option-catalog}); strategic
relationships are priced into BSA negotiations and surface in $r_c$;
BSA balancing falls out of the equalized-allowance accumulator
(\S\ref{sec:bsa-equalized-settlement}). A preference layer in the
optimizer would either double-count these mechanisms (if calibrated to
match them) or introduce un-auditable arbitrary weights (if not).

\paragraph{Volume kickers: monotone vs.\ period-scoped.}
The ``extends the rate card'' claim above holds cleanly only for a
\emph{monotone, per-shipment} kicker --- a lower marginal rate on
chargeable weight above a per-shipment break. That stays a non-decreasing
cost and folds into the existing weight-break card with no new machinery.
It does \emph{not} hold for a \emph{retroactive / all-units} kicker or any
kicker measured over a \emph{period} (monthly / IATA-season) or
\emph{network-wide} scope: there the discount on a shipment depends on
cumulative volume the per-batch MILP cannot see, and the cost-in-CW curve
becomes non-convex (incremental discount) or discontinuous with a downward
jump at the threshold (all-units rebate) --- breaking the convexity /
monotonicity the $\geq$-cut linearizer relies on. \textbf{This is the same
decomposition problem as equalized BSA settlement and is handled the same
way:} an external accumulator carries period-to-date volume per
incentivized contract across batches and feeds the per-batch solve a
\emph{control input} (a remaining-allowance or a per-kg price signal
$\delta_c$ that lowers the rate by a constant when near the threshold),
keeping every batch a clean, monotone, $\geq$-cut-safe MILP --- exactly as
$A_c$ does in \S\ref{sec:bsa-equalized-settlement}. The exact signal
(forecast-driven $\delta_c$ vs.\ Lagrangian target-volume coupling) is
deferred (\S\ref{sec:deferred}, item~\ref{item:volume-kicker-accumulator});
the architectural commitment --- \emph{kicker state lives outside the
per-batch objective, never endogenously} --- is fixed here.

\paragraph{Worked example.}
Tenant \texttt{ACME\_FWD}. HAWB $k$ from \texttt{Beta Corp} on TPE$\to$JFK,
service product \texttt{PRM\_AIR\_EXP} ($\text{carrier\_allow} = \{\text{CX, BR,
CV, LH, EK, KE}\}$).
\begin{table}[H]
\centering
\caption{Carrier-policy cascade worked example: tenant \texttt{ACME\_FWD}, HAWB on TPE$\to$JFK under \texttt{PRM\_AIR\_EXP}.}
\label{tab:carrier-policy-worked-example}
\begin{tabular}{p{3.2cm}p{4.5cm}p{6.5cm}}
\toprule
Layer & Rule contribution & Running state \\
\midrule
Tenant blacklist & deny \{KE\} & allow $= \{$CX, BR, CV, LH, EK$\}$ \\
Shipper-lane (\texttt{Beta Corp}, TPE$\to$JFK) & deny \{AA\} & AA was not in allow; tracked as deny for audit \\
Service product \texttt{PRM\_AIR\_EXP} & allow \{CX, BR, CV, LH, EK, KE\} & layer 3 intersection unchanged \\
Commodity overlay (no rule) & --- & no change \\
\bottomrule
\end{tabular}
\end{table}
Resolved: $C_k^{\text{allow}} = \{\text{CX, BR, CV, LH, EK}\}$,
$C_k^{\text{deny}} = \{\text{KE, AA}\}$. The MILP optimizes cost over
arcs whose every leg's $\text{carrier}^{\text{op}}(f) \in C_k^{\text{allow}}$
and $\notin C_k^{\text{deny}}$.

\paragraph{Where the rules engine lives.}
A separate component exposing
\texttt{resolve\_carrier\_policy(tenant\_id, shipment) $\to$
$(C^{\text{allow}}, C^{\text{deny}})$}. Its own LaTeX model
is in scope for Phase 1.

\paragraph{Storage and versioning.}
\texttt{policy\_rules} / \texttt{policy\_snapshots} /
\texttt{routing\_run\_policy\_bindings} schema
(\texttt{data\_model.md \S4}). Every solve binds to a snapshot at run start;
rule edits between snapshot and solve completion do not affect the in-flight
solve. Soft-delete preserves rows for audit.

\paragraph{Out of MVP scope.}
Time-windowed rules, conditional rules (``require freighter when commodity
is electronics''), and ML-learned allow/deny refinements from override
history --- all deferred to the rules-engine model and Phase 5
constraint learning. \emph{Note:} all such future extensions remain hard
allow/deny; no soft-preference dimension is reintroduced.

%%--------------------------------------------------------------------
\section{Per-shipment subgraph pre-filter}
\label{sec:prefilter}
%%--------------------------------------------------------------------

For each HAWB $k$, the per-shipment subgraph $A_k \subseteq A$ is constructed
by pruning arcs that fail any of eight predicates. This is Phase~1 step~11 of
the graph-construction logic (\texttt{air\_graph\_construction.md \S5}),
repeated here for completeness because it determines what the MILP sees.

\paragraph{Symbols used in this section.}
\begin{table}[H]
\centering
\caption{Symbols used in the per-shipment subgraph pre-filter.}
\label{tab:prefilter-syms}
\begin{tabular}{lp{9cm}p{3.5cm}}
\toprule
Symbol & Brief & Defined \\
\midrule
$A_k \subseteq A$ & HAWB $k$'s per-shipment subgraph (output of this section) & \S\ref{sec:variables} \\
$O_k,\, D_k^{\text{node}}$ & Origin / destination doors of $k$ & \S\ref{sec:hawb-params} \\
$\text{cargo\_type\_ok}(k, a)$ & Cargo-class equipment compatibility on every leg of $a$ & Eq.~\ref{eq:cargo-type-ok} \\
$\text{embargo\_ok}(k, a)$ & No active embargo matches $k$ on any leg of $a$ & Eq.~\ref{eq:embargo-feasibility} \\
$\text{lithium\_ok}(k, a)$ & Lithium PI / Section / aircraft-type acceptance on every leg & Eq.~\ref{eq:lithium-ok} \\
$\text{mode\_ok}(k, a)$ & Service-product mode-allow check (AIR for air arcs) & Eq.~\ref{eq:sp-mode-ok} \\
$\text{ac\_type\_ok}(k, a)$ & Service-product aircraft-type-allow check per leg & Eq.~\ref{eq:sp-ac-ok} \\
$\text{carrier\_ok}(k, a)$ & Carrier-policy cascade admissibility (Layer 3 in pre-filter) & Eq.~\ref{eq:carrier-ok-resolved} \\
$t^{\text{lo}}_n,\, t^{\text{hi}}_n$ & Propagated arrival-time window at node $n$ for HAWB $k$ & \S\ref{sec:fwd-time-propagation} \\
$\text{CO}_a^*,\, \mu_a,\, \text{ETA}_a$ & Per-arc effective cutoff / transit / scheduled arrival scalars & \S\ref{sec:air-arc-params} \\
$T_k^{\text{abs}}$ & Latest valid arrival time (fallback-arc arrival; predicate 6 reachability bound) & \S\ref{sec:hawb-params} \\
$a^{\text{fb}}_k$ & Fallback arc of $k$ --- exempt from this entire pre-filter & \S\ref{sec:fallback-arc} \\
$w_k,\, v_k$ & Actual weight / volume of $k$ (predicate 8 ULD physical fit) & \S\ref{sec:hawb-params} \\
$W_u,\, V_u$ & Per-ULD payload / volume caps (predicate 8) & \S\ref{sec:uld-spec} \\
$U_a$ & Admissible ULD types on \texttt{per\_uld\_pivot} arc $a$ & \S\ref{sec:variables} \\
$A^{\text{pu}}$ & Per-ULD-pivot arcs (predicate 8 scope) & \S\ref{sec:variables} \\
\bottomrule
\end{tabular}
\end{table}

\begin{enumerate}[noitemsep]
  \item \textbf{Lane / reachability.} Arc must lie on a path from $O_k$ to
        $D_k^{\text{node}}$ (verified by forward BFS from $O_k$ and backward
        BFS from $D_k^{\text{node}}$, retaining only arcs reachable from both
        ends).
  \item \textbf{Cargo type compatibility.} $\text{cargo\_type\_ok}(k, a)$ ---
        e.g.\ AVI requires AVI-capable equipment on every leg; DGR requires
        DGR-accepting carrier; PER requires cold-chain. Defined per leg:
        \begin{equation}
          \text{cargo\_type\_ok}(k, a) \;\defeq\;
            \forall f \in \text{legs}(a):
            \begin{cases}
              \text{True} & \ctype_k = \text{GEN} \\
              \text{dgr}(f) = 1 & \ctype_k = \text{DGR} \\
              \text{per}(f) = 1 & \ctype_k = \text{PER} \\
              \text{val\_capable}(f) = 1 & \ctype_k = \text{VAL} \\
              \text{avi\_capable}(f) = 1 & \ctype_k = \text{AVI} \\
              \text{hum\_capable}(f) = 1 & \ctype_k = \text{HUM}
            \end{cases}
          \label{eq:cargo-type-ok}
        \end{equation}
  \item \textbf{Embargo.} $\text{embargo\_ok}(k, a)$ (Eq.~\ref{eq:embargo-feasibility}).
  \item \textbf{Lithium.} $\text{lithium\_ok}(k, a)$ (Eq.~\ref{eq:lithium-ok}).
  \item \textbf{Service product (mode / carrier / aircraft type).}
        $\text{mode\_ok} \wedge \text{carrier\_ok} \wedge \text{ac\_type\_ok}$
        (Eqs.~\ref{eq:sp-mode-ok}--\ref{eq:sp-ac-ok},
        \ref{eq:carrier-ok-resolved}). Resolved per-HAWB by the carrier-policy
        cascade (\S\ref{sec:carrier-policy}).
  \item \textbf{Forward-time-window admission.} For each candidate arc $a =
        (n_1, n_2)$, run the propagation rules of
        \S\ref{sec:fwd-time-propagation}: (i) if $a$ is a POL-leaving air arc
        with cutoff $\text{CO}_a^*$ at $\text{tail}(a) = n_1$, admit $a$ only if
        $t^{\text{lo}}_{n_1} \leq \text{CO}_a^*$ (one-arc model: the cutoff sits
        on the air arc's tail, with no added transit; any positioning move is a
        preceding ground arc); (ii) if $a$ is an air arc, set its head window
        $t^{\text{lo}}_{n_2} = t^{\text{hi}}_{n_2} = \text{ETA}_a$;
        (iii) otherwise propagate the window forward. Additionally retain $a$
        only if a forward path from $n_2$ to $D_k^{\text{node}}$ exists with
        cumulative arrival $\leq T_k^{\text{abs}}$ (otherwise $a$ is dominated
        by the fallback on both cost and arrival). The fallback arc is exempt
        from this and every other pre-filter predicate.
  \item \textbf{Per-HAWB ULD physical fit.} For any $a \in A^{\text{pu}}$
        (\texttt{per\_uld\_pivot} arc), prune $a$ from $A_k$ if
        $w_k > \max_{u \in U_a} W_u$ or $v_k > \max_{u \in U_a} V_u$ ---
        i.e.\ no single ULD type admissible on $a$ physically holds the HAWB.
        This catches the cargo that the per-MAWB-aggregate C.5b cannot reject
        (see the C.5b remark). The HAWB may still route via
        \texttt{flat\_rate} / \texttt{min\_flat\_breaks} / \texttt{coload\_per\_kg}
        arcs where per-ULD claiming does not apply.
\end{enumerate}

\paragraph{Per-arc rejection logging --- root-cause attribution.}
The pre-filter records, for every arc $a$ that an HAWB's candidate set
considered but excluded, the \emph{first predicate} that rejected it.
For each $(k, a)$ rejection, the generator emits a structured record:
\begin{equation*}
  \text{reject\_record}(k, a) \;\defeq\; \bigl(k,\; a,\;
    \text{predicate}_{1\ldots 8},\; \text{evidence}\bigr)
\end{equation*}
where $\text{predicate}_{1\ldots 8}$ is the failing pre-filter predicate
index (1 = lane / reachability, 2 = cargo type, 3 = embargo, 4 =
lithium, 5 = service product, 6 = forward-time-window admission, 7 =
destination reachability, 8 = per-HAWB ULD fit) and $\text{evidence}$
is the structured trace of the failing check (e.g., for predicate 6:
the propagated lower bound at the cutoff node, the cutoff scalar, the
arc id; for predicate 4: the lithium PI and ac\_type that failed
acceptance). The aggregate log is emitted as
\texttt{prefilter\_rejections.jsonl} alongside other walking-skeleton
metrics (\S\ref{sec:walking-skeleton-instrumentation}).

\paragraph{Empty real-arc subgraph.}
If the real-arc set surviving pre-filter is empty ($A_k \setminus
\{a^{\text{fb}}_k\} = \emptyset$), the graph generator logs a structured
pre-solve warning (``HAWB $k$: zero real arcs survived pre-filter; routing
via fallback expected''). The HAWB still enters the MILP with $A_k =
\{a^{\text{fb}}_k\}$ (the fallback arc is always present per
\S\ref{sec:fallback-arc}); the optimal solution will set $x_{k,
a^{\text{fb}}_k} = 1$, contributing $C^{\text{fallback}}$ to cost.

Post-solve the orchestrator inspects the fallback set
(\S\ref{sec:output-diagnostics}) and the per-HAWB rejection records to
identify which predicate ate the candidate arcs --- the dominant rejection
predicate per HAWB is reported alongside $K^{\text{fb}}$ as
$\text{predicate}^{\text{dom}}_k$. The orchestrator then decides between
renegotiating constraints (e.g.\ relax service product), upgrading carrier
policy, procuring additional capacity, or escalating to the shipper.
There is no infeasibility branch.

%%--------------------------------------------------------------------
\section{Constraints}
\label{sec:constraints}
%%--------------------------------------------------------------------

Active constraint families: C.1, C.2, C.4, C.5/C.5b/C.5c, C.10, C.13,
C.14. (C.12 is retained below as a placeholder annotation --- locks are
handled by preprocessing outside the MILP, see
\S\ref{sec:locked-commitments}.) Time propagation, cargo cutoffs, and
hub MCT are enforced at graph build via forward-time-window propagation
(\S\ref{sec:fwd-time-propagation}), not in the MILP.

\paragraph{Symbols used in this section.}
Recap of the symbols that appear in C.1--C.14, grouped by role; full
definitions live in the cross-referenced homes.
\begin{table}[H]
\centering
\caption{Symbols used in the constraints section (recap of \S\ref{sec:variables}).}
\label{tab:constraints-syms}
\begin{tabular}{lp{9cm}p{3.5cm}}
\toprule
Symbol & Brief & Defined \\
\midrule
\multicolumn{3}{l}{\textit{Decision variables}} \\
$x_{k,a}$ & HAWB $k$ uses arc $a$ (binary) & \S\ref{sec:variables} \\
$z_{a,g}$ & MAWB $(a, g)$ instantiated (binary) & \S\ref{sec:variables} \\
$\text{CW}_{a,g}$ & MAWB chargeable weight & \S\ref{sec:cw-density-mixing} \\
$\text{Wt}_{a,g},\, \text{Wv}_{a,g}$ & Actual / volumetric weight aggregates on MAWB & \S\ref{sec:cw-density-mixing} \\
$\gamma_{a,g,b}$ & Break selector for \texttt{min\_flat\_breaks} (binary) & \S\ref{sec:variables} \\
$\text{BW}_{a,g,b}$ & Disaggregated billed weight at break $b$ & \S\ref{sec:variables} \\
$\eta_{a,g,u}$ & ULDs of type $u$ claimed on MAWB (integer) & \S\ref{sec:variables} \\
$\text{pivot}_{a,g}$ & Pivot-active billed weight (per-flight BSA) & \S\ref{sec:variables} \\
$\text{over}_c$ & Allowance overage on equalized BSA contract & \S\ref{sec:variables} \\
$t_k^{D_k^{\text{node}}}$ & Arrival time of $k$ at destination (C.10a) & \S\ref{sec:variables} \\
$\tau_k$ & Tardiness of $k$ against $\Delta_k$ & \S\ref{sec:variables} \\
\midrule
\multicolumn{3}{l}{\textit{Sets and indices}} \\
$K,\, N,\, A$ & HAWBs, nodes, arcs & \S\ref{sec:variables} \\
$A_k,\, N_k,\, K_a$ & Pre-filtered subgraph / nodes / HAWBs-on-arc & \S\ref{sec:variables} \\
$M$ & MAWB candidates $\{(a,g)\}$ & \S\ref{sec:variables} \\
$A^{\text{MFB}},\, A^{\text{pu}},\, C^{\text{eq}}$ & \texttt{min\_flat\_breaks} arcs / per-ULD-pivot arcs / equalized BSA contracts & \S\ref{sec:variables} \\
$B_a,\, U_a$ & Weight breaks on arc $a$; admissible ULD types on arc $a$ & \S\ref{sec:variables} \\
$g(k)$ & Consolidation group of $k$ & \S\ref{sec:g-of-k} \\
$\text{tail}(a),\, \text{head}(a)$ & Arc endpoints & \S\ref{sec:variables} \\
\midrule
\multicolumn{3}{l}{\textit{Parameters}} \\
$w_k,\, v_k,\, cw_k$ & Actual weight / volume / chargeable weight of $k$ & \S\ref{sec:hawb-params} \\
$O_k,\, D_k^{\text{node}}$ & Origin / destination doors of $k$ & \S\ref{sec:hawb-params} \\
$T_k^{\text{abs}},\, \Delta_k,\, T_k^{\text{SLA}}$ & Latest valid arrival (mandatory-finite, fallback-arc time) / effective soft deadline / SLA & \S\ref{sec:hawb-params}, \S\ref{sec:service-products}, \S\ref{sec:fallback-arc} \\
$t_k^{\text{rdy,early}},\, t_k^{\text{rdy,late}}$ & Cargo-ready window & \S\ref{sec:hawb-params} \\
$w_{\text{sp}(k)}$ & Service-product tardiness weight (USD/h\textsuperscript{2}) & \S\ref{sec:service-products} \\
$\text{value}_k,\, V^{\text{ref}},\, \mu_k,\, W_k$ & Declared shipment value, tenant reference value, value coefficient, combined tardiness weight & \S\ref{sec:hawb-params} \\
$C^{\text{fallback}}$ & Per-HAWB fallback-arc dollar cost (tenant-global) & \S\ref{sec:hawb-params}, \S\ref{sec:fallback-arc} \\
$\mu_a,\, \text{CO}_a^*,\, \text{ETA}_a$ & Air-arc transit / effective cutoff / scheduled ETA scalars (graph-build only) & \S\ref{sec:air-arc-params} \\
$\text{arr\_dest}(k, a)$ & Precomputed destination-arrival scalar on terminal arcs (C.10a) & \S\ref{sec:graph-construction} \\
$\delta_a,\, \delta^{\text{cust}}_k$ & Ground-arc dwell / per-HAWB customs dwell & \S\ref{sec:ground-arc-params}, \S\ref{sec:hawb-params} \\
$\varepsilon$ & Dunnage / strapping factor ($\approx 0.05$) & \S\ref{sec:air-arc-params} \\
$N_{a,u},\, W_u,\, V_u,\, \text{cap}_a$ & Per-contract ULD allotment / per-ULD physical cap / per-offer cap & \S\ref{sec:air-arc-params} \\
$\pi_a,\, r_c,\, A_c$ & BSA pivot weight / contract rate / allowance & \S\ref{sec:air-arc-params} \\
$\text{rate}_{a,b},\, \text{break}_{a,b}$ & Break-down rate and weight at break $b$ & \S\ref{sec:air-arc-params} \\
$\text{transit}(k, a)$ & Per-HAWB transit on arc & \S\ref{sec:variables} \\
$\text{CW}^{\text{ub}}_{a,g}$ & Per-MAWB chargeable-weight upper bound (big-M) & \S\ref{sec:variables} \\
\bottomrule
\end{tabular}
\end{table}

\subsection*{C.1 Flow conservation}

For each $k \in K$ and each node $n \in N_k$:
\begin{equation}
  \sum_{a \in A_k : \text{tail}(a) = n} x_{k,a}
  \;-\;
  \sum_{a \in A_k : \text{head}(a) = n} x_{k,a}
  \;=\;
  \begin{cases}
    +1 & n = O_k \\
    -1 & n = D_k^{\text{node}} \\
    0  & \text{otherwise}
  \end{cases}
  \label{eq:C1}
\end{equation}
Standard MCNF supply form (outflow $-$ inflow $= b(n)$ with $b(O_k) = +1$,
$b(D_k^{\text{node}}) = -1$, $b(n) = 0$ elsewhere). HAWB enters at origin
door, leaves at destination door, conserves at every intermediate node.

\subsection*{C.2 MAWB instantiation linkage}

For each MAWB $(a, g) \in M$ and each HAWB $k \in K_a$ with $g(k) = g$:
\begin{equation}
  x_{k,a} \;\leq\; z_{a,g}
  \tag{C.2a} \label{eq:C2a}
\end{equation}
For each MAWB $(a, g) \in M$:
\begin{equation}
  z_{a,g} \;\leq\; \sum_{k \in K_a : g(k) = g} x_{k,a}
  \tag{C.2b} \label{eq:C2b}
\end{equation}
Together: $z_{a,g} = 1 \iff \exists k \in K_a,\, g(k) = g,\, x_{k,a} = 1$.
C.2a forces MAWB activation when any HAWB rides the arc in that group; C.2b
prevents phantom activation.

\subsection*{C.4 Chargeable-weight aggregation (density mixing)}

For each MAWB $(a, g) \in M$:
\begin{align}
  \text{Wt}_{a,g} &= (1 + \varepsilon) \sum_{k \in K_a : g(k) = g} w_k \cdot x_{k,a}
    \tag{C.4a} \label{eq:C4a} \\
  \text{Wv}_{a,g} &= \sum_{k \in K_a : g(k) = g} v_k \cdot 167 \cdot x_{k,a}
    \tag{C.4b} \label{eq:C4b} \\
  \text{CW}_{a,g} &\geq \text{Wt}_{a,g}
    \tag{C.4c} \label{eq:C4c} \\
  \text{CW}_{a,g} &\geq \text{Wv}_{a,g}
    \tag{C.4d} \label{eq:C4d}
\end{align}
C.4c--d pin $\text{CW}_{a,g} = \max(\text{Wt}_{a,g}, \text{Wv}_{a,g})$ at
optimum by the monotonicity invariant (\S\ref{sec:objective}). For co-load
arcs no MAWB is instantiated; cost is per-HAWB on $cw_k$ in the objective.

\subsection*{C.5 Per-contract allotment cap (\texttt{per\_uld\_pivot})}

For each $a \in A^{\text{pu}}$ and each ULD type $u \in U_a$:
\begin{equation}
  \sum_{g \in G_a} \eta_{a,g,u} \;\leq\; N_{a,u}
  \label{eq:C5}
\end{equation}
For each MAWB $(a, g) \in M$ with $a \in A^{\text{pu}}$ and each $u \in U_a$
(per-MAWB activation tightening):
\begin{equation}
  \eta_{a,g,u} \;\leq\; N_{a,u} \cdot z_{a,g}
  \tag{C.5-act} \label{eq:C5-act}
\end{equation}
$N_{a,u}$ is the contracted allotment of type-$u$ ULD positions the forwarder
holds on arc $a$ (BSA contract data). C.5 sums across consolidation groups
--- multiple MAWBs sharing the same BSA contract share the allotment. C.5-act
ties ULD claims to MAWB activation: at LP relaxation, $z_{a,g} = 0$ now
forces $\eta_{a,g,u} = 0$ (no phantom ULD claims on inactive MAWBs);
integer-feasible solutions are unchanged but LP relaxation is tighter.

\paragraph{No flight-level physical-capacity constraint.}
The forwarder does not know $W_f$ and cannot plan against it. The constraints
that matter --- and that the forwarder \emph{does} know --- are per-contract
allotment (C.5), per-ULD physical capacity (C.5b), and per-offer caps where
the offer specifies one (C.5c). TACT/NAC offers without a cap are
request/confirm at booking, not a planning-time constraint.

\subsection*{C.5b Per-ULD physical capacity (item 13-A fix)}

For each MAWB $(a, g)$ with $a \in A^{\text{pu}}$:
\begin{align}
  \sum_{k \in K_a : g(k) = g} w_k \cdot x_{k,a}
    &\;\leq\; \sum_{u \in U_a} W_u \cdot \eta_{a,g,u}
    \tag{C.5b-w} \label{eq:C5b-w} \\
  \sum_{k \in K_a : g(k) = g} v_k \cdot x_{k,a}
    &\;\leq\; \sum_{u \in U_a} V_u \cdot \eta_{a,g,u}
    \tag{C.5b-v} \label{eq:C5b-v}
\end{align}
\textbf{Uses $w_k$ (actual mass), not $cw_k$ (chargeable).} $W_u$ is a
physical payload limit; the prior model used $cw_k$ here, double-counting
light-bulky cargo against both volume (C.5b-v) and weight. Worked LD3 case:
$W_u \approx 1588$ kg, $V_u \approx 4.5$ m\textsuperscript{3}; dense 1400 kg
/ 1 m\textsuperscript{3} HAWB plus light 30 kg / 3 m\textsuperscript{3} HAWB
$\to$ fits physically (actual 1430 kg, volume 4 m\textsuperscript{3}), but
the old-form $cw$ sum $1400 + 501 = 1901 > 1588$ wrongly rejects. C.5b-v is
independent --- light/bulky cargo (e.g.\ apparel $\sim 80$ kg/m\textsuperscript{3})
binds on volume long before weight; a 400 kg apparel HAWB at 5
m\textsuperscript{3} needs 2 LD3 by volume even though weight says 1.

\paragraph{Accepted looseness.}
C.5b is a per-MAWB aggregate bound on the total weight / volume against the
sum of ULD-type capacities. It does \emph{not} enforce that each individual
HAWB fits inside one ULD --- e.g.\ a single 3000 kg HAWB on $\eta = 2$ LD3
satisfies $3000 \leq 2 \cdot 1588 = 3176$ even though no LD3 individually
holds it. True bin-packing (HAWB-to-ULD-instance assignment) is deferred
(\S\ref{sec:deferred}). Pre-filter step 8 (\S\ref{sec:prefilter}) catches the
worst per-HAWB case (any single HAWB exceeding $\max_u W_u$ or $\max_u V_u$
on a \texttt{per\_uld\_pivot} arc), which is the failure mode worth blocking
at MVP.

\subsection*{C.5c Per-offer cap (when specified)}

For each arc $a$ with a specified offer-level cap (BSA secondary cap, GSA
tier limit, co-loader quote):
\begin{equation}
  \sum_{k \in K_a} w_k \cdot x_{k,a} \;\leq\; \text{cap}_a
  \label{eq:C5c}
\end{equation}
(Use $\text{cap}_a^{\text{cl}}$ for co-load arcs.) Offers without a specified
cap skip C.5c.

\paragraph{BSA arcs --- $\text{cap}_a$ is a dynamic period-pacing control input.}
For BSA arcs, $\text{cap}_a$ is \emph{not} static offer data: it is supplied
per solve by the external BSA pacing model (the same accumulator/orchestrator
that supplies the equalized allowance $A_c$, \S\ref{sec:bsa-equalized-settlement}),
which sees the whole period and the period-to-date consumption the myopic
per-batch MILP cannot. This is where the period-awareness lives --- not as an
in-MILP period ceiling (\S\ref{sec:objective}, C.13a ``No period overage
cap''). The controller steers usage by period pacing:
\begin{itemize}[noitemsep]
  \item \textbf{Running behind commitment} (under-consuming $P_{c,t}$, period
        early): relax --- set $\text{cap}_a$ non-binding so the firm allotment
        is fully available, and the $A_c$ free-allowance already pulls cargo
        onto the BSA (sunk capacity should not be wasted to spot).
  \item \textbf{Running hot / near period end} (over-consuming, or limited
        flights left): tighten --- set $\text{cap}_a$ \emph{below} the firm
        allotment to throttle BSA loading and reserve / divert to spot,
        e.g.\ to hold capacity for forecast high-priority demand or avoid
        burning overage on low-value cargo.
\end{itemize}
\textbf{Scope of the lever.} $\text{cap}_a$ (a weight cap) can only throttle
\emph{at or below} the firm per-flight allotment, because C.5 ($N_{a,u}$) is a
separate hard physical cap that still binds --- $\min(\text{cap}_a,\ \text{C.5
physical capacity})$ is operative. Authorizing usage \emph{above} the firm
allotment (requesting free-sale space when behind) is therefore an adjustment
to $N_{a,u}$ (C.5), not to $\text{cap}_a$; the controller may set both. Either
way these are exogenous control inputs --- the MILP body is unchanged, it
simply reads the per-batch $\text{cap}_a$ and $N_{a,u}$ values.

\subsection*{C.10 Destination arrival, soft deadline, quadratic tardiness}

\paragraph{Destination arrival (C.10a).}
HAWB $k$'s arrival at its destination node is a precomputed function of the
chosen terminal arc. Let
\[
  \mathcal{A}^{\text{last}}_k \;\defeq\; \{a \in A_k : \text{head}(a) =
    \text{POD}_k\} \;\cup\; \{a^{\text{fb}}_k\}
\]
be the set of arcs that terminate $k$'s routing --- air arcs landing at the
HAWB's port of discharge, plus the fallback arc. Flow conservation
guarantees exactly one $a \in \mathcal{A}^{\text{last}}_k$ has $x_{k,a} = 1$
in any feasible solution. The continuous variable $t_k^{D_k^{\text{node}}}
\geq 0$ is pinned by
\begin{equation}
  t_k^{D_k^{\text{node}}} \;=\;
    \sum_{a \in \mathcal{A}^{\text{last}}_k} \text{arr\_dest}(k, a) \cdot x_{k,a}
  \label{eq:C10a}
\end{equation}
where $\text{arr\_dest}(k, a)$ is a graph-build-precomputed scalar:
\begin{itemize}[noitemsep]
  \item air arc with $\text{head}(a) = \text{POD}_k$: $\text{arr\_dest}(k, a)
        = \text{ETA}_a + \Delta^{\text{post}}_k$, where $\Delta^{\text{post}}_k$
        is the sum of post-POD ground-arc transits (cartage, CFS dwell,
        customs, final delivery) on $k$'s destination tail;
  \item fallback arc: $\text{arr\_dest}(k, a^{\text{fb}}_k) = T_k^{\text{abs}}$.
\end{itemize}
Intermediate-node arrival times do not appear in the MILP --- their
feasibility (cargo cutoffs at POLs, MCT at hubs) is enforced at graph build
by forward-time-window propagation (\S\ref{sec:fwd-time-propagation}).

\paragraph{Soft deadline and tardiness (C.10b).}
For each $k \in K$:
\begin{equation}
  \tau_k \;\geq\; t_k^{D_k^{\text{node}}} - \Delta_k
  \label{eq:C10b}
\end{equation}
with $\tau_k \geq 0$ (domain, C.14) and $\Delta_k = \min(T_k^{\text{dead}},
T_k^{\text{SLA}})$ (\S\ref{sec:hawb-params}). Tardiness enters the objective
\emph{quadratically} via the auxiliary $\text{pen}_k \geq W_k \cdot \tau_k^2$,
where $W_k = w_{\text{sp}(k)} \cdot \mu_k$ combines the service-product
tardiness weight and the value coefficient $\mu_k = \text{value}_k /
V^{\text{ref}}$. The objective contributes $+ \sum_k \text{pen}_k$.

\paragraph{Why quadratic.}
Linear penalty is indifferent between (one HAWB late by $2T$ hours) and (two
HAWBs late by $T$ hours each) at identical $W$. Quadratic prefers the latter
($2 W T^2 < W (2T)^2 = 4 W T^2$), pushing the optimizer toward
\emph{reasonable} lateness when slippage is unavoidable rather than
concentrating the slip on one shipment. Value scaling via $\mu_k$ ensures a
high-value HAWB receives proportionally larger marginal pressure.

\paragraph{Linearization.}
$\tau_k^2$ is convex, so the quadratic is represented by an outer
approximation: $\text{pen}_k$ is constrained from below by a family of
linear cuts at tangent breakpoints $\hat\tau_{k,j} = \alpha_j (T_k^{\text{abs}}
- \Delta_k)$ placed \emph{per HAWB}, relative to that HAWB's feasible
tardiness range. Detail in \S\ref{sec:lin-tardiness}. Default
$\alpha = \{0, 0.25, 0.5, 0.75, 1.0\}$ (5 cuts per HAWB); the top breakpoint
$\alpha = 1.0$ coincides with the fallback-arc tardiness, so no extrapolation
occurs on wide-backstop products.

\paragraph{Finding S Ch~1 framing.}
When $T_k^{\text{SLA}}$ derives from a service product, this is a planning
bound, not a contractual guarantee. MVP binds against deterministic
$\text{arr\_dest}(k, a)$ scalars; later phase will swap for the TT-Service
P85--P90 end-to-end quantile precomputed per terminal arc.

\subsection*{C.12 Locked commitments --- preprocessing, not constraints}

Locks are resolved at preprocessing (\S\ref{sec:locked-commitments}); the
MILP body is lock-agnostic. The locked-state pre-MILP reachability check is
an early warning, not a feasibility gate --- if no forward path exists, the
MILP will route the HAWB via the fallback arc absent orchestrator
intervention.

\subsection*{C.13 BSA settlement}

\subsubsection*{C.13a Equalized-settlement allowance}

For each equalized-settlement contract $c \in C^{\text{eq}}$ with allowance $A_c$:
\begin{align}
  \text{over}_c &\;\geq\; \text{chargeable}(c) - A_c
    \tag{C.13a} \label{eq:C13a} \\
  \text{over}_c &\;\geq\; 0 \quad \text{(domain, C.14)}
\end{align}
where
\[
  \text{chargeable}(c) \;\defeq\; \sum_{(a, g) :\, a \in A_c^{\text{MAWB}},\, g \in G_a} \text{CW}_{a,g}.
\]
Cost contribution = $r_c \cdot \text{over}_c$ (free up to $A_c$, then $r_c$/kg).
The sunk portion $r_c \cdot P_{c,t}$ is constant (booked outside the MILP);
dropped from the optimizer's argmin, surfaced in cost reporting separately.

\paragraph{No period overage cap --- the per-flight allotment is the limit.}
Weight above $P_{c,t}$ is \emph{priced} ($r_c \cdot \text{over}_c$), not
capped: $P_{c,t}$ is the take-or-pay basis (sunk), and the forwarder is free
to use --- and pay overage on --- the ULD positions it actually holds. The
operative usage limit on a BSA is therefore the \emph{per-flight allotment}
$N_{a,u}$ enforced by C.5 ($\sum_g \eta_{a,g,u} \leq N_{a,u}$ per contracted
arc $a$, ULD type $u$), plus per-ULD physical capacity C.5b; these already
bound $\text{over}_c$. A period weight ceiling of the form $\alpha_c P_{c,t}$
would wrongly forbid using contracted positions and is \emph{not} imposed.

\subsubsection*{C.13b Per-flight settlement pivot}

For each MAWB $(a, g) \in M$ with $a \in A^{\text{pu}}$ and $\text{settlement}_a
= \texttt{per\_flight}$:
\begin{align}
  \text{pivot}_{a,g} &\;\geq\; \text{CW}_{a,g}
    \tag{C.13b-1} \label{eq:C13b-1} \\
  \text{pivot}_{a,g} &\;\geq\; \pi_a \cdot \sum_{u \in U_a} \eta_{a,g,u}
    \tag{C.13b-2} \label{eq:C13b-2}
\end{align}
Cost contribution = $r_a \cdot \text{pivot}_{a,g}$. Since $r_a > 0$ the LP
relaxation pins $\text{pivot}_{a,g} = \max(\text{CW}_{a,g},\, \pi_a \sum_u
\eta_{a,g,u})$.

\subsection*{C.14 Domain (with explicit upper bounds where they tighten the LP)}
\label{sec:domain}

\begin{align*}
  x_{k,a}      &\in \{0, 1\}                              & & \forall k \in K,\, a \in A_k \\
  z_{a,g}      &\in \{0, 1\}                              & & \forall (a, g) \in M \\
  \gamma_{a,g,b} &\in \{0, 1\}                            & & \forall (a, g) \in M,\, a \in A^{\text{MFB}},\, b \in B_a \\
  \eta_{a,g,u} &\in \mathbb{Z} \cap [0,\, N_{a,u}]        & & \forall (a, g) \in M,\, a \in A^{\text{pu}},\, u \in U_a \\
               & & & \text{(per-MAWB cap --- tighter than aggregate C.5)} \\
  \text{CW}_{a,g} &\in [0,\, \text{CW}^{\text{ub}}_{a,g} \cdot z_{a,g}] & & \forall (a, g) \in M \\
               & & & \text{(empty bucket} \Rightarrow \text{CW} = 0) \\
  \text{Wt}_{a,g},\,\text{Wv}_{a,g} &\in [0,\, \text{CW}^{\text{ub}}_{a,g}] & & \\
  \text{BW}_{a,g,b} &\in [0,\, \max(\text{CW}^{\text{ub}}_{a,g},\, \max_{b' \in B_a} \text{break}_{a,b'})] & & \forall b \in B_a \\
                & & & \text{(widened bound; see \S\ref{sec:bigm-tightening})} \\
  \text{pivot}_{a,g} &\in [0,\, \max(\text{CW}^{\text{ub}}_{a,g},\, \pi_a \sum_u N_{a,u})] & & \forall (a, g) \in M,\, a \in A^{\text{pu}} \\
  \text{over}_c &\in [0,\, \sum_{(a,g) :\, a \in A_c^{\text{MAWB}}} \text{CW}^{\text{ub}}_{a,g}] & & \forall c \in C^{\text{eq}} \\
  \tau_k        &\in [0,\, \max(0,\, T_k^{\text{abs}} - \Delta_k)] & & \forall k \in K \\
  \text{pen}_k  &\in [0,\, W_k \cdot (T_k^{\text{abs}} - \Delta_k)^2] & & \forall k \in K \\
  t_k^{D_k^{\text{node}}} &\in \Bigl[\min_{a \in \mathcal{A}^{\text{last}}_k} \text{arr\_dest}(k, a),\, T_k^{\text{abs}}\Bigr] & & \forall k \in K
\end{align*}
The $t_k^{D_k^{\text{node}}}$ upper bound $T_k^{\text{abs}}$ is always finite
(every HAWB carries a finite $T_k^{\text{abs}}$, \S\ref{sec:hawb-params})
and is attained on the fallback arc (\S\ref{sec:fallback-arc}). The
cargo-ready upper bound on origin time is enforced at graph build via the
forward-time-window propagation
($[t^{\text{lo}}_{O_k}, t^{\text{hi}}_{O_k}] = [t_k^{\text{rdy,early}},
t_k^{\text{rdy,late}}]$, \S\ref{sec:fwd-time-propagation}); no MILP variable
$t_k^{O_k}$ exists.

The per-MAWB chargeable-weight upper bound,
\begin{equation}
  \text{CW}^{\text{ub}}_{a,g} \;\defeq\; (1 + \varepsilon) \sum_{k \in K_a : g(k) = g} \max\bigl(w_k,\, v_k \cdot 167\bigr),
  \label{eq:cw-ub}
\end{equation}
is the tightest valid upper bound from C.4a--d (sum of the larger of actual
and volumetric for each eligible HAWB, plus dunnage). All per-MAWB
big-M values in the break-disaggregation use a scalar derived from
$\text{CW}^{\text{ub}}_{a,g}$ (\S\ref{sec:bigm-tightening}).

%%--------------------------------------------------------------------
\section{Objective Function}
\label{sec:objective}
%%--------------------------------------------------------------------

\paragraph{Symbols used in this section.}
\begin{table}[H]
\centering
\caption{Symbols used in the objective function (recap of \S\ref{sec:variables}).}
\label{tab:objective-syms}
\begin{tabular}{lp{9cm}p{3.5cm}}
\toprule
Symbol & Brief & Defined \\
\midrule
\multicolumn{3}{l}{\textit{Decision variables}} \\
$x_{k,a}$ & HAWB $k$ uses arc $a$ & \S\ref{sec:variables} \\
$z_{a,g}$ & MAWB $(a, g)$ instantiated & \S\ref{sec:variables} \\
$\text{over}_c$ & Allowance overage on equalized BSA contract $c$ & \S\ref{sec:variables} \\
$\tau_k$ & Tardiness of $k$ against $\Delta_k$ & \S\ref{sec:variables} \\
\midrule
\multicolumn{3}{l}{\textit{Sets and indices}} \\
$K,\, K_a$ & HAWBs; HAWBs-on-arc $a$ & \S\ref{sec:variables} \\
$A_k$ & Pre-filtered subgraph of $k$ & \S\ref{sec:variables} \\
$A^{\text{coload}},\, A^{\text{ground}}$ & Co-load and ground arc sets & \S\ref{sec:variables} \\
$M$ & MAWB candidates & \S\ref{sec:variables} \\
$C^{\text{eq}}$ & Equalized-settlement BSA contracts & \S\ref{sec:variables} \\
$F$ & Physical flights (graph-build only) & \S\ref{sec:variables} \\
\midrule
\multicolumn{3}{l}{\textit{Parameters and derived terms}} \\
$\text{cost}^{\text{MAWB}}_{a,g}$ & Per-MAWB airline-billable cost, dispatched on $\text{rate\_family}_a$ & \S\ref{sec:lin-bucket} \\
$c^{\text{MAWB}}_{\text{fix}}$ & Per-active-MAWB documentation overhead ($\approx \$50$, \texttt{MARKET RESEARCH NEEDED}) & \S\ref{sec:air-arc-params} ``Other'' \\
$m_a^{\text{cl}}$ & Co-load per-kg rate & \S\ref{sec:air-arc-params} \\
$cw_k$ & Per-HAWB chargeable weight $\max(w_k,\, v_k \cdot 167)$ & \S\ref{sec:hawb-params} \\
$w_k$ & Per-HAWB actual weight & \S\ref{sec:hawb-params} \\
$c_a^{\text{flat}},\, c_a^{\text{kg}}$ & Ground-arc flat / per-kg handling charges & \S\ref{sec:ground-arc-params} \\
$r_c$ & BSA contract per-kg rate & \S\ref{sec:air-arc-params} \\
$\text{surcharge\_cost}(k, a)$ & Path-A per-HAWB-per-arc surcharge total & Eq.~\ref{eq:surcharge-cost} \\
$\sigma^{\text{uld}}_{a,u}$ & Graph-build-precomputed per-ULD surcharge rate on MAWB-arc $a$ for ULD type $u$ (build-up rate at origin $\text{tail}(a)$ + breakdown rate at destination $\text{head}(a)$, counted once per arc); coefficient on $\eta_{a,g,u}$ & Eq.~\ref{eq:sigma-uld} \\
$\eta_{a,g,u}$ & ULDs of type $u$ claimed on MAWB $(a,g)$ (integer) & \S\ref{sec:variables} \\
$w_{\text{sp}(k)}$ & Tardiness weight for the service product bound to $k$ (USD/h\textsuperscript{2}) & \S\ref{sec:service-products} \\
$\mu_k,\, W_k$ & Value coefficient $\mu_k = \text{value}_k / V^{\text{ref}}$ (dimensionless); combined weight $W_k = w_{\text{sp}(k)} \cdot \mu_k$ (USD/h\textsuperscript{2}) & \S\ref{sec:hawb-params} \\
$\text{pen}_k$ & PWL outer-approximation of $W_k \cdot \tau_k^2$ (continuous $\geq 0$) & \S\ref{sec:variables}, \S\ref{sec:lin-tardiness} \\
$C^{\text{fallback}}$ & Per-HAWB fallback-arc dollar cost (tenant-global) & \S\ref{sec:hawb-params}, \S\ref{sec:fallback-arc} \\
$a^{\text{fb}}_k$ & Fallback arc of HAWB $k$ ($O_k \to D_k^{\text{node}}$) & \S\ref{sec:fallback-arc} \\
\bottomrule
\end{tabular}
\end{table}

Minimize total cost plus quadratic tardiness penalty plus fallback usage cost:
\begin{align}
  \min \quad
  & \sum_{(a,g) \in M} \text{cost}^{\text{MAWB}}_{a,g}
    \tag{bucket cost on MAWBs} \\
  +\;
  & \sum_{(a,g) \in M} c^{\text{MAWB}}_{\text{fix}} \cdot z_{a,g}
    \tag{MAWB fixed-charge (issuance + customs filing)} \\
  +\;
  & \sum_{a \in A^{\text{coload}}} \sum_{k \in K_a} m_a^{\text{cl}} \cdot cw_k \cdot x_{k,a}
    \tag{co-load per-kg cost} \\
  +\;
  & \sum_{a \in A^{\text{ground}},\, k \in K_a} (c_a^{\text{flat}} + c_a^{\text{kg}} \cdot w_k) \cdot x_{k,a}
    \tag{ground handling / cartage} \\
  +\;
  & \sum_{c \in C^{\text{eq}}} r_c \cdot \text{over}_c
    \tag{BSA overage on equalized contracts} \\
  +\;
  & \sum_{k \in K,\, a \in A_k} \text{surcharge\_cost}(k, a) \cdot x_{k,a}
    \tag{Path-A per-arc surcharges} \\
  +\;
  & \sum_{(a,g) \in M,\, u \in U_a} \sigma^{\text{uld}}_{a,u} \cdot \eta_{a,g,u}
    \tag{Path-B per-ULD surcharges (linear in $\eta_{a,g,u}$), Eq.~\ref{eq:sigma-uld}} \\
  +\;
  & \sum_{k \in K} \text{pen}_k
    \tag{quadratic tardiness (PWL outer-approximation), C.10 + \S\ref{sec:lin-tardiness}} \\
  +\;
  & \sum_{k \in K} C^{\text{fallback}} \cdot x_{k,\, a^{\text{fb}}_k}
    \tag{fallback-arc usage cost, \S\ref{sec:fallback-arc}}
\end{align}

\paragraph{Monotonicity invariant.}
All cost contributors have non-negative coefficients on $\text{CW}_{a,g}$
(and on $\text{Wt}, \text{Wv}, \text{pivot}, \text{over}, \tau, \text{pen}$).
This is the
correctness condition that lets C.4c--d be inequalities ($\text{CW} \geq
\text{Wt},\, \text{CW} \geq \text{Wv}$) rather than the explicit $\max$: cost
minimization drives $\text{CW}$ down to $\max(\text{Wt}, \text{Wv})$ exactly.
\textbf{If any future rate-family extension introduces a negative-coefficient
term on $\text{CW}_{a,g}$ (rebate, marketing credit), C.4c--d must be
tightened to equality via a PWL formulation.} The current model does not
permit such terms; a post-solve assertion suite checks this invariant on
every solve (\texttt{TEST\_PLAN.md \S10}).

\paragraph{Bucket cost term --- dispatched on $\text{rate\_family}_a$.}
$\text{cost}^{\text{MAWB}}_{a,g}$ is the airline-billable cost of MAWB
$(a, g)$, defined per family in \S\ref{sec:lin-bucket} below.

\paragraph{MAWB fixed-charge term.}
$c^{\text{MAWB}}_{\text{fix}}$ (USD per active MAWB, defined in
\S\ref{sec:air-arc-params} under ``Other'') captures the per-MAWB documentation
overhead: customs manifest filing (AMS for US imports $\approx \$30$/MAWB; ENS
for EU imports of similar magnitude), forwarder MAWB issuance handling, and
related fixed administrative costs that scale with the number of MAWBs, not
with weight. Wired to the existing activation binary $z_{a,g}$
(\S\ref{sec:variables}). Without this term the MILP has no penalty for
splitting a single consolidatable batch across multiple $(a, g)$ MAWBs at
identical variable cost --- the consolidation decision becomes degenerate.
Placeholder $\$50$/MAWB; \texttt{MARKET RESEARCH NEEDED}.

\paragraph{Sell rate / margin out of scope.}
The MILP outputs cost only. Sell rate is set by a separate quote engine at
booking time (cost-plus-margin, market-rate comparison, competitor signals);
once a HAWB is booked, the price is fixed and routing is pure cost
minimization. Margin = sell $-$ cost is composed by the quote engine for
operator-facing screens (\texttt{shipment\_attributes.md \S2.5}).

%%--------------------------------------------------------------------
\section{Linearization Details}
\label{sec:linearization}
%%--------------------------------------------------------------------

\paragraph{Symbols used in this section.}
\begin{table}[H]
\centering
\caption{Symbols used in the linearization details (recap of \S\ref{sec:variables}).}
\label{tab:linearization-syms}
\begin{tabular}{lp{9cm}p{3.5cm}}
\toprule
Symbol & Brief & Defined \\
\midrule
\multicolumn{3}{l}{\textit{Decision variables (existing)}} \\
$z_{a,g}$ & MAWB $(a, g)$ instantiated (binary) & \S\ref{sec:variables} \\
$\text{CW}_{a,g}$ & MAWB chargeable weight & \S\ref{sec:cw-density-mixing} \\
$\gamma_{a,g,b}$ & Break selector for \texttt{min\_flat\_breaks} (binary) & \S\ref{sec:variables} \\
$\text{BW}_{a,g,b}$ & Disaggregated billed weight at break $b$ & \S\ref{sec:variables} \\
$\eta_{a,g,u}$ & ULDs of type $u$ claimed on MAWB (integer) & \S\ref{sec:variables} \\
$\text{pivot}_{a,g}$ & Pivot-active billed weight (per-flight BSA) & \S\ref{sec:variables} \\
$\tau_k,\, \text{pen}_k$ & Tardiness and its quadratic-penalty PWL value & \S\ref{sec:variables} \\
\multicolumn{3}{l}{\textit{Auxiliaries introduced here}} \\
$c_{a,g}$ & Linearization auxiliary for $\text{cost}^{\text{MAWB}}_{a,g}$ ($\geq 0$); used in the objective in place of the $\max$ form & \S\ref{sec:lin-bucket} \\
\midrule
\multicolumn{3}{l}{\textit{Sets and indices}} \\
$M,\, A^{\text{MFB}},\, B_a,\, A^{\text{pu}},\, U_a$ & MAWB candidates; \texttt{min\_flat\_breaks} arcs; weight breaks; per-ULD-pivot arcs; admissible ULDs & \S\ref{sec:variables} \\
\midrule
\multicolumn{3}{l}{\textit{Parameters}} \\
$\hat\tau_{k,j}$ & Per-HAWB tangent-breakpoint grid for the quadratic-tardiness PWL, $\hat\tau_{k,j} = \alpha_j (T_k^{\text{abs}} - \Delta_k)$ (precomputed scalar, \emph{not} a variable); default $\alpha = \{0, 0.25, 0.5, 0.75, 1.0\}$ & \S\ref{sec:lin-tardiness} \\
$m_a,\, \text{min\_chg}_a,\, \text{cap}_a$ & \texttt{flat\_rate} per-kg rate / minimum charge / cap & \S\ref{sec:air-arc-params} \\
$\text{rate}_{a,b},\, \text{break}_{a,b}$ & Per-break rate and weight floor (\texttt{min\_flat\_breaks}) & \S\ref{sec:air-arc-params} \\
$\pi_a,\, r_a$ & BSA pivot weight; per-kg BSA rate & \S\ref{sec:air-arc-params} \\
$\text{settlement}_a$ & BSA settlement mode (\texttt{per\_flight} / \texttt{equalized}) & \S\ref{sec:air-arc-params} \\
$\text{CW}^{\text{ub}}_{a,g}$ & Per-MAWB chargeable-weight upper bound (big-M) & \S\ref{sec:variables} \\
$w_{\text{sp}(k)},\, \mu_k,\, W_k$ & Tardiness weights (service-product / value coefficient / combined) & \S\ref{sec:hawb-params}, \S\ref{sec:service-products} \\
\midrule
\multicolumn{3}{l}{\textit{Derived terms}} \\
$\text{cost}^{\text{MAWB}}_{a,g}$ & Per-MAWB airline-billable cost, dispatched on $\text{rate\_family}_a$ (closed-form per family in this section) & \S\ref{sec:lin-bucket} \\
\bottomrule
\end{tabular}
\end{table}

\subsection{Bucket cost linearization}
\label{sec:lin-bucket}

\paragraph{\texttt{flat\_rate}.}
$\text{cost}^{\text{MAWB}}_{a,g} = \max(\text{min\_chg}_a \cdot z_{a,g},\,
m_a \cdot \text{CW}_{a,g})$. Two inequalities, no binary --- introduce
auxiliary $c_{a,g} \geq 0$:
\[
  c_{a,g} \;\geq\; \text{min\_chg}_a \cdot z_{a,g}, \qquad
  c_{a,g} \;\geq\; m_a \cdot \text{CW}_{a,g}.
\]
Objective uses $c_{a,g}$. The $z_{a,g}$ factor on the minimum charge keeps an
empty bucket at zero cost (when $z = 0$ the upper-link bound forces
$\text{CW} = 0$, so both right-hand sides are $0$).

\paragraph{\texttt{min\_flat\_breaks} (TACT / SCR / CCR --- item 7).}
The IATA next-break-down rule, Eq.~\ref{eq:rate-fn-minflat}:
\begin{equation}
  \text{cost}^{\text{MAWB}}_{a,g} \;=\; \min_{b \in B_a} \text{rate}_{a,b} \cdot \max(\text{CW}_{a,g},\, \text{break}_{a,b}).
  \label{eq:rate-fn-minflat}
\end{equation}
The ``round up to a higher break for a lower rate'' case \emph{must} be
representable --- it is the whole point of break-down rating. Encoded as:
\begin{align}
  \sum_{b \in B_a} \gamma_{a,g,b} \;&=\; z_{a,g}
    \tag{exactly one break IF active} \\
  \text{BW}_{a,g,b} \;&\leq\; M^{\text{BW}}_{a,g} \cdot \gamma_{a,g,b}
    \quad \forall b \tag{force $\text{BW} = 0$ when $\gamma = 0$} \\
  \text{BW}_{a,g,b} \;&\geq\; \text{break}_{a,b} \cdot \gamma_{a,g,b}
    \quad \forall b \tag{when $\gamma = 1$: $\text{BW} \geq \text{break}$} \\
  \text{BW}_{a,g,b} \;&\geq\; \text{CW}_{a,g} - M^{\text{BW}}_{a,g} \cdot (1 - \gamma_{a,g,b})
    \quad \forall b \tag{when $\gamma = 1$: $\text{BW} \geq \text{CW}$} \\
  \text{cost}^{\text{MAWB}}_{a,g} \;&=\; \sum_{b \in B_a} \text{rate}_{a,b} \cdot \text{BW}_{a,g,b}
\end{align}
Three inequalities on $\text{BW}_b$ (not four). \textbf{No $\text{BW}_b \leq
\text{CW}$.} That constraint would force $\text{BW}_{b^*} = \text{CW}$ for
the chosen break; combined with $\text{BW}_{b^*} \geq \text{break}_{b^*}$ it
would make any selection with $\text{break}_{b^*} > \text{CW}$ infeasible,
\emph{banning the round-up-to-higher-break case}. Without it: $\text{BW}_b =
\max(\text{CW},\, \text{break}_b)$ when $\gamma_b = 1$; $\text{BW}_b = 0$
when $\gamma_b = 0$; minimization over the break binaries naturally recovers
$\min_b \text{rate}_b \cdot \max(\text{CW},\, \text{break}_b)$ = IATA cost.

\paragraph{Worked check.}
90 kg HAWB; breaks $(b = 45, \text{rate} = \$10/\text{kg})$ and
$(b = 100, \text{rate} = \$8/\text{kg})$:
\begin{itemize}[noitemsep]
  \item $\gamma_{45} = 1$: $\text{BW}_{45} = \max(90, 45) = 90$. Cost $= 10 \cdot 90 = \$900$.
  \item $\gamma_{100} = 1$: $\text{BW}_{100} = \max(90, 100) = 100$. Cost $= 8 \cdot 100 = \$800$.
\end{itemize}
Min over breaks: $\$800$ --- matches IATA (round up to $b = 100$ for the
cheaper rate). The earlier 4-inequality form (with the spurious $\text{BW}_b
\leq \text{CW}$) would have forced $\$900$. Caught by the Session 15
linearization-agent critique; fix verified.

\paragraph{Empty bucket --- $\Sigma \gamma = z$.}
The exactly-one-if-active form ($\sum_b \gamma_{a,g,b} = z_{a,g}$, not
$\sum_b \gamma = 1$) handles empty MAWBs: $z_{a,g} = 0 \Rightarrow$ all
$\gamma_b = 0 \Rightarrow$ all $\text{BW}_b = 0 \Rightarrow$ cost contribution
$= 0$, without a phantom minimum charge.

$M^{\text{BW}}_{a,g}$ is the per-MAWB-tight big-M --- see
\S\ref{sec:bigm-tightening}.

\paragraph{\texttt{per\_uld\_pivot}, \texttt{per\_flight} settlement.}
$\text{cost}^{\text{MAWB}}_{a,g} = r_a \cdot \text{pivot}_{a,g}$ with
$\text{pivot}_{a,g}$ from C.13b-1, C.13b-2. Standard two-inequality $\max$
linearization; tight because $r_a > 0$.

\paragraph{\texttt{per\_uld\_pivot}, \texttt{equalized} settlement.}
$\text{cost}^{\text{MAWB}}_{a,g} = 0$ for equalized \texttt{per\_uld\_pivot}
arcs: \emph{all} cost on these arcs is realized through $r_c \cdot
\text{over}_c$ in the objective and the C.13a allowance overage. The
chargeable weight contributes to $\text{chargeable}(c) = \sum_{(a,g):\, a \in
A_c^{\text{MAWB}}} \text{CW}_{a,g}$ (a constraint accounting quantity, not a
cost) which feeds C.13a. The pivot floor is handled at the \emph{contract}
level by $A_c$ and $\text{over}_c$, not per-MAWB; C.13b is \emph{not}
generated for equalized arcs. Per-ULD physical capacity C.5b still applies.

\subsection{Pivot weight linearization (per-flight BSA)}
\label{sec:lin-pivot}

For each MAWB $(a, g)$ with $a \in A^{\text{pu}}$ and $\text{settlement}_a =
\texttt{per\_flight}$, the native cost is
\[
  r_a \cdot \max\bigl(\text{CW}_{a,g},\, \pi_a \cdot \sum_u \eta_{a,g,u}\bigr).
\]
Replaced by $r_a \cdot \text{pivot}_{a,g}$ with the two inequalities C.13b-1,
C.13b-2 above. Adds one continuous variable and two inequalities per
$(a, g)$ on $A^{\text{pu}}$. P1: a disaggregated formulation
($\eta_{a,g,u} = \sum_n \eta^c_{a,g,u,n}$ with per-ULD-slot binaries) yields
a tighter LP relaxation at the cost of $O(N_{a,u})$ extra binaries; defer
until LP-gap-source instrumentation indicates this family dominates.

\subsection{Quadratic tardiness linearization (outer approximation)}
\label{sec:lin-tardiness}

The C.10 quadratic penalty $W_k \cdot \tau_k^2$ is replaced in the MILP by
$\text{pen}_k$ (continuous $\geq 0$, \S\ref{sec:variables}) constrained from
below by a fixed family of $J$ tangent cuts. The breakpoints are placed
\emph{per HAWB}, relative to that HAWB's own feasible tardiness range
$R_k = T_k^{\text{abs}} - \Delta_k$:
\begin{equation*}
  \hat\tau_{k,j} = \alpha_j \cdot R_k,
  \qquad \alpha \in \{0,\, 0.25,\, 0.5,\, 0.75,\, 1.0\},
  \qquad j = 1, \ldots, J.
\end{equation*}
For each $k \in K$ and each breakpoint $\hat\tau_{k,j}$, the cut from the
tangent of $f(\tau) = W_k \cdot \tau^2$ at $\tau = \hat\tau_{k,j}$ is
\begin{equation}
  \text{pen}_k \;\geq\; W_k \cdot \bigl(\hat\tau_{k,j}^2 + 2\hat\tau_{k,j} (\tau_k - \hat\tau_{k,j})\bigr)
            \;=\; W_k \cdot \bigl(2\hat\tau_{k,j} \cdot \tau_k - \hat\tau_{k,j}^2\bigr)
  \label{eq:lin-tardiness}
\end{equation}
The minimization of $\sum_k \text{pen}_k$ in the objective forces
$\text{pen}_k$ down to the upper envelope of these cuts, which is the
piecewise-linear outer approximation of the convex quadratic. As $J \to
\infty$, the approximation $\to$ the true quadratic from below. No binary or
SOS-2 needed --- the family of $\geq$ constraints alone suffices because
$\tau^2$ is convex. The coefficients $\hat\tau_{k,j}$ are precomputed scalars
at graph build (both $\alpha_j$ and $R_k$ are known), so the cut remains
linear in the decision variable $\tau_k$.

\paragraph{Why a relative grid.}
A fixed \emph{absolute} grid (e.g.\ $\{0,4,12,36,96\}$ h) breaks on
wide-backstop products. The last breakpoint caps the tangent family, so for
any $\tau_k$ beyond it the approximation degenerates to the single tangent at
that breakpoint and extrapolates \emph{linearly} while the true penalty grows
quadratically. For GEN with a $\sim$30-day cancellation backstop ($R_k \sim
720$ h) this under-prices a very-late real route by $\gtrsim 4\times$ at the
upper end, and the optimizer then prefers that route over the
correctly-priced fallback arc (\S\ref{sec:fallback-arc}), whose tardiness is
exactly $R_k$. Anchoring the grid to $R_k$ removes the extrapolation
entirely: the top breakpoint $\alpha_J = 1.0$ lands on $\hat\tau_{k,J} = R_k$,
the same point the fallback arc fires, so late-real-route vs.\ fallback is
always a like-for-like comparison. The $\alpha_1 = 0$ cut gives the floor
$\text{pen}_k \geq 0$.

\paragraph{Why even spacing.}
The outer-approximation gap on interval $[\hat\tau_{k,j}, \hat\tau_{k,j+1}]$
peaks at the midpoint at $W_k (\hat\tau_{k,j+1} - \hat\tau_{k,j})^2 / 4
= W_k (\alpha_{j+1} - \alpha_j)^2 R_k^2 / 4$ --- quadratic in the interval
\emph{width}. With a fixed point budget the worst-case gap is therefore set by
the widest interval, and is minimized by equal widths. Even spacing gives a
uniform worst-case gap of $W_k R_k^2 / 64$ across all four intervals; a
low-skewed alternative such as $\{0, 0.05, 0.15, 0.5, 1.0\}$ would buy
resolution near zero at the cost of a $4\times$ larger tail gap
($W_k R_k^2 \cdot 0.0625$ on $[0.5 R_k, R_k]$) --- precisely the region the
relative grid exists to price correctly. The low-end resolution lost to even
spacing is recovered automatically for tight-backstop products, since the
grid rescales to each HAWB's range.

\paragraph{Cost (constraint count).}
$J \cdot |K|$ additional linear inequalities and $|K|$ additional continuous
variables. With $J = 5$ and $|K| = 100$ that is $500$ inequalities (vs.\ $100$
under the prior linear C.10), a $\sim 5\%$ increase in total constraint count
at MVP scale --- unchanged from the absolute grid, since the relative grid
moves breakpoints, not their number. Recorded in \S\ref{sec:scaling-roadmap}.

\paragraph{Approximation tightness.}
At a grid point $\hat\tau_{k,j}$ the cut is exact: $\text{pen}_k = W_k
\hat\tau_{k,j}^2$. Between grid points the LP relaxation may under-estimate the
true quadratic by at most $W_k R_k^2 / 64$ (even-grid midpoint bound);
integer-feasibility recovers the correct $\tau_k$ value, and the cuts at
adjacent breakpoints bound $\text{pen}_k$ from below on both sides. The bound
scales with $R_k^2$, so absolute gap is tight for short backstops and looser
for long ones --- acceptable because long-backstop penalties are themselves
soft commercial estimates, not hard time bounds.

\paragraph{Tenant override.}
The $\alpha$ vector is a tenant-config list, not a hardcoded constant.
Operators with empirical lateness distributions can refine it --- e.g.\ a
6-point $\{0, 0.1, 0.25, 0.5, 0.75, 1.0\}$ to add low-end slope fidelity
without widening any interval past $0.25 R_k$, at the cost of one extra cut
per HAWB. The grid stays relative to $R_k$; only the partition of $[0,1]$
changes.

\subsection{Linearization summary}
\label{sec:lin-summary}

\begin{table}[H]
\centering
\caption{Linearization summary --- mechanism, trigger, and technique per rate family / nonlinearity.}
\label{tab:lin-summary}
\begin{tabular}{p{4.0cm}p{4.0cm}p{6.5cm}}
\toprule
Mechanism & Trigger & Technique \\
\midrule
$\text{CW} = \max(\text{Wt}, \text{Wv})$ & all MAWBs & Two $\geq$ inequalities (C.4c, C.4d). Correctness via monotonicity invariant. \\
$\max(\text{min\_chg} \cdot z,\, m \cdot \text{CW})$ & \texttt{flat\_rate} & Aux $c \geq$ both. Empty-bucket: $z = 0 \Rightarrow \text{CW} = 0 \Rightarrow c = 0$. \\
$\min_b \text{rate}_b \cdot \max(\text{CW},\, \text{break}_b)$ & \texttt{min\_flat\_breaks} & Binary $\gamma_b$ exactly-one-if-active ($\sum \gamma = z$) $+$ \textbf{3-inequality} $\text{BW}_b$ disaggregation (no $\text{BW} \leq \text{CW}$). Verified worked example. \\
$\max(\text{CW},\, \pi \sum \eta)$ & \texttt{per\_uld\_pivot}, \texttt{per\_flight} & Two $\geq$ inequalities (C.13b). \\
$\max(0,\, \text{chargeable} - A)$ & \texttt{per\_uld\_pivot}, \texttt{equalized} & One $\geq$ $+$ non-negativity (C.13a). \\
$\max(0,\, t_k^{D^{\text{node}}} - \Delta)$ & soft tardiness $\tau_k$ & One $\geq$ $+$ non-negativity (C.10b). \\
$W_k \cdot \tau_k^2$ (quadratic penalty) & C.10b quadratic & Outer-approximation of convex quadratic via $J$ tangent cuts on $\text{pen}_k$ at a fixed breakpoint grid; no binary needed (\S\ref{sec:lin-tardiness}). \\
Destination arrival from terminal arc & C.10a & Linear definition $t_k^{D_k^{\text{node}}} = \sum x_{k,a} \cdot \text{arr\_dest}(k,a)$ over terminal arcs only; no big-M (\S\ref{sec:fwd-time-propagation} pushes time feasibility to graph build). \\
MAWB activation & C.2 & $\Sigma x$-bounded; no auxiliary. \\
Bucket aggregate upper-link & C.14 & $\text{CW} \leq \text{CW}^{\text{ub}} \cdot z$ ties bucket cost vars to MAWB activation; tightens LP. \\
\bottomrule
\end{tabular}
\end{table}

\subsection{Per-constraint tight big-M values}
\label{sec:bigm-tightening}

Loose big-M values weaken the LP relaxation and slow solving. Per Wolsey
(\emph{Integer Programming} 1998 \S1.7) and Bertsimas--Tsitsiklis
(\emph{LP Intro} \S11), the tight $M$ for an indicator constraint is the
maximum amount the constraint must be relaxed when the indicator deactivates
--- constraint-specific, not horizon-wide.

\begin{table}[H]
\centering
\caption{Per-constraint tight big-M values.}
\label{tab:bigm-tightening-tab}
\begin{tabular}{p{2.8cm}p{4.0cm}p{7.5cm}}
\toprule
Constraint & Tight $M$ & Rationale \\
\midrule
$M^{\text{BW}}_{a,g}$ (TACT break disaggregation) & $\max\bigl(\text{CW}^{\text{ub}}_{a,g},\; \max_{b \in B_a} \text{break}_{a,b}\bigr)$ where $\text{CW}^{\text{ub}}_{a,g}$ is from Eq.~\ref{eq:cw-ub} & Per-MAWB tight: bounds both $\text{BW}_b \leq M \gamma_b$ and $\text{BW}_b \geq \text{CW} - M(1 - \gamma_b)$. \textbf{The $\max$ with $\max_b \text{break}_{a,b}$ is load-bearing}: without it, any break with $\text{break}_{a,b} > \text{CW}^{\text{ub}}_{a,g}$ becomes infeasible (the upper-link $\text{BW}_b \leq M$ conflicts with $\text{BW}_b \geq \text{break}_{a,b}$), silently banning the IATA round-up-to-higher-break case --- exactly the case the 3-inequality form was designed to enable. The §\ref{sec:appendix-tact} worked example (280 kg MAWB facing breaks 45/100/300/500/1000, optimum at break 300 = \$1260) requires this widening to be reachable. \\
\bottomrule
\end{tabular}
\end{table}

\paragraph{Implementation note.}
Big-M values are precomputed at MILP build from per-MAWB parameters.
Solver-side presolve will further tighten through bounds propagation; the
model-side values give it a strong starting point.

%%--------------------------------------------------------------------
\section{Output and Diagnostics}
\label{sec:output-diagnostics}
%%--------------------------------------------------------------------

Because the MILP is always feasible by construction (the fallback arc
guarantees a valid solution for every HAWB, \S\ref{sec:fallback-arc}),
solver-infeasibility is no longer the rescue signal. The orchestrator reads
the solve output and reacts to fallback usage and tardiness instead.

\paragraph{Reported quantities per solve.}
The solver output exposes aggregate cost numbers, the fallback set, and a
per-HAWB attribution panel:
\begin{table}[H]
\centering
\caption{Quantities reported per solve, plus the per-HAWB diagnostic surface.}
\label{tab:output-quantities}
\begin{tabular}{lp{9.0cm}}
\toprule
Field & Meaning \\
\midrule
$z^{\text{routing}}$ & Real routing cost: sum of all objective terms
\emph{excluding} fallback-arc cost and tardiness penalty. The number
operators see on planning screens as ``cost to serve.'' \\
$z^{\text{tardiness}}$ & Total quadratic tardiness penalty: $\sum_k \text{pen}_k$
(approximating $\sum_k W_k \cdot \tau_k^2$). \\
$z^{\text{fallback}}$ & Fallback usage cost: $\sum_k C^{\text{fallback}}
\cdot x_{k,\, a^{\text{fb}}_k}$. Equals $C^{\text{fallback}} \cdot
|K^{\text{fb}}|$ where $K^{\text{fb}} = \{k : x_{k,\, a^{\text{fb}}_k} = 1\}$. \\
$z^{\text{total}}$ & Total solver objective: $z^{\text{routing}} + z^{\text{tardiness}} + z^{\text{fallback}}$. \\
$K^{\text{fb}}$ & List of HAWB IDs that took the fallback arc this solve. \\
$|K^{\text{fb}}|$ & Count of fallback HAWBs. \\
$\tau_k^{\text{hr}}$ & Per-HAWB tardiness in hours: $\tau_k$ (the soft-tardiness decision variable, \S\ref{sec:variables}). One row per HAWB; surfaced directly in the operator-facing per-shipment view. \\
$c_k^{\text{attr}}$ & Per-HAWB attributed cost in USD --- see Per-HAWB cost attribution below. One row per HAWB. \\
\bottomrule
\end{tabular}
\end{table}

\paragraph{Per-HAWB cost attribution (proportional-to-CW rule).}
The MILP returns per-MAWB cost via $\text{cost}^{\text{MAWB}}_{a,g}$
(\S\ref{sec:lin-bucket}); operator workflows (quote-desk margin review,
customer billing reconciliation, KAM QBR analysis) need per-HAWB cost.
Attribution is computed post-solve, not inside the MILP, by the
\emph{proportional-to-chargeable-weight} rule. For each HAWB $k$:
\begin{align}
  c_k^{\text{attr}} \;\defeq\;
    & \sum_{(a, g) :\, a \in A_k,\, g(k) = g,\, x_{k,a} = 1}
      \text{cost}^{\text{MAWB}}_{a,g} \cdot
      \frac{\text{cw}^{\text{attr}}_{k, a}}{\text{CW}_{a, g}}
    \tag{air bucket cost share} \\
  +\;
    & \sum_{a \in A^{\text{coload}},\, x_{k,a} = 1}
      m_a^{\text{cl}} \cdot cw_k
    \tag{co-load per-kg, direct} \\
  +\;
    & \sum_{a \in A^{\text{ground}},\, x_{k,a} = 1}
      (c_a^{\text{flat}} + c_a^{\text{kg}} \cdot w_k)
    \tag{ground arcs, direct} \\
  +\;
    & \sum_{a \in A_k,\, x_{k,a} = 1}
      \text{surcharge\_cost}(k, a)
    \tag{Path-A surcharges, direct} \\
  +\;
    & c^{\text{MAWB,attr}}_{k}
    \tag{MAWB fixed-charge share} \\
  +\;
    & f^{\text{Path-B},\text{attr}}_{k}
    \tag{Path-B per-flight share} \\
  +\;
    & \sum_{c \in C^{\text{eq}}}
      r_c \cdot \text{over}_c \cdot
      \frac{\text{cw}^{\text{attr}}_{k, c}}{\text{chargeable}(c)}
    \tag{BSA overage share, equalized}
  \label{eq:per-hawb-attribution}
\end{align}
where the per-HAWB chargeable-weight share on MAWB $(a, g)$ is
$\text{cw}^{\text{attr}}_{k, a} = \max(w_k,\, v_k \cdot 167) / \sum_{k' \in
K_a : g(k') = g,\, x_{k', a} = 1} \max(w_{k'},\, v_{k'} \cdot 167)$
times $\text{CW}_{a, g}$ (i.e.\ HAWB $k$'s share of MAWB chargeable
weight equals its per-HAWB pre-aggregation contribution as a fraction of
the sum across activated HAWBs in the same group). MAWB-level
fixed-charge $c^{\text{MAWB,attr}}_{k}$ and Path-B per-flight shares
$f^{\text{Path-B},\text{attr}}_{k}$ are allocated proportionally by the
same CW share. BSA overage on equalized contracts is allocated across
HAWBs riding the contract by their CW share of the contract-level
$\text{chargeable}(c)$ that exceeds the allowance.

The rule is exact at the aggregate ($\sum_k c_k^{\text{attr}} =
z^{\text{routing}}$) and operationally defensible: the airline bills the
MAWB at $\text{CW}$, the forwarder bills the shipper proportional to that
HAWB's contribution to the MAWB's billing weight. Alternative attribution
rules (Shapley-value, proportional-to-actual-weight,
proportional-to-volume) are deferred (\S\ref{sec:deferred},
item~\ref{item:per-hawb-cost-attribution}).

\paragraph{Rescue signal.}
Any $k \in K^{\text{fb}}$ is a rescue event. The orchestrator
(\texttt{build\_plan.md}) inspects each fallback HAWB and reads
$\text{predicate}^{\text{dom}}_k$ --- the dominant pre-filter
rejection predicate --- from the pre-filter rejection log
(\S\ref{sec:prefilter}). Routing the rescue is then a function of the
dominant predicate:
\begin{itemize}[noitemsep]
  \item predicate 2 (cargo type) / 3 (embargo) / 4 (lithium) failure:
        usually a data-curation issue or a genuine no-supply case;
        operator triages with supply procurement;
  \item predicate 5 (service product) / carrier-policy denied all
        carriers: renegotiate the policy or upgrade the service
        product;
  \item predicate 6 (forward-time-window) failure: cutoff missed;
        renegotiate cargo-ready window with shipper or accept later
        $T_k^{\text{abs}}$;
  \item predicate 7 (destination reachability) failure: the HAWB's
        backstop is too tight for any feasible route; escalate to
        customer-success for SLA-breach handling.
\end{itemize}
The rescue is structured per-HAWB, with attribution, rather than a
batch-level ``solver returned INFEASIBLE.''

\paragraph{Diagnostic invariants asserted post-solve.}
On every solve, a post-solve assertion suite checks
(\texttt{TEST\_PLAN.md \S10}):
\begin{itemize}[noitemsep]
  \item $\text{pen}_k = W_k \cdot \tau_k^2$ at integer feasibility (within
        the PWL grid approximation gap; the assertion uses a relative
        tolerance from the grid).
  \item $z^{\text{fallback}} = C^{\text{fallback}} \cdot |K^{\text{fb}}|$
        (linear identity, exact).
  \item For every $k \in K^{\text{fb}}$: $t_k^{D_k^{\text{node}}} =
        T_k^{\text{abs}}$ (fallback arrival is exactly $T_k^{\text{abs}}$).
  \item For every $k \notin K^{\text{fb}}$: $\sum_{a \in A_k :\, a \neq
        a^{\text{fb}}_k} x_{k,a} = 1$ at each cut node (cargo took a real
        route).
\end{itemize}
Violation of any invariant indicates a silent bug --- raise a
\texttt{invariant\_violations.json} entry and fail the solve loudly.

\paragraph{Operator-facing presentation.}
Planning screens display $z^{\text{routing}}$ as the primary cost. Tardiness
and fallback are shown as flags with counts and per-HAWB drill-down (the
flag taxonomy from \texttt{agent\_design.md}: SLA risk, capacity risk,
disruption). $C^{\text{fallback}}$ is never shown as a dollar number on
planning screens --- it is a mathematical device for guaranteeing
feasibility, not an actual price.

%%--------------------------------------------------------------------
\section{Open Items and Future Extensions}
\label{sec:deferred}
%%--------------------------------------------------------------------

Recorded for the next phase of work. Items below are \emph{not} in MVP.

\subsection*{Deferred P1 items (worth doing later)}

\begin{enumerate}
  \item \label{item:tt-quantile-binding} \textbf{TT-Service P85--P90 quantile
        binding.} Replace deterministic arc transits with a TT-Service
        end-to-end P85--P90 quantile (Finding S Ch~1 hook;
        \texttt{transit\_time\_model.md} for the Phase 2 quantile service).
        \textbf{Supported migration path (only):} precompute the end-to-end
        P-quantile arrival per terminal arc and bind it into
        $\text{arr\_dest}(k, a)$ entering C.10a. Leave the forward-time-window
        propagation on a deterministic statistic (e.g.\ mean) used only for
        cutoff and reachability admission. \textbf{Do not} lift the
        propagation to per-arc quantile arithmetic: the sum of P-quantiles is
        not the P-quantile of the sum (convolution problem), so naive
        propagation under-estimates end-to-end uncertainty. End-to-end
        quantile arithmetic requires per-path analysis --- which falls out
        naturally from column generation if and when the formulation pivots
        to path-based (\texttt{scalability.md}). Three structurally distinct
        migrations exist: (a) quantile-bound arrival in $\text{arr\_dest}$
        (arc-based-compatible, this item); (b) expected $\tau^2$ over a
        scenario distribution (requires scenario decomposition or per-arc
        moment propagation); (c) chance-constrained $\Pr(\text{arrival} \leq
        T_k^{\text{abs}}) \geq p$ or CVaR-on-tardiness (path-based natural,
        see item~\ref{item:probabilistic-transit}). Naming the chosen path
        explicitly prevents committing to (c) and shipping only (a).
  \item \label{item:offload-priority} \textbf{Offload priority into the
        MILP.} Finding S Ch~2: $\text{supply\_class} \in \{\texttt{confirmed},
        \texttt{best\_effort}\}$ as a per-offer attribute $+$ a
        $\text{confirmed\_only}$ product attribute $+$ a subgraph pre-filter
        predicate (express $\Rightarrow$ confirmed-capacity options only).
        Tier 1 air products sell offload priority more than door-to-door
        clock; the MVP catalog doesn't surface this yet.
  \item \label{item:sla-endpoint-and-hops} \textbf{SLA endpoint and max-hops
        product attributes.} Finding S Ch~3: A2A vs D2D SLA endpoint on
        service products; $\text{max\_hops}$ / $\text{direct\_only}$ product
        attributes for express tiers.
  \item \label{item:quadratic-tardiness} \textbf{Quadratic tardiness ---
        PROMOTED to MVP.} The $W_k \cdot \tau_k^2$ penalty (PWL outer
        approximation, \S\ref{sec:lin-tardiness}) replaced the prior linear
        $w \cdot \tau_k$ form. Combined with the per-HAWB fallback arc
        (\S\ref{sec:fallback-arc}), the model is always feasible and
        prefers \emph{reasonable} lateness over concentrated breach. Retained
        in this list for changelog traceability only.
  \item \label{item:volume-kicker-accumulator} \textbf{Period-scoped volume
        kickers via external accumulator.} Retroactive / all-units / period-
        or network-scoped volume-discount kickers (\S\ref{sec:carrier-policy},
        ``Volume kickers: monotone vs.\ period-scoped''). The cumulative
        threshold is not separable into a per-batch solve and the cost curve is
        non-convex (incremental) or discontinuous (all-units), so it cannot
        enter the per-batch objective directly. \textbf{Handled identically to
        equalized BSA settlement} (\S\ref{sec:bsa-equalized-settlement}): an
        external accumulator maintains period-to-date volume per incentivized
        contract and feeds the per-batch MILP a control input --- a remaining
        allowance, or a per-kg price signal $\delta_c$ that lowers the rate by
        a constant when within a band of the threshold, keeping each batch
        monotone and $\geq$-cut-safe. Open design: the exact $\delta_c$ (a
        forecast of remaining period volume, vs.\ a Lagrangian multiplier on a
        target-volume coupling constraint dualized across the rolling horizon).
        The retro down-jump modeled exactly in-objective (discontinuous PWL,
        Vielma--Ahmed--Nemhauser) is a further deferral --- the $\delta_c$
        control input approximates it and avoids the integrality cost.
        Methodology anchor: \texttt{references/air-cargo-allotment-contracts.md}
        (Gupta 2008; Amaruchkul 2018).
  \item \label{item:split-shipment} \textbf{Split shipment.} One HAWB across
        two flights / MAWBs (capacity / partial-roll). Option 1 (virtual
        sub-HAWBs at intake, parent-HAWB binding for billing) preferred ---
        the MILP is unchanged.
  \item \label{item:pairwise-dg-segregation} \textbf{Pairwise DG segregation
        at the ULD layer.} IATA segregation tables; pairwise, non-transitive
        --- not a partition over groups, so it lives at the ULD layer, not
        the MAWB layer. Adds a per-ULD constraint family.
  \item \label{item:temp-band-refinement} \textbf{Temperature-band
        refinement within PER} (e.g., pharma sub-bands $+2$ to $+8\,^\circ$C
        vs $+15$ to $+25\,^\circ$C). The structural form for this is a
        \emph{poset} on temperature regimes (pharma $\preceq$ chilled
        $\preceq$ ambient, where $\preceq$ means ``validation envelope
        contains''): HAWBs share a MAWB if their temperature tuples agree
        on the infimum. Still a partition (the equivalence class is
        ``shipments that agree on the infimum''), but the partition's
        cells are coarser than naive sub-banding. P1 design pass for
        $g(\cdot)$ refinement.
  \item \label{item:miqp-tardiness} \textbf{Direct MIQP tardiness when
        HiGHS gains MIQP support.} The PWL outer-approximation of
        $W_k \tau_k^2$ (\S\ref{sec:lin-tardiness}) is a lower bound on
        the true quadratic and biases the optimizer's argmin toward
        $\tau_k$ values between grid points (where penalty is
        underestimated). When HiGHS adds MIQP capability (per HiGHS
        roadmap, $\sim$1-year horizon), replace the PWL family with a
        direct $W_k \tau_k^2$ term in the objective and drop the
        $\text{pen}_k$ auxiliary. Eliminates the grid-calibration
        problem entirely.
  \item \label{item:forecast-aware-accumulator} \textbf{Forecast-aware BSA
        accumulator.} The current equalized-BSA accumulator $A_c$
        (\S\ref{sec:bsa-equalized-settlement}) does not see remaining
        period demand. If the period is half over and $A_c$ is
        half-consumed but the remaining half-period has 2$\times$ the
        expected demand (seasonal / lane-tide), the marginal pricing is
        wrong: this batch consumes against an allowance that will be
        exhausted by future demand. Forecast-aware form:
        $A_c^{\text{effective}} = A_c \cdot (\text{remaining period} /
        \text{expected remaining demand fraction})$. Surface as the
        free-quota in C.13a. P1 once demand forecasts are available.
  \item \label{item:slot-symmetry-breaking} \textbf{Slot-symmetry
        breaking on per-ULD-slot disaggregation.} If the per-ULD-slot
        disaggregated form $\eta_{a,g,u,n}$ (per
        \S\ref{sec:lin-pivot} P1 note) is ever promoted, the
        $n$-th ULD slot is interchangeable with the $(n+1)$-th --- a
        classical labeling symmetry. Break it with ordering constraints
        $\eta_{a,g,u,n} \geq \eta_{a,g,u,n+1}$ for $n = 1, \ldots,
        N_{a,u} - 1$. Cheap and standard (Margot 2010).
  \item \label{item:per-flight-fx-pinning} \textbf{Per-run FX pinning for
        audit reproducibility.} MVP uses a single FX table at solve; no
        per-run pinning.
  \item \label{item:multi-seg-pu-pwl} \textbf{Multi-tier per-ULD over-pivot
        rate.} Contingency only --- not found as a standard tariff form in
        the sources reviewed (IATA TACT, carrier BUC programmes,
        pivot-weight literature). If a carrier quotes one, $C_{a,g,u}$ would
        gain a PWL with SOS-2 or per-segment binaries.
  \item \label{item:lock-buyout} \textbf{Lock-buyout as a MILP decision.}
        MVP handles lock-break as an orchestrator decision between runs
        (\S\ref{sec:locked-commitments}); P1 could promote to a MILP binary
        $b_k$ with cost $\kappa_k^{\text{break}}$.
  \item \label{item:bin-packing} \textbf{Full 3D bin-packing inside ULDs.}
        MVP uses C.5b aggregate caps + pre-filter step 8. A future
        bin-packing model handles HAWBs with awkward piece dimensions
        precisely.
  \item \label{item:cfs-storage-demurrage} \textbf{CFS storage / demurrage
        cost in the objective.} MVP omits CFS dwell cost; optimizer may
        prefer a cheaper later flight without seeing offsetting storage
        cost. Sourced rate notes (2026-05-17 research): $\sim\$12$/CBM/day
        with $\$70$/day minimum at major US CFS (Imperial CFS,
        \url{https://www.imperialcfs.com/IPI/Warehouse}); demurrage
        post-free-period $\$5$--$15$/CBM/day; free periods 3--5 days export,
        5--7 days import (FreightAmigo summary,
        \url{https://www.freightamigo.com/en/blog/logistics/understanding-storage-fees-in-logistics-what-you-need-to-know/}).
        For a 5 CBM shipment dwelling 5 days past free at $\$12$/CBM/day =
        $\$300$ --- material on a 50-shipment/day book where $\sim 10\%$
        land in storage.
  \item \label{item:probabilistic-clearance} \textbf{Probabilistic clearance
        dwell with exam-hold probability.} MVP uses deterministic
        $\delta^{\text{cust}}_k$. P1: Bernoulli exam-hold event with
        destination-country-specific probability and conditional dwell
        distribution. Integrates with item~\ref{item:tt-quantile-binding}.
  \item \label{item:probabilistic-transit} \textbf{Probabilistic transit
        time objective.} Extend to $\Pr(\text{arrival} \leq T_k^{\text{abs}})
        \geq p$ as a chance constraint, or via Conditional Value-at-Risk on
        tardiness.
  \item \label{item:uld-availability} \textbf{ULD availability at CFS.} MVP
        assumes empty ULDs are always delivered to CFS before cargo cutoff;
        P1: model the ULD pool as a constraint.
  \item \label{item:cfs-dock-capacity} \textbf{CFS dock door capacity.} MVP
        assumes infinite CFS throughput; P1: model dock door and labor
        capacity.
  \item \label{item:multi-uld-types-per-shipment} \textbf{Multiple ULD
        types per HAWB.} MVP forces a HAWB to a single ULD type per
        \texttt{per\_uld\_pivot} arc; P1: allow splitting across LD3 + LD7.
  \item \label{item:alliance-slot-sharing} \textbf{Carrier alliance slot
        sharing.} MVP treats each contract independently; P1: model alliance
        Cargo pool slot exchanges where ULD positions on partner carriers can
        substitute.
  \item \label{item:spot-availability-uncertainty} \textbf{Spot capacity
        availability uncertainty.} MVP treats spot as always available; real
        spot has day-of-departure availability risk. P1: stochastic
        constraint or explicit booking-rejection events.
  \item \label{item:in-transit-hub-customs} \textbf{In-transit hub customs
        as a separate arc.} MVP folds in-transit clearance dwell into
        $\delta_a$ of the deconsolidation-dwell arc as data only; P1 models
        it as a distinct arc with its own cost and stochastic dwell.
  \item \label{item:per-hawb-cost-attribution} \textbf{Per-HAWB cost
        attribution under consolidation --- PROMOTED to MVP.} MILP returns
        per-MAWB cost; post-optimization attribution to per-HAWB is wired
        into the diagnostic output (\S\ref{sec:output-diagnostics}) via
        the proportional-to-CW rule. Refined allocation schemes (Shapley,
        proportional-to-actual-weight, proportional-to-volume) remain P1
        as alternative attribution modes; the MVP default is the simple
        proportional-to-CW rule. Retained in this list for changelog
        traceability only.
  \item \label{item:charter-and-bor} \textbf{Charter and broker-of-record
        procurement modes.} MVP catalog covers contracted, equalized,
        per-flight-pivot, and co-load only. P1: charter (full-aircraft tender)
        and broker-of-record handoff as additional supply-option types.
  \item \label{item:cwlinearizer-interface} \textbf{CWLinearizer interface
        for future rate-family extensions.} The current C.4c--d $\geq$
        inequality form (\S\ref{sec:cw-density-mixing}, \S\ref{sec:objective})
        is exact only under the \emph{monotonicity invariant}: every
        objective coefficient on $\text{CW}_{a,g}$ (and on $\text{Wt}$,
        $\text{Wv}$, $\text{pivot}$, $\text{over}$, $\tau$, $\text{pen}$)
        is non-negative. A future rate family with a \emph{negative}
        coefficient --- e.g.\ a $\$2$/kg rebate above 5000 kg, a
        peak-season volume kicker, a marketing credit per shipment ---
        would break the invariant: the minimization would drive
        $\text{CW}_{a,g}$ to whichever inequality is binding in the
        negative term's direction, producing a wrong cost. The post-solve
        assertion in \S\ref{sec:output-diagnostics} catches the
        violation but only after the wrong cost is acted on.

  \textbf{Design sketch.} Introduce a \texttt{CWLinearizer} abstraction
  with two modes:
    \begin{itemize}[noitemsep]
      \item \emph{Inequality mode} (current): C.4c--d as $\geq$. Exact
            under monotonicity invariant. Cheapest LP.
      \item \emph{Equality mode}: $\text{CW} = \text{Wt} + s_{\text{Wt}}$,
            $\text{CW} = \text{Wv} + s_{\text{Wv}}$ with complementarity
            $s_{\text{Wt}} \cdot s_{\text{Wv}} = 0$ linearized by an
            auxiliary binary $\delta_{a,g}$ picking the binding side. Two
            additional inequalities and one binary per MAWB; LP-tight.
    \end{itemize}

  \textbf{Catalog-time validation.} At each rate-family registration
  (a tenant adds an offer to the supply catalog), the system computes
  $\partial \text{cost}^{\text{MAWB}}_{a,g} / \partial \text{CW}_{a,g}$
  (and the analogous partials on $\text{Wt}$, $\text{Wv}$,
  $\text{pivot}$, $\text{over}$, $\tau$, $\text{pen}$). If any partial is
  negative, the registration is rejected (or auto-promotes the rate
  family's MAWB-arc binary $z_{a,g}$ to use Equality-mode CW
  linearization) before the offer ever enters a solve. This shifts the
  invariant from a post-solve assertion to a pre-solve guarantee.
  Implementation lives in the supply-catalog data-layer
  (\texttt{data\_model.md}); the LaTeX continues to default to
  Inequality mode under the catalog-time guarantee.
  \item \label{item:awb-stock-management} \textbf{AWB stock management.}
        MVP assumes AWB numbers are always available; real operations
        constrain by IATA stock allocation per forwarder per origin. P1:
        per-tenant AWB-stock accumulator with re-stock triggers.
  \item \label{item:time-windowed-carrier-rules} \textbf{Time-windowed /
        seasonal carrier-policy rules.} Hajj / monsoon / Lunar New Year
        embargoes and capacity rules. Deferred to the separate rules-engine
        model; consumed via the pre-filter carrier-policy predicate when
        ready.
\end{enumerate}

%%--------------------------------------------------------------------
\section{Tractability and Scaling Roadmap}
\label{sec:scaling-roadmap}
%%--------------------------------------------------------------------

This section records the tractability levers, decomposition options, and
production-deployment strategies for the air MILP. None are active model
changes; all are P1-or-later, gated on empirical solve-time data from
production. The intent is to keep the analysis with the model so future-you
doesn't re-derive it. Citations are inline where the lever maps to an
established OR result.

\subsection*{Base-scale estimate (MVP)}

At $|K| = 100$ HAWBs, $|A| = 500$, $|A_k|_{\text{avg}} \approx 15$ (10\%
pre-filter survival assumed; to be measured),
$|A^{\text{MAWB}}| \approx 250$, $|A^{\text{MFB}}| \approx 80$,
$|A^{\text{pu}}| \approx 40$, $|U_a|_{\text{avg}} \approx 2.5$,
$|G_a|_{\text{avg}} \approx 1.5$, $|M| \approx 200$:

\begin{itemize}[noitemsep]
  \item \textbf{Binaries.} $x_{k,a}: \approx 1{,}500$.
        $z_{a,g}: \approx 300$.
        $\gamma_{a,g,b}: \approx 720$.
        \textbf{Total $\sim 2{,}500$ binaries.}
  \item \textbf{Integer (general).} $\eta_{a,g,u}: \approx 150$, bounded
        $[0, N_{a,u}]$ (typically $\leq 10$).
  \item \textbf{Continuous.} $t_k^{D_k^{\text{node}}}: \approx 100$ (one per
        HAWB). $\text{CW} + \text{Wt} + \text{Wv}: \approx 900$. $\text{BW}:
        \approx 720$. $\text{pivot} + \text{over} + \tau + \text{pen}: \approx
        350$ (adds $|K| \approx 100$ for $\text{pen}_k$).
        \textbf{Total $\sim 2{,}100$ continuous} (down from $\sim 4{,}500$
        prior to the J19 reshape: $\sim 2{,}400$ intermediate-node $t_k^n$
        variables removed by forward-time-window propagation,
        \S\ref{sec:fwd-time-propagation}).
  \item \textbf{Constraints.} C.1: $\approx 2{,}500$. C.2a/b: $\approx 3{,}000$.
        C.4: $\approx 1{,}200$. C.5 $+$ C.5b $+$ C.5c: $\approx 500$.
        C.10a (destination-arrival definition, one row per HAWB):
        $\approx 100$. C.10b ($\tau$ definition + $J = 5$ PWL cuts on
        $\text{pen}_k$): $\approx 600$. C.13: $\approx 80$. \emph{Fallback
        arc adds $|K|$ rows to C.1 (already counted) and one binary
        $x_{k, a^{\text{fb}}_k}$ per HAWB.}
        \textbf{Total $\sim 8{,}000$ constraints} (down from $\sim 10{,}500$
        prior to the J19 reshape: $\sim 1{,}500$ C.6 rows and $\sim 300$
        C.9 rows eliminated).
  \item \textbf{Expected HiGHS solve time:} seconds to a few minutes
        single-component, well-tuned.
  \item \textbf{Phase 2 ($|K| = 500$):} $\approx 12{,}500$ binaries, $\approx
        50{,}000$ constraints --- minutes to $\sim 1$ hour monolithic;
        decomposition mandatory.
  \item \textbf{Phase 2 ($|K| = 1500$):} monolithic infeasible without
        column generation (see \texttt{scalability.md}).
\end{itemize}

\subsection*{Walking-skeleton instrumentation (mandatory from v4 of the build)}
\label{sec:walking-skeleton-instrumentation}

Measure these on every solve; output to structured log files alongside
solutions. They answer the scale questions empirically before they bite.
See \texttt{TEST\_PLAN.md \S10}.

\begin{table}[H]
\centering
\small
\caption{Walking-skeleton instrumentation --- metrics output on every solve.}
\label{tab:walking-skeleton-metrics}
\begin{tabular}{p{4.7cm}p{4.0cm}p{6.2cm}}
\toprule
Metric & Output & Purpose \\
\midrule
$|A_k|$ histogram + per-predicate drop counts (the 8 pre-filter predicates) & \texttt{pre\_filter\_stats.jsonl} per HAWB & Determines $x_{k,a}$ binary count. The single most important number. \\
$|G_a|$ distribution + activated bucket count $\sum z_{a,g}$ & \texttt{bucket\_stats.jsonl} per solve & Drives $|M|$ and $\gamma$ count estimates. \\
LP-vs-MIP gap by constraint family & \texttt{lp\_gap\_breakdown.json} & Identifies which family loosens the LP most. \\
Forward-time-window pruning rate, \emph{subdivided by cause} --- cutoff vs reachability vs ETA-anchor vs dispatch-feasibility (\S\ref{sec:fwd-time-propagation}) & \texttt{fwd\_time\_window\_drops.jsonl} & Confirms time feasibility is binding pre-MILP and exposes which time bound (cutoff stack, $T_k^{\text{abs}}$, ETA, dispatch lead time) is the dominant pruning driver. \\
Internal-MCT data-bug compliance --- assert $\mu_a \geq \sum_{\text{legs}} (\text{block time} + \text{MCT})$ for every multi-leg arc emitted & \texttt{internal\_mct\_compliance.jsonl} & One-line invariant. Catches data-bug class where the graph generator emits a multi-leg arc whose $\mu_a$ doesn't include the internal MCT --- the MILP would silently route through. \\
LP-gap contribution by constraint family, \emph{distinguishing C.13b-1 from C.13b-2} & part of \texttt{lp\_gap\_breakdown.json} & The pivot LP-gap source matters: C.13b-2 (per-flight pivot floor) is the place where fractional $\eta$ in LP relaxation lets the optimizer artificially lower the pivot. If C.13b-2 dominates total LP gap on BSA-heavy instances, the per-ULD-slot disaggregation (deferred) moves up the priority list. \\
LP-gap contribution from break-disaggregation family & part of \texttt{lp\_gap\_breakdown.json} & If on TACT-heavy instances the break-disaggregation dominates total LP gap, evaluate either (a) commercial solver with native indicators, or (b) per-arc convex-hull formulation (Sherali-Adams lift) for small-break-count cases. \\
Connected components of $H$ (HAWB-sharing-arc) --- shadow & \texttt{component\_decomposition\_shadow.jsonl} & Indicator for when decomposition becomes mandatory. \\
Super-shipment equivalence classes --- shadow & \texttt{aggregation\_potential} & Decides default for \texttt{consolidation\_mode}. \\
BSA-contract cross-coupling fraction & \texttt{bsa\_coupling\_fraction} & Predicts Lagrangian-relaxation value for Phase 2 decomposition. \\
Post-solve invariant assertions & \texttt{invariant\_violations.json} & $\text{CW} = \max(\text{Wt}, \text{Wv})$; $\tau_k = \max(0, \text{lateness})$; $\sum_g \eta \leq N$; no $z = 1$ with empty bucket. Catches silent bugs. \\
HiGHS phase breakdown & callback-derived \texttt{solve\_phase\_times.json} & Tells us where time actually goes. \\
Fractional-$x$ frequency at LP root --- fraction of $x_{k,a}$ fractional at root LP, and fraction in the ``hard'' $[0.4, 0.6]$ range & \texttt{lp\_root\_fractional\_x.jsonl} & Formulation-quality signal. High fractional-$x$ rate means the arc-based LP relaxation is intrinsically weak for this instance; direct signal for ``pivot to column generation,'' better than wall-clock alone. \\
Per-arc activated $z_{a,g}$ distribution --- histogram of $|\{g : z_{a,g} = 1\}|$ across arcs $a$ & \texttt{activated\_z\_distribution.jsonl} & Formulation-simplification signal. If the distribution is mostly 0 or 1 (single-group MAWBs everywhere), the $(a, g)$ formulation is overkill and could be collapsed to a simpler partition; if many arcs have $|\{g\}| \geq 2$, consolidation is doing real work. \\
Fractional fallback usage --- $\sum_k x_{k, a^{\text{fb}}_k}^{\text{LP root}} - |K^{\text{fb}}|^{\text{MIP}}$ per solve & \texttt{fractional\_fallback\_usage} & $C^{\text{fallback}}$ sizing check (\S\ref{sec:hawb-params}). If $> 1$ persistently, the LP relaxation is exploiting fallback smearing and $C^{\text{fallback}}$ is set too low for some HAWBs --- recompute the sizing formula at the per-tenant level. \\
\bottomrule
\end{tabular}
\end{table}

\subsection*{Simplification levers (P1, ordered by impact)}

\begin{enumerate}
  \item \label{item:scale-prefilter} \textbf{Pre-filter aggressiveness has
        100$\times$ leverage; instrument first.} The pre-filter applies eight
        predicates (\S\ref{sec:prefilter}). Each typically drops 30--60\% of
        arcs; compounded survival $\sim 1$--5\% post-filter. Pre-filter
        quality has more leverage than any single solver-tuning change.
        Required instrumentation: log $|A_k|$ post-filter per HAWB. If
        survival rate $> 5\%$, fix the filter before any other scale work.
  \item \label{item:scale-hawb-aggregation} \textbf{HAWB aggregation
        (\texttt{consolidation\_mode = preprocess}).} At Phase 2 load,
        HAWBs sharing $(g, O_k, D_k^{\text{node}}, \text{deadline tier},
        \text{sp}(k))$ are interchangeable from the MILP's standpoint.
        Treat each equivalence class as a super-HAWB with combined $w_k$,
        $v_k$. Largest single Phase-2 tractability win.
        \textbf{Promoted to walking-skeleton-instrumented} via the
        \texttt{aggregation\_potential} shadow metric --- decide
        default-on / default-off from empirical data, not a priori.
  \item \label{item:scale-bucket-dominance} \textbf{Per-MAWB-candidate
        dominance pre-filter.} Pre-MILP, drop arc $a' \in A^{\text{MAWB}}$
        if $\exists a \in A^{\text{MAWB}}$ with same $G_a$, same O-D,
        $\mu_a \leq \mu_{a'}$, $\text{cap}_a \geq \text{cap}_{a'}$, and
        $\text{rate}_a(w) \leq \text{rate}_{a'}(w)$ for all feasible $w$.
        Hard pre-filter step. Estimate 20--40\% offer pruning at MVP.
        Source: standard variable fixing via reduced-cost dominance.
  \item \label{item:scale-component-decomposition} \textbf{Component-wise
        solve on the supply graph $H$.} After subgraph construction, build
        $H$ where HAWBs are connected if they share any MAWB-arc (and group)
        or BSA contract; solve each connected component as an independent
        MILP. Expected $\sim 3$--10 components per tenant batch at MVP scale.
        Combined with HAWB aggregation, $|K| = 500$ split into 5 components
        of $\approx 100$, each $\sim 30$s under simplifications: $\sim 30$s
        total parallel. Source: standard MCNF decomposition
        (Ahuja--Magnanti--Orlin \emph{Network Flows} \S17). BSA-contract
        coupling (\(C^{\text{eq}}\) chargeable accumulator across components)
        is the typical breaker --- measure \texttt{bsa\_coupling\_fraction}
        first; if high, plan for Lagrangian relaxation of just C.13a.
  \item \label{item:scale-warm-start} \textbf{Warm-start re-solves from
        prior solution.} Rolling-horizon re-solves overlap heavily with
        previous solves (only locked prefixes and new HAWBs change). HiGHS
        supports MIP warm-start via \texttt{Highs::setSolution}. Expected
        3--10$\times$ speedup on re-solves vs cold MIP. Source: HiGHS docs;
        empirical results in Erera et al.~2017 on rolling-horizon service
        network design.
  \item \label{item:scale-mip-gap} \textbf{Production MIP gap $\geq 1\%$,
        never solve to 0\%.} Freight routing literature (Powell ADP corpus,
        Topaloglu, Erera) consistently treats 1--5\% gap as
        production-acceptable: cost variance from sell-rate negotiation, FX,
        fuel surcharge volatility, and override behavior dwarfs the
        optimization gap. Recommended: 1\% for premium / SLA-critical
        batches; 3\% for standard / multimodal-economy. Expected speedup vs
        exact: 3--5$\times$ on $|K| = 300$.
\end{enumerate}

\subsection*{Strategy notes (production deployment)}

\begin{enumerate}[resume]
  \item \label{item:strat-column-generation} \textbf{Column generation ---
        deferred until empirical $> 15$ min at production scale.}
        Dantzig--Wolfe / column generation (pricing subproblem $=$
        shortest-path-with-resource-constraints per HAWB on $A_k$ with
        capacity duals from master) is the classical heavy weapon for large
        MCNF. Premature at MVP; trigger only when monolithic exceeds 15 min
        on production-typical instances. See \texttt{scalability.md} for
        the SPPRC pricing-subproblem sketch. Source: Desrosiers \& Lübbecke
        (2005) primer.
  \item \label{item:strat-commercial-solver} \textbf{Commercial solver
        (Gurobi / CPLEX / COPT) transition.} HiGHS handles MVP scale at 1\%
        gap. Commercial worth evaluating when (a) per-tenant $|K| > 500$
        sustained, (b) re-solve latency budget $< 30$s for urgent rescue
        shipments, or (c) LP relaxation gap on time-binding instances
        consistently $> 10\%$ (Gurobi cuts + branching typically
        2--5$\times$ faster than HiGHS on big-M MCNF per Mittelmann
        benchmarks). Business case: Gurobi licensing $\sim \$10$K/year/seat
        justifiable above $\sim \$1$M ARR / 5+ active tenants. Until then,
        HiGHS plus the simplifications above.
  \item \label{item:strat-sla-pickup-anchor} \textbf{SLA pickup-time anchor
        --- operational tradeoff, not a bug.} C.10b anchors against the
        \emph{parameter} $t_k^{\text{rdy,early}}$ (via $\Delta_k$ which
        derives from $T_k^{\text{SLA}} = t_k^{\text{rdy,early}} +
        T^{\text{SLA}}_{\text{sp}(k)}$). Operationally defensible: shipper
        SLAs are measured from when cargo is offered for pickup, not from
        when the forwarder elects to dispatch. Side effect: the optimizer
        is not penalized for paths whose effective pickup is later in the
        ready window if that buys cheaper downstream slots. Acceptable
        unless tenant policy requires earliest-pickup discipline; in that
        case, anchor $\Delta_k$ relative to a per-arc-effective-pickup
        attribute carried into $\text{arr\_dest}(k, a)$ at graph build.
\end{enumerate}

\paragraph{Punch list summary.}
\begin{itemize}[noitemsep]
  \item \textbf{Ship as-is} if instance size $\leq 100$ HAWBs / tenant /
        solve.
  \item For 100--500 HAWBs per tenant: implement
        items~\ref{item:scale-prefilter},~\ref{item:scale-hawb-aggregation},~\ref{item:scale-bucket-dominance},~\ref{item:scale-component-decomposition},~\ref{item:scale-warm-start},~\ref{item:scale-mip-gap}.
  \item Beyond $\sim 750$ HAWBs per tenant: column generation
        (item~\ref{item:strat-column-generation}) or commercial solver
        (item~\ref{item:strat-commercial-solver}).
  \item Pre-filter instrumentation (item~\ref{item:scale-prefilter}) is the
        first thing to wire up regardless of scale.
\end{itemize}

%%--------------------------------------------------------------------
\appendix
%%--------------------------------------------------------------------

\section{Air Cargo Rate Types --- Worked Examples}
\label{sec:appendix-rate-examples}

Worked examples for the three rate acronyms in the \texttt{min\_flat\_breaks}
rate family (\S\ref{sec:supply-option-catalog}). All rate-card numbers below
are \emph{illustrative}, not published tariffs; the \emph{structure} --- weight
breaks and the next-break-down rule --- is real.

\paragraph{The next-break-down rule.}
An air cargo rate card lists weight \emph{breaks} --- e.g.\ 45, 100, 300, 500,
1000 kg --- each with a lower per-kg rate than the band below. The shipper
may always declare a \emph{higher} break weight when that is cheaper: the
billed cost is the minimum, over all breaks $b$, of $\text{rate}_b \cdot
\max(\text{CW},\, \text{break}_b)$.

\subsection{TACT --- The Air Cargo Tariff (General Cargo Rate)}
\label{sec:appendix-tact}

\textbf{TACT} (``The Air Cargo Tariff'') is IATA's published tariff. A TACT
\emph{General Cargo Rate} (GCR) is the default rate on a lane for ordinary
cargo, with no commodity discount.

Illustrative TPE$\to$JFK GCR card:
\begin{table}[H]
\centering
\caption{Illustrative TPE$\to$JFK TACT General Cargo Rate card.}
\label{tab:tact-card}
\begin{tabular}{lcccccc}
\toprule
Break (kg)   & 0 & 45 & 100 & 300 & 500 & 1000 \\
Rate (\$/kg) & 6.00 & 5.50 & 4.80 & 4.20 & 3.80 & 3.20 \\
\bottomrule
\end{tabular}
\end{table}

A 280 kg MAWB sits in the 100--300 kg band (rate \$4.80): billed at its own
weight that is $280 \times 4.80 = \$1{,}344$. But declaring the next break
down --- 300 kg at \$4.20 --- costs $300 \times 4.20 = \$1{,}260$, which is
cheaper. \textbf{Billed: \$1{,}260}, declaring 300 kg. That is the
next-break-down rule.

\subsection{SCR --- Specific Commodity Rate}
\label{sec:appendix-scr}

An \textbf{SCR} is a discounted rate for a particular commodity (identified
by a 4-digit IATA commodity code) on a particular lane. Lower than the GCR,
same weight-break structure, requires commodity-code declaration.

Illustrative SCR card, same lane, garments commodity code:
\begin{table}[H]
\centering
\caption{Illustrative TPE$\to$JFK SCR (garments commodity code) rate card.}
\label{tab:scr-card}
\begin{tabular}{lcccc}
\toprule
Break (kg)   & 0 & 100 & 500 & 1000 \\
Rate (\$/kg) & 4.50 & 3.90 & 3.30 & 2.90 \\
\bottomrule
\end{tabular}
\end{table}

A 460 kg MAWB is in the 100--500 kg band (\$3.90): $460 \times 3.90 =
\$1{,}794$. The next break down --- 500 kg at \$3.30 --- is $500 \times 3.30 =
\$1{,}650$, cheaper. \textbf{Billed: \$1{,}650.}

\subsection{CCR --- Class Rate}
\label{sec:appendix-ccr}

A \textbf{CCR} prices a class of cargo as a percentage of the GCR. Surcharge
classes (valuables, live animals, human remains) bill above 100\%; a few
classes (e.g.\ newspapers) bill below. The class percentage scales every GCR
break rate uniformly; the next-break-down rule applies to the scaled card.

Illustrative live-animal class at 150\% of GCR:
\begin{table}[H]
\centering
\caption{Illustrative live-animal CCR at 150\% of GCR --- rate card.}
\label{tab:ccr-card}
\begin{tabular}{lcccccc}
\toprule
Break (kg) & 0 & 45 & 100 & 300 & 500 & 1000 \\
\midrule
GCR rate (\$/kg) & 6.00 & 5.50 & 4.80 & 4.20 & 3.80 & 3.20 \\
CCR rate $= 1.50 \times$ GCR (\$/kg) & 9.00 & 8.25 & 7.20 & 6.30 & 5.70 & 4.80 \\
\bottomrule
\end{tabular}
\end{table}

For a 280 kg live-animal MAWB:
\begin{table}[H]
\centering
\caption{Live-animal CCR --- billed-weight calculation per break for a 280~kg MAWB.}
\label{tab:ccr-calc}
\begin{tabular}{lcccccc}
\toprule
Break (kg)         & 0 & 45 & 100 & 300 & 500 & 1000 \\
\midrule
Billed weight (kg) & 280 & 280 & 280 & 300 & 500 & 1000 \\
CCR rate (\$/kg)   & 9.00 & 8.25 & 7.20 & 6.30 & 5.70 & 4.80 \\
Cost (\$)          & 2{,}520 & 2{,}310 & 2{,}016 & \textbf{1{,}890} & 2{,}850 & 4{,}800 \\
\bottomrule
\end{tabular}
\end{table}

The minimum is \textbf{\$1{,}890} at the 300 kg break --- billed as 300 kg of
live-animal cargo. Equals $1.50 \times \$1{,}260$, the GCR next-break-down
cost for the same 280 kg (\S\ref{sec:appendix-tact}): since the multiplier
scales every break uniformly, billing the CCR card directly and scaling the
GCR result give the same answer.

%%====================================================================
\section{Forward-time-window propagation $\equiv$ big-M MILP (proof)}
\label{app:fwd-prop-equiv}
%%====================================================================

This formalizes the equivalence asserted in \S\ref{sec:fwd-time-propagation}:
for deterministic transits, the graph-build per-outbound admission and the
big-M MILP time-propagation family define the \emph{same} feasible route set.

\paragraph{Setup and assumptions.}
Fix a HAWB and drop its subscript. \textbf{(A1) Determinism:} each arc has a
fixed $\tau_a = \text{transit}(k,a) \geq 0$, with $\tau_a > 0$ on physical
transport arcs. \textbf{(A2) Scheduled air arcs:} an air arc $a=(n_1,n_2)$
leaving a POL has a fixed cutoff $\text{CO}_a^*$, scheduled departure
$\text{STD}_a$, block time $\mu_a$, and scheduled arrival
$\text{ETA}_a = \text{STD}_a + \mu_a$; the cutoff is the only admission gate it
imposes, and conditional on admission its head time is pinned to
$\text{ETA}_a$. \textbf{(A2$'$) Scheduling consistency:} $\text{CO}_a^* \leq
\text{STD}_a$ (cutoff precedes departure), so whenever the cutoff is cleared,
$\text{ETA}_a = \text{STD}_a + \mu_a \geq \text{CO}_a^* \geq t^{\text{lo}}_{n_1}$
--- the air head is at-or-after the earliest tail arrival. \textbf{(A3) Origin
release:} $t_O \geq t^{\text{rdy,early}}$. \textbf{(A4) Acyclicity:} the
candidate subgraph is acyclic. This is delivered by A1+A2$'$: arrival time
strictly increases along every arc ($\tau_a > 0$ on transport arcs; an air arc
moves $t^{\text{lo}}_{n_1} \leq \text{CO}_a^* \leq \text{ETA}_a$ forward), so no
node repeats. A4 is used only in the Corollary (path-faithfulness), not in the
Lemma. A \emph{route} is a simple $O \to D^{\text{node}}$ path $P$ with
incidence vector $x^P$ ($x_a = 1 \iff a \in P$).

The graph-build \emph{earliest-arrival} along $P$ is the forward sum, with air
arcs resetting to ETA:
\[
t^{\text{lo}}_n =
\begin{cases}
t^{\text{rdy,early}}, & n = O,\\
t^{\text{lo}}_{n_1} + \tau_a, & a=(n_1,n)\ \text{ground/dwell},\\
\text{ETA}_a, & a=(n_1,n)\ \text{air}.
\end{cases}
\]
Per-outbound admission keeps a POL air arc $a$ iff $t^{\text{lo}}_{n_1} \leq
\text{CO}_a^*$. The MILP family carries node-time variables
$t_n \in [t^{\text{rdy,early}}, T^{\text{abs}}]$ and \emph{$x$-gated} rows:
ground $t_{n_2} \geq t_{n_1} + \tau_a - M(1-x_a)$; gated air-head equality
$\text{ETA}_a - M(1-x_a) \leq t_{n_2} \leq \text{ETA}_a + M(1-x_a)$; air cutoff
$t_{n_1} \leq \text{CO}_a^* + M(1-x_a)$. Gating every row (not just the cutoff)
is necessary: an ungated air-head equality would impose a single $\text{ETA}$
on a node that heads two candidate flights, and ungated ground rows would
couple unselected arcs' tails.

\paragraph{Lemma (earliest arrival is the lower envelope).}
For $x = x^P$, every $\{t_n\}$ satisfying the on-path ($x_a = 1$) ground and
air-head rows obeys $t_n \geq t^{\text{lo}}_n$ for all $n \in P$, and
$t = t^{\text{lo}}$ satisfies those rows.

\emph{Proof.} A simple path has distinct, linearly ordered nodes, so induction
along $P$ is well-founded (no global acyclicity needed). Base: $t_O \geq
t^{\text{rdy,early}} = t^{\text{lo}}_O$. Ground step: $t_{n_2} \geq t_{n_1} +
\tau_a \geq t^{\text{lo}}_{n_1} + \tau_a = t^{\text{lo}}_{n_2}$ ($\tau_a \geq
0$). Air step: the active equality gives $t_{n_2} = \text{ETA}_a =
t^{\text{lo}}_{n_2}$; by A2$'$ this value is $\geq t^{\text{lo}}_{n_1}$, so the
envelope stays non-decreasing. Setting $t = t^{\text{lo}}$ meets every on-path
row with equality. \hfill$\square$

\paragraph{Proposition (per-route equivalence).}
Take $M \geq T^{\text{abs}} + \max_a \tau_a$ (so every row of an arc with
$x_a = 0$ is slack given $t_n \in [t^{\text{rdy,early}}, T^{\text{abs}}]$). For
a route $P$,
\[
\big(\exists\,\{t_n\}\ \text{MILP-feasible with } x = x^P\big)
\iff
\big(t^{\text{lo}}_{n_1(a)} \leq \text{CO}_a^*\ \ \forall\ \text{air-at-POL } a \in P\big).
\]

\emph{Proof.} \textbf{Deactivation:} for $a \notin P$, $x_a = 0$ makes all
three of its rows slack at the chosen $M$ (each RHS slack exceeds the
$[t^{\text{rdy,early}}, T^{\text{abs}}]$ span). So only on-path rows bind.
\textbf{($\Leftarrow$)} Assume the RHS; set $t := t^{\text{lo}}$ on nodes of
$P$ and $t_n := t^{\text{rdy,early}}$ off $P$. By the Lemma the on-path
ground/air-head rows hold; each on-path cutoff holds since $t^{\text{lo}}_{n_1}
\leq \text{CO}_a^*$; off-path rows are slack. Feasible.
\textbf{($\Rightarrow$)} Given a feasible witness, each on-path cutoff gives
$t_{n_1} \leq \text{CO}_a^*$ and the Lemma gives $t_{n_1} \geq
t^{\text{lo}}_{n_1}$; chaining, $t^{\text{lo}}_{n_1} \leq \text{CO}_a^*$.
\hfill$\square$

\paragraph{Corollary (pre-projection).}
Restrict attention to the time rows above. Under C.1 flow conservation with
unit $O \to D$ balance, integrality $x \in \{0,1\}$, and A4 acyclicity, the
feasible $x$ are exactly the simple $O \to D$ path incidence vectors (a unit
integral $O$-to-$D$ flow on an acyclic graph decomposes into a single simple
path; no disjoint circulation survives acyclicity). The other constraint
families (C.2 MAWB linkage, C.5 capacity, \ldots) intersect this set further
but never \emph{admit} a route the time rows reject, so the time-row
projection is the relevant one. Admission is conjunctive over a route's arcs,
and both formulations apply the identical per-arc test $t^{\text{lo}}_{n_1}
\leq \text{CO}_a^*$ against \emph{each} air arc's own cutoff (here
$t^{\text{lo}}_{n_1}$ is the per-route value, i.e.\ the label of $n_1$'s
incoming arc on $P$). Hence
\[
\operatorname{proj}_x\big(\text{time-row feasible region}\big) =
\{\,x^P : P\ \text{graph-admitted}\,\}.
\]
With destination reachability folded in (Remark 4), $A_k = \bigcup_{P\ \text{end-to-end feasible}} P$. \hfill$\square$

\paragraph{Remarks.}
\textbf{(1) The ETA-fix is necessary; the STD-floor is equivalent only under
A2$'$.} Replacing the gated air-head equality with a generic block-time row
$t_{n_2} \geq t_{n_1} + \mu_a$ lets early cargo drive $t_{n_2}$ below the
schedule, over-admitting arcs. A scheduled-departure floor $t_{n_1} \geq
\text{STD}_a$ with block propagation yields $t_{n_2} \geq \text{ETA}_a$ (not
$=$); this matches the ETA-fix \emph{as a projection onto $x$} only under A2$'$
(where the cutoff already implies the floor), not as an identity of the
continuous $t$-regions. \textbf{(2) Earliest arrival, not the full window.}
Cutoff admission is a pure upper bound, so only the lower envelope
$t^{\text{lo}}$ can bind it; $t^{\text{hi}}$ is irrelevant to this test. Hard
delivery windows $t_n \leq \overline{w}_n$ compose by re-running the Lemma.
\textbf{(3) Determinism is load-bearing.} (A1) makes $t^{\text{lo}}_n$ a
deterministic scalar equal to the binding lower bound. Under stochastic
transits the test becomes $\Pr(\text{arrival} \leq \text{CO}_a^*) \geq p$ and
$t^{\text{lo}}$ a sum of per-arc quantiles; since the quantile of a sum $\neq$
the sum of quantiles, the per-arc projection breaks. The equivalence holds
strictly when a \emph{single} deterministic statistic drives both cutoff
admission and destination arrival (mixing mean-transit cutoffs with a
quantile arrival re-opens the gap) --- the boundary flagged in
\S\ref{sec:deferred}, item~\ref{item:tt-quantile-binding}.
\textbf{(4) Dispatch is a per-arc row; reachability is not.} The dispatch
check $t^{\text{rdy,early}} \leq \text{CO}_a^* - \lambda^{\text{disp}}_k$ is
another per-arc upper-bound gate and composes verbatim. Destination
reachability $t^{\text{lo}}_{D} \leq T^{\text{abs}}$ is \emph{not} a per-arc
row: it is a path-existence predicate (a min over downstream paths from
$\text{head}(a)$ to $D$). The graph build enforces it as a separate backward
pass, and $A_k$ equals the forward-admitted set \emph{intersected with} the
backward-reachable set, which is exactly $\bigcup_{P\ \text{end-to-end feasible}} P$.

\end{document}
